% Le Cameroun dans la coopération multilatérale — ECM Tle Tech — Cours signé OnBuch+
% Programme officiel MINESEC. Compiler avec : tectonic N03-cameroun-cooperation-multilaterale.tex
\newcommand{\DOCMATIERE}{ECM}
\newcommand{\DOCNIVEAU}{Tle Tech}
\newcommand{\DOCMODULE}{ECM – classe de Terminale technique}
\newcommand{\DOCLECON}{N03}
\newcommand{\DOCTITRE}{Le Cameroun dans la coopération multilatérale}
% OnBuch+ — préambule « cours signés » ENRICHI (compilé avec tectonic / XeLaTeX)
% Système visuel « Pop » d'OnBuch+ : fond crème (jamais blanc), bordures encre
% épaisses, ombres « dures » décalées, titres Archivo Black en capitales.
% Version MATHÉMATIQUES : graphes (pgfplots/tikz), boîtes pédagogiques
% (définition, propriété, méthode, attention, le savais-tu, exercice résolu,
% exercice, TP, bilan), page de garde et signature OnBuch+ ancrée en bas.
%
% Macros attendues dans main.tex AVANT \input{preamble} :
%   \DOCMATIERE  \DOCNIVEAU  \DOCMODULE  \DOCLECON (numéro)  \DOCTITRE
\documentclass[11pt]{article}

\usepackage{fontspec}
\usepackage[french]{babel}
\usepackage[a4paper,margin=2cm,top=3.4cm,bottom=2.8cm,headheight=1.4cm,footskip=1.3cm]{geometry}
\usepackage{amsmath,amssymb,amsfonts}
\usepackage{mathtools} % \xmapsto, \coloneqq... souvent utilisés par les modèles
\usepackage{cancel} % simplifications barrées (bilans de forces)
% \abs{z} (module, valeur absolue) et \norm{u} (norme), utilisés par les modèles
\providecommand{\abs}{}\let\abs\relax\DeclarePairedDelimiter{\abs}{\lvert}{\rvert}
\providecommand{\norm}{}\let\norm\relax\DeclarePairedDelimiter{\norm}{\lVert}{\rVert}
\usepackage[table]{xcolor} % [table] : fournit aussi \rowcolors (alternance de lignes)
\usepackage{pagecolor}
\usepackage{tikz}
\usetikzlibrary{arrows.meta,positioning,calc,shapes.geometric,shapes.misc,shapes.arrows,decorations.pathmorphing,decorations.pathreplacing,decorations.markings,patterns,shadows,mindmap,trees,fit,backgrounds,matrix,intersections,angles,quotes}
\usepackage{pgfplots}
\usepackage[version=4]{mhchem} % équations nucléaires \ce{^{235}_{92}U}
\usepackage[european,siunitx]{circuitikz} % schémas électriques
\pgfplotsset{compat=1.18}
\usepgfplotslibrary{fillbetween,groupplots} % aires entre courbes (calcul intégral)
\usepackage{enumitem}
\usepackage{fancyhdr}
% [most] charge toutes les bibliothèques tcolorbox (listings, minted, raster,
% documentation...) et gonfle inutilement la mémoire TeX sur un document long
% (~300 boîtes) : on ne charge que ce qui est réellement utilisé.
\usepackage[skins,breakable]{tcolorbox}
\usepackage{graphicx}
\usepackage{colortbl}
\usepackage{booktabs}
\usepackage{tabularx}
\usepackage{multirow}
\usepackage{diagbox} % case d'en-tête diagonale des tableaux à double entrée
% Colonnes extensibles centrée / gauche / droite (C, L, R), souvent utilisées par les modèles
\newcolumntype{C}{>{\centering\arraybackslash}X}
\newcolumntype{L}{>{\raggedright\arraybackslash}X}
\newcolumntype{R}{>{\raggedleft\arraybackslash}X}
\usepackage{array}
\usepackage{multicol}
\usepackage{siunitx}
% parse-numbers=false : accepte aussi \SI{2,5 \times 10^{-3}}{...}
\sisetup{locale=FR,output-decimal-marker={,},group-minimum-digits=4,parse-numbers=false}
\DeclareSIUnit\ohm{\ensuremath{\symup{\Omega}}} % U+2126 absent de la police
\DeclareSIUnit\division{div} % réglages d'oscilloscope (ms/div, V/div)
\DeclareSIUnit\tr{tr} % tours (vitesse de rotation, ex. \SI{1500}{\tr\per\minute})
\DeclareSIUnit\molar{M} % concentration molaire (ex. \SI{3}{\molar}, \SI{1}{\micro\molar})
\DeclareSIUnit\an{an} % année (le modèle confond parfois avec \year, un compteur TeX) — ex. \SI{5}{\kilo\meter\per\an}
\DeclareSIUnit\FCFA{FCFA} % franc CFA (ex. \SI{500}{\FCFA\per\kilo\gram})
\DeclareSIUnit\degreeBrix{\textdegree Brix} % teneur en sucre d'un sirop/jus (réfractométrie)
\DeclareSIUnit\month{mois} % le modèle utilise parfois \month, un compteur TeX, comme unité de durée
\DeclareSIUnit\microliter{\micro\liter} % le modèle écrit parfois \microliter en un seul token
\DeclareSIUnit\million{millions} % dénombrement (ex. \SI{15}{\million\per\mL} spermatozoïdes)
\usepackage{qrcode}
% \degree : commande fréquemment utilisée par les modèles pour un angle ou une
% température en dehors de \SI{}{} (non fournie par les paquets chargés).
\providecommand{\degree}{^{\circ}}
% \qed (fin de démonstration), utilisé par les modèles sans amsthm
\providecommand{\qed}{\hfill$\square$}
% \barre : double barre des valeurs interdites dans un tableau de variations (tkz-tab)
\providecommand{\barre}{\ensuremath{\|}}
% environnement proof (amsthm non chargé) utilisé par les modèles
\ifdefined\proof\else\newenvironment{proof}[1][Démonstration]{\par\noindent\textit{#1.}\ }{\qed\par}\fi
\usepackage{lastpage}
\usepackage{hyperref}
\hypersetup{hidelinks}

% ============================================================
% Palette « Pop » OnBuch+ — mêmes hex que la classe P (plus_ui.dart)
% ============================================================
\definecolor{poporange}{HTML}{F59321}
\definecolor{popdark}{HTML}{E07A0C}
\definecolor{popcream}{HTML}{FFF6E8}
\definecolor{popink}{HTML}{1C1714}
\definecolor{poppurple}{HTML}{7A5AE0}
\definecolor{popgreen}{HTML}{2E9E6A}
\definecolor{poppink}{HTML}{E0457B}
\definecolor{popblue}{HTML}{2F7DE1}
\definecolor{popgold}{HTML}{C98A12}
\definecolor{popmuted}{HTML}{8A7F76}
\definecolor{popline}{HTML}{EADFD2}
\definecolor{popgray}{HTML}{8A7F76}
\colorlet{popgrey}{popgray}
% Alias tolérants (le modèle nomme parfois les couleurs différemment)
\colorlet{popwhite}{white}
\colorlet{popblack}{popink}
\colorlet{popred}{poppink}
\colorlet{popyellow}{popgold}
% Teintes claires (fonds de tableaux/encadrés) — le modèle nomme parfois
% les couleurs avec un suffixe L (ex. popblueL) sans les avoir définies.
\colorlet{poporangeL}{poporange!20}
\colorlet{popdarkL}{popdark!20}
\colorlet{poppurpleL}{poppurple!20}
\colorlet{popgreenL}{popgreen!20}
\colorlet{poppinkL}{poppink!20}
\colorlet{popblueL}{popblue!20}
\colorlet{popgoldL}{popgold!20}
\colorlet{popredL}{popred!20}
\colorlet{popgrayL}{popgray!20}
% Teintes claires pour fonds de tableaux / zones
\colorlet{popblueL}{popblue!10!white}
\colorlet{poporangeL}{poporange!14!white}
\colorlet{popgreenL}{popgreen!12!white}
\colorlet{poppurpleL}{poppurple!12!white}
\colorlet{poppinkL}{poppink!12!white}
\colorlet{popgoldL}{popgold!14!white}

\pagecolor{popcream}

% ============================================================
% Typographie
% ============================================================
\defaultfontfeatures{Ligatures=TeX}
\setmainfont{PlusJakartaSans-Medium.ttf}[Path=../../../fonts/,BoldFont=PlusJakartaSans-Bold.ttf]
% Maths en sans-serif (Fira Math) avec chiffres et lettres droites de Plus
% Jakarta, pour que les formules se fondent dans le texte.
\usepackage[math-style=ISO,bold-style=ISO]{unicode-math}
\setmathfont{FiraMath-Regular.otf}[Path=../../../fonts/]
\setmathfont{PlusJakartaSans-Medium.ttf}[Path=../../../fonts/,range={up/num,up/{Latin,latin},\mathup}]
\usepackage{icomma} % virgule décimale sans espace en mode math
% Caractères mathématiques Unicode absents des polices de texte (Plus Jakarta,
% Archivo Black) : rendus par leur équivalent mathématique, en texte comme en
% formule — sinon ils s'affichent en carré vide (ex. « Arithmétique dans ℤ »).
\usepackage{newunicodechar}
\newunicodechar{ℝ}{\ensuremath{\mathbb{R}}}
\newunicodechar{ℤ}{\ensuremath{\mathbb{Z}}}
\newunicodechar{ℂ}{\ensuremath{\mathbb{C}}}
\newunicodechar{ℕ}{\ensuremath{\mathbb{N}}}
\newunicodechar{ℚ}{\ensuremath{\mathbb{Q}}}
\newunicodechar{𝔻}{\ensuremath{\mathbb{D}}}
\newunicodechar{α}{\ensuremath{\alpha}}
\newunicodechar{β}{\ensuremath{\beta}}
\newunicodechar{γ}{\ensuremath{\gamma}}
\newunicodechar{δ}{\ensuremath{\delta}}
\newunicodechar{ε}{\ensuremath{\varepsilon}}
\newunicodechar{θ}{\ensuremath{\theta}}
\newunicodechar{λ}{\ensuremath{\lambda}}
\newunicodechar{μ}{\ensuremath{\mu}}
\newunicodechar{σ}{\ensuremath{\sigma}}
\newunicodechar{τ}{\ensuremath{\tau}}
\newunicodechar{φ}{\ensuremath{\varphi}}
\newunicodechar{ψ}{\ensuremath{\psi}}
\newunicodechar{ω}{\ensuremath{\omega}}
\newunicodechar{Σ}{\ensuremath{\Sigma}}
% majuscules produites par \MakeUppercase dans les titres (α-aminés -> Α)
\newunicodechar{Α}{\ensuremath{\alpha}}
\newunicodechar{Β}{\ensuremath{\beta}}
\newunicodechar{Γ}{\ensuremath{\Gamma}}
\newunicodechar{Δ}{\ensuremath{\Delta}}
\newunicodechar{Ε}{\ensuremath{\varepsilon}}
\newunicodechar{Θ}{\ensuremath{\Theta}}
\newunicodechar{Λ}{\ensuremath{\Lambda}}
\newunicodechar{Μ}{\ensuremath{\mu}}
\newunicodechar{Τ}{\ensuremath{\tau}}
\newunicodechar{Φ}{\ensuremath{\Phi}}
\newunicodechar{Ψ}{\ensuremath{\Psi}}
\newunicodechar{Ω}{\ensuremath{\Omega}}
\newunicodechar{↦}{\ensuremath{\mapsto}}
\newunicodechar{∈}{\ensuremath{\in}}
\newunicodechar{∉}{\ensuremath{\notin}}
\newunicodechar{∧}{\ensuremath{\wedge}}
\newunicodechar{⇔}{\ensuremath{\Leftrightarrow}}
\newunicodechar{⟺}{\ensuremath{\Longleftrightarrow}}
\newunicodechar{⇒}{\ensuremath{\Rightarrow}}
\newunicodechar{⟹}{\ensuremath{\Longrightarrow}}
\newunicodechar{∘}{\ensuremath{\circ}}
\newunicodechar{∪}{\ensuremath{\cup}}
\newunicodechar{∩}{\ensuremath{\cap}}
\newunicodechar{≡}{\ensuremath{\equiv}}
\newunicodechar{⊂}{\ensuremath{\subset}}
\newunicodechar{⊥}{\ensuremath{\perp}}
\newunicodechar{∃}{\ensuremath{\exists}}
\newunicodechar{∀}{\ensuremath{\forall}}
\newunicodechar{∛}{\ensuremath{\sqrt[3]{\,}}}
\newunicodechar{∜}{\ensuremath{\sqrt[4]{\,}}}
\newunicodechar{⃗}{\ensuremath{{}^{\to}}}
\newunicodechar{ⁿ}{\ifmmode^{n}\else\textsuperscript{\ensuremath{n}}\fi}
\newunicodechar{ˣ}{\ifmmode^{x}\else\textsuperscript{\ensuremath{x}}\fi}
\newunicodechar{ᵇ}{\ifmmode^{b}\else\textsuperscript{\ensuremath{b}}\fi}
\newunicodechar{ʳ}{\ifmmode^{r}\else\textsuperscript{\ensuremath{r}}\fi}
\newunicodechar{ᵘ}{\ifmmode^{u}\else\textsuperscript{\ensuremath{u}}\fi}
\newunicodechar{ʸ}{\ifmmode^{y}\else\textsuperscript{\ensuremath{y}}\fi}
\newunicodechar{ᵃ}{\ifmmode^{a}\else\textsuperscript{\ensuremath{a}}\fi}
\newunicodechar{ᵉ}{\ifmmode^{e}\else\textsuperscript{\ensuremath{e}}\fi}
\newunicodechar{ᵏ}{\ifmmode^{k}\else\textsuperscript{\ensuremath{k}}\fi}
\newunicodechar{ᵖ}{\ifmmode^{p}\else\textsuperscript{\ensuremath{p}}\fi}
\newunicodechar{ᴺ}{\ifmmode^{N}\else\textsuperscript{\ensuremath{N}}\fi}
\newunicodechar{ᶜ}{\ifmmode^{c}\else\textsuperscript{\ensuremath{c}}\fi}
\newunicodechar{ʰ}{\ifmmode^{h}\else\textsuperscript{\ensuremath{h}}\fi}
\newunicodechar{ᵠ}{\ifmmode^{q}\else\textsuperscript{\ensuremath{q}}\fi}
\newunicodechar{ₙ}{\ifmmode_{n}\else\textsubscript{\ensuremath{n}}\fi}
\newunicodechar{ₐ}{\ifmmode_{a}\else\textsubscript{\ensuremath{a}}\fi}
\newunicodechar{ᵢ}{\ifmmode_{i}\else\textsubscript{\ensuremath{i}}\fi}
\newunicodechar{ᵣ}{\ifmmode_{r}\else\textsubscript{\ensuremath{r}}\fi}
\newunicodechar{ₖ}{\ifmmode_{k}\else\textsubscript{\ensuremath{k}}\fi}
\newunicodechar{ᵦ}{\ifmmode_{\beta}\else\textsubscript{\ensuremath{\beta}}\fi}
\newunicodechar{ₓ}{\ifmmode_{x}\else\textsubscript{\ensuremath{x}}\fi}
\newfontfamily\popdisplay{ArchivoBlack-Regular.ttf}[Path=../../../fonts/]
\newfontfamily\poplogo{SpaceGrotesk-Bold.ttf}[Path=../../../fonts/]
\newfontfamily\popsemi{PlusJakartaSans-SemiBold.ttf}[Path=../../../fonts/]
\newfontfamily\popxbold{PlusJakartaSans-ExtraBold.ttf}[Path=../../../fonts/]
\newfontfamily\popxboldls{PlusJakartaSans-ExtraBold.ttf}[Path=../../../fonts/,LetterSpace=3.0]

\setlist[enumerate]{leftmargin=1.7em,itemsep=5pt,topsep=4pt}
\setlist[itemize]{leftmargin=1.7em,itemsep=4pt,topsep=4pt,label={\color{poporange}\textbullet}}
\setlength{\parindent}{0pt}
\setlength{\parskip}{5pt}
\renewcommand{\arraystretch}{1.25}

% pgfplots : style Pop par défaut
\pgfplotsset{
  every axis/.append style={axis line style={popink,line width=1pt},tick style={popink},
    grid style={popline,line width=0.5pt},label style={font=\small},tick label style={font=\footnotesize}},
  popcurve/.style={poporange,line width=1.8pt,no markers,smooth},
}
% Même style disponible dans un \draw TikZ simple (hors pgfplots)
\tikzset{popcurve/.style={poporange,line width=1.8pt}}
\tikzset{popgrid/.style={popline,very thin}}
\colorlet{popcurve}{poporange} % le nom du style est parfois utilisé comme couleur

% ============================================================
% Pill / squiggle
% ============================================================
\newcommand{\poppill}[3]{%
  \tikz[baseline=-0.55ex]\node[draw=popink,line width=1.2pt,rounded corners=2.6pt,%
    fill=#2,inner xsep=6.5pt,inner ysep=3pt]{%
    \popxboldls\fontsize{7}{7}\selectfont\color{#3}\MakeUppercase{#1}};%
}
\newcommand{\popsquiggle}[1]{%
  \tikz[baseline=-0.4ex]{
    \draw[#1,line width=2.6pt,line cap=round]
      plot [smooth] coordinates {(0,0.09) (0.34,0.01) (0.68,0.21) (1.02,0.01) (1.36,0.10)};
  }%
}
% Mot-clé surligné (marqueur orange sous le texte)
\newcommand{\cle}[1]{{\popxbold\color{popdark}#1}}

% ============================================================
% En-tête / pied
% ============================================================
\pagestyle{fancy}
\fancyhf{}
\renewcommand{\headrulewidth}{0pt}
\renewcommand{\footrulewidth}{0pt}
\lhead{\raisebox{-0.15ex}{\poplogo\fontsize{13}{13}\selectfont\color{popink}OnBuch\color{poporange}+}}
\rhead{\poppill{\DOCMATIERE\ \textbullet\ \DOCNIVEAU\ \textbullet\ Leçon \DOCLECON}{white}{popink}}
\lfoot{\popxbold\fontsize{7.6}{7.6}\selectfont\color{popmuted}\MakeUppercase{Cours signé OnBuch+}\ \textbullet\ {\popsemi\DOCTITRE}}
\rfoot{\tikz[baseline=-0.6ex]\node[rounded corners=8.5pt,fill=popink,minimum height=17pt,minimum width=17pt,inner xsep=5pt,inner sep=0]{\popxbold\fontsize{8}{8}\selectfont\color{popcream}\thepage\,/\,\pageref*{LastPage}};}

% ============================================================
% Titres
% ============================================================
\newcommand{\chaptitre}[2]{%
  \begin{tcolorbox}[enhanced,colback=poporange,colframe=popink,boxrule=2.6pt,arc=11pt,
      drop shadow=popink,left=15pt,right=15pt,top=12pt,bottom=11pt,before skip=3pt,after skip=18pt]
    \poppill{#2}{popink}{popcream}\par\vspace{9pt}
    {\raggedright\hyphenpenalty=10000\exhyphenpenalty=10000\popdisplay\fontsize{22}{24}\selectfont\color{popcream}\MakeUppercase{#1}\par}
  \end{tcolorbox}
}
% Section numérotée : pastille orange avec le numéro + titre Archivo
\newcounter{popsec}
\newcommand{\coursec}[1]{%
  \stepcounter{popsec}\setcounter{popsub}{0}%
  \par\vspace{16pt}\noindent
  \begin{minipage}{\linewidth}
  \tikz[baseline=-0.7ex]\node[draw=popink,line width=1.6pt,fill=poporange,rounded corners=4pt,minimum size=22pt,inner sep=0]{\popdisplay\fontsize{12}{12}\selectfont\color{popcream}\thepopsec};%
  \hspace{8pt}{\popdisplay\fontsize{15}{17}\selectfont\color{popink}\MakeUppercase{#1}}\par
  \vspace{4pt}\noindent{\color{poporange}\rule{36pt}{3.6pt}}
  \end{minipage}\par\nopagebreak\vspace{8pt}
}
\newcounter{popsub}
\newcommand{\courssub}[1]{%
  \stepcounter{popsub}%
  \par\vspace{9pt}\noindent{\popxbold\fontsize{11.5}{13}\selectfont\color{popink}{\color{poporange}\thepopsec.\thepopsub}\hspace{6pt}#1}\par\nopagebreak\vspace{4pt}
}

% ============================================================
% Boîtes pédagogiques — même recette : fond blanc, bordure épaisse colorée,
% ombre dure encre, badge plein. Titre optionnel [#1].
% ============================================================
\tcbset{popbase/.style={breakable,enhanced,colback=white,boxrule=2.3pt,arc=9pt,
  shadow={3pt}{-3pt}{0pt}{popink},left=14pt,right=14pt,top=11pt,bottom=11pt,before skip=11pt,after skip=15pt,
  fontupper=\popsemi\fontsize{10.6}{14.5}\selectfont\color{popink},
  fontlower=\popsemi\fontsize{10.6}{14.5}\selectfont\color{popink}}}
% Coche dessinée (le glyphe ✓ n'existe pas dans les polices du document)
\newcommand{\popcheck}[1]{\tikz[baseline=-0.5ex]\draw[#1,line width=1.6pt,line cap=round,line join=round](-0.13,0.0)--(-0.04,-0.09)--(0.13,0.1);}
\providecommand{\coche}{\popcheck{popgreen}}
\providecommand{\resultat}[1]{\par\smallskip\noindent\fcolorbox{popgreen}{popgreenL}{\parbox{\dimexpr\linewidth-2\fboxsep-2\fboxrule\relax}{\textbf{Résultat.} #1}}\par\smallskip}
\newcommand{\popboxhead}[3]{\poppill{#1}{#2}{white}\ifx\relax#3\relax\else\hspace{6pt}{\popxbold\fontsize{10.6}{12}\selectfont\color{popink}#3}\fi\par\vspace{7pt}}

\newtcolorbox{aretenir}[1][]{popbase,colframe=popgreen,before upper={\popboxhead{À retenir}{popgreen}{#1}}}
\newtcolorbox{exemplebox}[1][]{popbase,colframe=poporange,before upper={\popboxhead{Exemple}{poporange}{#1}}}
\newtcolorbox{definition}[1][]{popbase,colframe=popblue,before upper={\popboxhead{Définition}{popblue}{#1}}}
\newtcolorbox{propriete}[1][]{popbase,colframe=poppurple,before upper={\popboxhead{Propriété}{poppurple}{#1}}}
\newtcolorbox{methode}[1][]{popbase,colframe=popgold,before upper={\popboxhead{Méthode}{popgold}{#1}}}
\newtcolorbox{attention}[1][]{popbase,colframe=poppink,before upper={\popboxhead{Attention}{poppink}{#1}}}
\newtcolorbox{experience}[1][]{popbase,colframe=popblue,colback=popblueL,before upper={\popboxhead{Expérience}{popblue}{#1}}}
% « Le savais-tu ? » : carte violette pleine (contexte, culture, Cameroun)
\newtcolorbox{savaistu}[1][]{popbase,colback=poppurple,colframe=popink,
  fontupper=\popsemi\fontsize{10.4}{14.2}\selectfont\color{white},
  before upper={\poppill{Le savais-tu\,?}{white}{poppurple}\ifx\relax#1\relax\else\hspace{6pt}{\popxbold\color{white}#1}\fi\par\vspace{7pt}}}
% Exercice résolu : énoncé en haut, \tcblower puis la solution sur fond crème
\newcounter{popexo}\newcounter{popexr}
\newtcolorbox{exoresolu}[1][]{popbase,colframe=poporange,colbacklower=popcream,
  segmentation style={popink,line width=1.2pt,dashed},
  before upper={\stepcounter{popexr}\popboxhead{Exercice résolu \thepopexr}{poporange}{#1}},
  before lower={\poppill{Solution}{popink}{popcream}\par\vspace{6pt}}}
% Exercice (non résolu en place) — la correction est dans la partie Corrigés
\newtcolorbox{exercice}[1][]{popbase,colframe=popink,
  before upper={\stepcounter{popexo}\popboxhead{Exercice \thepopexo}{popink}{#1}}}
% Corrigé
\newtcolorbox{corrige}[1][]{popbase,colframe=popgreen,colback=popgreenL,
  before upper={\popboxhead{Corrigé}{popgreen}{#1}}}
% Objectifs / prérequis / bilan
\newtcolorbox{objectifs}{popbase,colframe=popink,colback=poporangeL,
  before upper={\popboxhead{Objectifs de la leçon}{popink}{}}}
\newtcolorbox{prerequis}{popbase,colframe=popink,colback=popblueL,
  before upper={\popboxhead{Prérequis}{popblue}{}}}
\newtcolorbox{bilan}{popbase,colframe=popink,colback=white,boxrule=2.8pt,
  before upper={\popboxhead{Fiche bilan}{poporange}{L'essentiel à connaître pour l'examen}}}
% Figure encadrée avec légende
\newtcolorbox{popfigure}[1][]{enhanced,breakable=false,colback=white,colframe=popline,boxrule=1.4pt,arc=8pt,
  left=10pt,right=10pt,top=10pt,bottom=8pt,before skip=10pt,after skip=12pt,halign=center,
  lower separated=false,halign lower=center,
  fontlower=\popsemi\fontsize{9}{11.5}\selectfont\color{popmuted},
  #1}
\newcommand{\legende}[1]{\par\vspace{6pt}{\popsemi\fontsize{9}{11.5}\selectfont\color{popmuted}#1}}

% ============================================================
% Signature OnBuch+
% ============================================================
\newcommand{\poplogoinline}[1]{{\poplogo\fontsize{#1}{#1}\selectfont\color{popink}OnBuch\color{poporange}+}}
\newcommand{\signatureonbuch}{%
  \begin{tcolorbox}[enhanced,colback=popink,colframe=popink,boxrule=2.2pt,arc=10pt,
      drop shadow=poporange,left=14pt,right=14pt,top=10pt,bottom=10pt,before skip=0pt,after skip=0pt]
    \tikz[baseline=-0.7ex]\node[circle,fill=popgreen,draw=popcream,line width=1.3pt,minimum size=22pt,inner sep=0]{\popcheck{white}};%
    \hspace{9pt}\begin{minipage}[c]{0.62\linewidth}
      {\popxbold\fontsize{12}{13}\selectfont\color{popcream}Signé\ \textbullet\ {\poplogo OnBuch\color{poporange}+}}\par\vspace{2pt}
      {\raggedright\popsemi\fontsize{8}{10.5}\selectfont\color{popline}\MakeUppercase{Équipe pédagogique OnBuch+}\\\MakeUppercase{Conforme au programme officiel MINESEC}\par}
    \end{minipage}\hfill
    {\poplogo\fontsize{22}{22}\selectfont\color{popcream}OnBuch\color{poporange}+}
  \end{tcolorbox}%
}

% ============================================================
% Page de garde
% #1 titre · #2 matière · #3 niveau · #4 module · #5 n° leçon · #6 durée ·
% #7 description / objectifs (texte ou itemize)
% ============================================================
\newcommand{\pagegarde}[7]{%
  \thispagestyle{empty}
  \newgeometry{a4paper,margin=1.7cm,top=1.6cm,bottom=1.4cm}%
  \begin{tikzpicture}[remember picture,overlay]
    \fill[poporange!18] ([xshift=-3.2cm,yshift=-3.6cm]current page.north east) circle (4.6cm);
    \fill[poppurple!14] ([xshift=2.4cm,yshift=7.6cm]current page.south west) circle (3.8cm);
  \end{tikzpicture}%
  {\poplogo\fontsize{30}{30}\selectfont\color{popink}OnBuch\color{poporange}+}\hfill
  \raisebox{0.6ex}{\poppill{Cours signé \textbullet\ Premium}{popink}{popcream}}\par
  \vspace{1pt}\popsquiggle{poppurple}\par
  \vspace{8pt}
  \begin{tcolorbox}[enhanced,colback=poporange,colframe=popink,boxrule=2.8pt,arc=13pt,
      drop shadow=popink,left=18pt,right=18pt,top=12pt,bottom=13pt,after skip=10pt]
    \poppill{#2 \textbullet\ #3}{popink}{popcream}\hspace{5pt}\poppill{#4}{white}{popink}\par\vspace{8pt}
    {\popdisplay\fontsize{14}{14}\selectfont\color{popink}LEÇON #5}\par\vspace{4pt}
    {\raggedright\hyphenpenalty=10000\exhyphenpenalty=10000\popdisplay\fontsize{25}{27}\selectfont\color{popcream}\MakeUppercase{#1}\par}
    \vspace{8pt}
    {\popxbold\fontsize{9.5}{9.5}\selectfont\color{popink}\MakeUppercase{Durée conseillée : #6}}
  \end{tcolorbox}
  \begin{tcolorbox}[enhanced,colback=white,colframe=popink,boxrule=2.2pt,arc=10pt,
      drop shadow=popink,left=16pt,right=16pt,top=10pt,bottom=10pt,after skip=8pt]
    \poppill{Dans cette leçon}{poppurple}{white}\par\vspace{5pt}
    \popsemi\fontsize{9.3}{12.6}\selectfont\color{popink}#7
  \end{tcolorbox}
  \noindent\begin{minipage}[t]{0.487\linewidth}
    \begin{tcolorbox}[enhanced,colback=popgreen,colframe=popink,boxrule=2.2pt,arc=10pt,
        drop shadow=popink,left=11pt,right=11pt,top=10pt,bottom=10pt,height=3.1cm,valign=center]
      \tcbox[colback=white,colframe=popink,boxrule=1.3pt,arc=5pt,boxsep=0pt,nobeforeafter,
        left=3pt,right=3pt,top=3pt,bottom=3pt,valign=center]{\qrcode[height=1.9cm]{https://play.google.com/store/apps/details?id=com.luvvix.onbuch&hl=fr}}%
      \hspace{8pt}\begin{minipage}[c]{0.5\linewidth}
        {\popdisplay\fontsize{11}{12.5}\selectfont\color{white}\MakeUppercase{Télécharge OnBuch}}\par\vspace{4pt}
        {\raggedright\popsemi\fontsize{8.4}{11}\selectfont\color{white}Gratuit \textbullet\ Google Play\\Tuteur IA, annales, concours, résultats\par}
      \end{minipage}
    \end{tcolorbox}
  \end{minipage}\hfill
  \begin{minipage}[t]{0.487\linewidth}
    \begin{tcolorbox}[enhanced,colback=poppurple,colframe=popink,boxrule=2.2pt,arc=10pt,
        drop shadow=popink,left=11pt,right=11pt,top=10pt,bottom=10pt,height=3.1cm,valign=center]
      \tikz[baseline=-0.7ex]\node[circle,fill=white,draw=popink,line width=1.3pt,minimum size=1.6cm,inner sep=0]{\popdisplay\fontsize{24}{24}\selectfont\color{poppurple}+};%
      \hspace{8pt}\begin{minipage}[c]{0.6\linewidth}
        {\popdisplay\fontsize{11}{12.5}\selectfont\color{white}\MakeUppercase{Passe à OnBuch+}}\par\vspace{4pt}
        {\raggedright\popsemi\fontsize{8.4}{11}\selectfont\color{white}Tous les cours signés, Tuteur IA illimité, fiches PDF — onglet OnBuch+ de l'app\par}
      \end{minipage}
    \end{tcolorbox}
  \end{minipage}
  \vspace{8pt}
  \popcontenu
  \vfill
  \signatureonbuch
  \restoregeometry
  \newpage
}

% Rangée « Ce cours contient » (page de garde)
\newcommand{\poptile}[3]{% #1 couleur #2 gros texte #3 légende
  \begin{tcolorbox}[enhanced,colback=#1,colframe=popink,boxrule=2pt,arc=9pt,shadow={2.5pt}{-2.5pt}{0pt}{popink},
      left=6pt,right=6pt,top=9pt,bottom=9pt,halign=center,valign=center,height=1.75cm,width=\linewidth,nobeforeafter]
    {\popdisplay\fontsize{12.5}{13}\selectfont\color{white}\MakeUppercase{#2}}\par\vspace{3pt}
    {\centering\popsemi\fontsize{7.8}{9.5}\selectfont\color{white}#3\par}
  \end{tcolorbox}}
\newcommand{\popcontenu}{%
  \par\noindent{\popxboldls\fontsize{7.5}{7.5}\selectfont\color{popmuted}CE COURS SIGNÉ CONTIENT}\par\vspace{7pt}
  \noindent\begin{minipage}[t]{0.235\linewidth}\poptile{popblue}{Cours}{complet, illustré, pas à pas}\end{minipage}\hfill
  \begin{minipage}[t]{0.235\linewidth}\poptile{poporange}{Méthodes}{et exercices résolus}\end{minipage}\hfill
  \begin{minipage}[t]{0.235\linewidth}\poptile{poppink}{Exercices}{type examen + corrigés}\end{minipage}\hfill
  \begin{minipage}[t]{0.235\linewidth}\poptile{popgreen}{Fiche bilan}{l'essentiel à retenir}\end{minipage}\par}

% Sommaire
\newtcolorbox{sommaire}{enhanced,colback=white,colframe=popink,boxrule=2.2pt,arc=10pt,
  drop shadow=popink,left=18pt,right=18pt,top=13pt,bottom=13pt,before skip=6pt,after skip=20pt,
  before upper={\poppill{Sommaire}{poppurple}{white}\par\vspace{9pt}\popsemi\fontsize{10.6}{14.5}\selectfont\color{popink}}}

% Signature de fin de cours — poussée tout en bas de la dernière page
\newcommand{\findecours}{%
  \par\vfill\nopagebreak
  \begin{center}\popsquiggle{poporange}\end{center}\vspace{4pt}
  \signatureonbuch
}
\AtBeginDocument{\def\llbracket{\mathopen{[\mkern-3.5mu[}}\def\rrbracket{\mathclose{]\mkern-3.5mu]}}} % intervalles d'entiers (glyphe ⟦ absent de la police)
% Glyphes absents de Fira Math : redéfinis à partir de glyphes disponibles
\AtBeginDocument{%
  \def\setminus{\mathbin{\backslash}}\let\smallsetminus\setminus
  \def\vdots{\mathord{\vbox{\baselineskip3.2pt\lineskiplimit0pt\kern4pt\hbox{.}\hbox{.}\hbox{.}}}}%
  \def\ddots{\mathinner{\mkern1mu\raise6.5pt\hbox{.}\mkern2mu\raise3.6pt\hbox{.}\mkern2mu\raise0.7pt\hbox{.}\mkern1mu}}%
  \def\triangle{\mathord{\scalebox{0.9}{$\symup{\Delta}$}}}%
  \def\nabla{\mathord{\raisebox{\depth}{\scalebox{1}[-1]{$\symup{\Delta}$}}}}%
  \def\checkmark{\popcheck{popdark}}%
  \def\star{\mathbin{\ast}}\def\bigstar{\mathord{\ast}}%
  \def\diamond{\mathbin{\scalebox{0.6}{$\Diamond$}}}%
  \def\bigcup{\mathop{\mathchoice{\raisebox{-0.3em}{\scalebox{1.7}{$\cup$}}}{\scalebox{1.25}{$\cup$}}{\cup}{\cup}}\displaylimits}%
  \def\bigcap{\mathop{\mathchoice{\raisebox{-0.3em}{\scalebox{1.7}{$\cap$}}}{\scalebox{1.25}{$\cap$}}{\cap}{\cap}}\displaylimits}%
  \def\lesseqqgtr{\mathrel{\lessgtr}}\def\bigoplus{\mathop{\oplus}\displaylimits}%
}
\providecommand{\ssi}{si et seulement si} % abréviation fréquente
\AtBeginDocument{\providecommand{\wideparen}{\overparen}} % arc de cercle

%% FIGFIX-BEGIN v3 — correctifs globaux des figures (docs/onbuch-generation/tools/figfix)
\usepackage{adjustbox}\usepackage{etoolbox}
% toute figure TikZ/pgfplots plus large que la ligne est réduite à la largeur disponible
\BeforeBeginEnvironment{tikzpicture}{\begin{adjustbox}{max width=\linewidth}}
\AfterEndEnvironment{tikzpicture}{\end{adjustbox}}
% légendes pgfplots lisibles par-dessus les courbes
\pgfplotsset{every axis/.append style={legend style={fill=white,fill opacity=0.92,text opacity=1,draw=black!15,inner sep=3pt}}}
% étiquettes d'arêtes (syntaxe edge["..."]) avec fond blanc
\tikzset{every edge quotes/.append style={fill=white,inner sep=1.2pt}}
% bulles de cartes mentales : texte à la ligne dans la bulle, bulle un peu plus grande
\tikzset{every concept/.append style={text width=2.0cm,align=center,minimum size=2.3cm,inner sep=2pt}}
%% FIGFIX-END
\begin{document}

\pagegarde{Le Cameroun dans la coopération multilatérale}{ECM}{Tle Tech}{ECM – classe de Terminale technique}{N03}{13 séances (fiche de progression harmonisée)}{Cette leçon analyse la stratégie diplomatique du Cameroun à travers ses engagements multilatéraux (ONU, UA, CEMAC) et bilatéraux (France, Chine, USA, etc.), évalue l'impact de l'aide au développement et de la gestion des conflits, pour comprendre les leviers et les limites de son insertion internationale.
\vspace{4pt}
\begin{multicols}{2}\small
\begin{itemize}
\item À la fin de cette leçon, je sais situer le Cameroun dans l'architecture des Nations Unies et évaluer sa contribution aux opérations de maintien de la paix.
\item Je sais analyser les enjeux de la coopération Cameroun-Union Européenne (Accords de Cotonou/Samoa) et les apports de la Francophonie, du Commonwealth et de l'OCI.
\item Je sais comparer les modèles de partenariat avec la France, le Royaume-Uni, la Chine, le Japon, les USA et le Brésil (investissements, aide, transfert de technologie).
\item Je sais expliquer le rôle du Cameroun au sein de l'Union Africaine et des mécanismes de paix en Afrique centrale (COPAX).
\item Je sais construire et commenter un croquis cartographique des flux de coopération (aide, investissements, migrations) du Cameroun.
\item Je sais rédiger une argumentation structurée sur les problèmes posés par la dépendance, l'endettement et la souveraineté dans les relations internationales du Cameroun.
\end{itemize}\end{multicols}}

\begin{sommaire}
\begin{multicols}{2}
\begin{enumerate}
\item 1. Le Cameroun à l'ONU : engagement universel et maintien de la paix
\item 2. Coopération régionale : Union Africaine, CEMAC/CEEAC et le COPAX
\item 3. Partenariats historiques et institutionnels : UE, Francophonie, Commonwealth, OCI
\item 4. Grands partenaires bilatéraux : France, Royaume-Uni, Chine, Japon, USA, Brésil
\item 5. L'aide au développement : acteurs, instruments, efficacité et conditionnalités (Dossier 3)
\item 6. Gestion des conflits et sécurité collective : Contribution camerounaise (TD2)
\item 7. Problèmes et défis des relations internationales du Cameroun (Dossier 5)
\item Activité d'intégration
\item Méthodes et exercices résolus
\item Exercices
\item Corrigés des exercices
\item Fiche bilan
\end{enumerate}
\end{multicols}
\end{sommaire}

\begin{objectifs}
\begin{itemize}
\item À la fin de cette leçon, je sais situer le Cameroun dans l'architecture des Nations Unies et évaluer sa contribution aux opérations de maintien de la paix.
\item Je sais analyser les enjeux de la coopération Cameroun-Union Européenne (Accords de Cotonou/Samoa) et les apports de la Francophonie, du Commonwealth et de l'OCI.
\item Je sais comparer les modèles de partenariat avec la France, le Royaume-Uni, la Chine, le Japon, les USA et le Brésil (investissements, aide, transfert de technologie).
\item Je sais expliquer le rôle du Cameroun au sein de l'Union Africaine et des mécanismes de paix en Afrique centrale (COPAX).
\item Je sais construire et commenter un croquis cartographique des flux de coopération (aide, investissements, migrations) du Cameroun.
\item Je sais rédiger une argumentation structurée sur les problèmes posés par la dépendance, l'endettement et la souveraineté dans les relations internationales du Cameroun.
\item Je sais mobiliser des données chiffrées (APD, IDE, dettes, effectifs militaires ONU) pour justifier une analyse géopolitique.
\item Je sais appliquer les concepts de coopération multilatérale/bilatérale, d'aide au développement et de sécurité collective à des cas concrets camerounais.
\end{itemize}
\end{objectifs}

\begin{prerequis}
\begin{itemize}
\item Notions d'organisation internationale (ONU, UA) et de droit international humanitaire (classe de 4ème / 1ère).
\item Géographie physique et humaine du Cameroun (ressources, population, position stratégique golfe de Guinée).
\item Histoire de la décolonisation et de la construction de l'État camerounais (réunification, politique d'union nationale).
\item Notions de base en économie : PIB, balance commerciale, dette publique, Investissements Directs Étrangers (IDE).
\item Maîtrise de la lecture de cartes thématiques (flux, choroplèthes) et de graphiques statistiques (diagrammes en barres, camemberts).
\end{itemize}
\end{prerequis}

\newpage

\coursec{Le Cameroun à l'ONU : engagement universel et maintien de la paix}

\courssub{Situation d'accroche et problématique}
Imagine un jeune ingénieur camerounais, fraîchement diplômé de l'\emph{École Nationale Supérieure Polytechnique (ENSP)} de Yaoundé. Il travaille sur le chantier colossal du barrage de \textbf{Nachtigal}, financé par un consortium où la \textbf{France} (via l'AFD et Proparco) et la \textbf{Banque Africaine de Développement (BAD)} jouent les premiers rôles. Pour se rendre sur le site, il conduit un véhicule assemblé au \textbf{Japon}, utilise un \textbf{smartphone} conçu aux \textbf{USA} intégrant des composants \textbf{chinois}, et consulte ses mails via une bourse d'études postulée auprès du \textbf{Commonwealth} ou de la \textbf{Francophonie}.

Cette scène banale révèle une vérité géopolitique majeure : le développement du Cameroun ne se joue pas en vase clos. Il est le fruit d'une \textbf{diplomatie multi-vectorielle} tissée patiemment depuis 1960. Comment le Cameroun, État pivot de l'Afrique centrale, utilise-t-il l'enceinte onusienne — cœur du multilatéralisme universel — pour transformer son potentiel en leviers concrets de sécurité, de formation et de prestige, tout en gérant les contradictions inhérentes à sa position de \emph{contributeur net de sécurité} mais \emph{bénéficiaire net d'aide} ?

\courssub{Historique : de l'adhésion (1960) aux grands organes onusiens}

\paragraph{Une entrée précoce dans la communauté des Nations.} Le 20 septembre 1960, à peine quatre mois après l'accession à la souveraineté internationale (1er janvier 1960), le Cameroun devient le \textbf{99e État membre} de l'ONU. Ce geste fondateur traduit la volonté des pères de la nation (Ahmadou Ahidjo, André-Marie Mbida) d'inscrire le jeune État dans la \textbf{légalité internationale} et de se doter d'une tribune pour défendre l'intégrité territoriale face aux velléités sécessionnistes (rébellion UPC, question du Bakassi).

\paragraph{Au Conseil de sécurité : deux mandats de membre non permanent.}
\begin{itemize}
    \item \textbf{1974--1975} : Premier mandat. Le Cameroun y défend la cause de la décolonisation (Sahara occidental, Afrique australe) et la \emph{Nouvel Ordre Économique International (NOEI)}.
    \item \textbf{2002--2003} : Second mandat, sous la présidence de \textbf{Paul Biya}. Point d'orgue : la gestion du contentieux \textbf{Bakassi} face au Nigeria. Le Cameroun obtient gain de cause à la \textbf{CIJ} (arrêt du 10 octobre 2002) et utilise son siège pour assurer le suivi de l'exécution de l'arrêt (Commission mixte Cameroun-Nigéria, \emph{Greentree Agreement} 2006).
\end{itemize}

\paragraph{À l'Assemblée générale et à la CIJ.}
\begin{itemize}
    \item \textbf{Assemblée générale} : Participation active aux votes de principe (décolonisation, Palestine, désarmement). Le Cameroun a fourni des vice-présidents de session et présidé la \emph{Commission des questions politiques spéciales et de la décolonisation} (4e commission).
    \item \textbf{Cour internationale de Justice (CIJ)} : Au-delà de l'affaire \emph{Cameroun c. Nigeria}, le Cameroun a été partie dans l'affaire \emph{Certaines activités nucléaires (Îles Marshall c. Cameroun)} (2016) et a déposé des déclarations d'intervention (affaire \emph{Tchad c. Libye} pour la bande d'Aozou). Des juristes camerounais (ex. \textbf{Me Henry Fossung}, \textbf{Pr. Maurice Kamto}) ont siégé ou plaidé, renforçant l'expertise juridique nationale.
\end{itemize}

\courssub{Dossier 2 : Participation aux OMP — Analyse quantitative et qualitative}

\begin{definition}[Opération de Maintien de la Paix (OMP)]
Une OMP est une opération déployée par l'ONU avec le consentement des parties principales, sous commandement onusien, pour \textbf{maintenir ou restaurer la paix et la sécurité internationales}.
\begin{itemize}
    \item \textbf{Chapitre VI de la Charte} : Règlement pacifique des différends (missions d'observation, interposition classique, consentement total).
    \item \textbf{Chapitre VII} : Action en cas de menace pour la paix, de rupture de la paix ou d'acte d'agression (imposition de la paix, \emph{peace enforcement}, mandat robuste, protection des civils \textbf{par tous les moyens nécessaires}).
\end{itemize}
Les OMP modernes sont multidimensionnelles : militaire, police, civile (droits de l'homme, élections, État de droit, DDR -- Désarmement, Démobilisation, Réinsertion).
\end{definition}

\begin{methode}[Calculer le « taux de contribution » d'un pays à une OMP]
Cet indicateur mesure l'engagement réel par rapport à la demande onusienne.
\begin{enumerate}
    \item Identifier l'\textbf{effectif autorisé} par la résolution du Conseil de sécurité (plafond \emph{force strength}).
    \item Identifier l'\textbf{effectif réellement déployé} par le pays (données mensuelles \emph{UN Force Generation Service}).
    \item Appliquer la formule : 
    \[ \text{Taux de contribution (\%)} = \frac{\text{Effectif national déployé}}{\text{Effectif total autorisé pour la mission}} \times 100 \]
    \item Pour un pays contributeur : $\frac{\text{Effectif national}}{\text{Effectif total de la mission}} \times 100$ donne la \textbf{part nationale dans la mission}.
\end{enumerate}
\end{methode}

\begin{attention}[Contribution financière vs Contribution en personnel]
Ne pas confondre deux logiques distinctes :
\begin{itemize}
    \item \textbf{Contribution financière (quotes-parts)} : Calculée selon la capacité de paiement (PIB, revenu par tête). Le Cameroun paie une quote-part minime (\textasciitilde 0,015\% du budget ordinaire, \textasciitilde 0,1\% du budget OMP). Il est souvent en \textbf{arriérés de paiement}.
    \item \textbf{Contribution en personnel (TCC/PCC)} : \emph{Troop Contributing Countries} / \emph{Police Contributing Countries}. Le Cameroun figure \textbf{régulièrement dans le Top 20 mondial} et \textbf{Top 5 africain} (souvent 3e--4e derrière Rwanda, Égypte, Ghana).
\end{itemize}
\textbf{Paradoxe} : Le Cameroun « paie » sa dette onusienne en \textbf{sang et en expertise} plutôt qu'en devises. Cela lui confère un \emph{capital diplomatique} (droit de regard, postes de commandement) disproportionné par rapport à son poids financier.
\end{attention}

\paragraph{Déploiements majeurs : effectifs, pertes, grades (Données 2000--2023).}

\begin{center}
\begin{tabularx}{\linewidth}{|l|X|X|X|}\hline
\textbf{Mission (Pays)} & \textbf{Période / Mandat} & \textbf{Contingent camerounais (Spécificités)} & \textbf{Pertes / Grades clés} \\ \hline
\textbf{MINUSTAH / BINUH (Haïti)} & 2004--2017 / 2019-- & \textbf{Police (FPU/IPU)} : Gendarmerie mobile, police judiciaire. Apport clé : rétablissement ordre public, formation PNH. & ~15 morts (dont séisme 2010). Commissaires divisionnaires chefs de composante police. \\ \hline
\textbf{MINUSCA (RCA)} & 2014--présent & \textbf{Infanterie (Bataillons), Gendarmerie (FPU), Hôpital niveau 2, Etat-major}. 1er contributeur africain par moments. & \textasciitilde 40 morts. \textbf{Gén. Martin Tumenta Chomu} (Cdt Force 2015-16), Gén. Balla Keita (Cdt Force 2021-23). \\ \hline
\textbf{MONUSCO (RDC)} & 2010--présent & \textbf{Infanterie, Génie, Hôpital, Police}. Déploiement Est (Nord-Kivu, Ituri). & \textasciitilde 25 morts. Officiers supérieurs à l'Etat-major (G5, J3). \\ \hline
\textbf{MINUSMA (Mali)} & 2013--2023 (retrait) & \textbf{Infanterie, Gendarmerie, Hôpital, Aviation (hélicos Mi-17/24)}. Retrait ordonné 2023. & \textasciitilde 30 morts (IED, attaques complexes). Gén. Kévin Bikandi (Cdt Secteur Nord). \\ \hline
\textbf{UNISFA (Abyéi)} & 2011--présent & \textbf{Infanterie, Observateurs militaires}. Zone frontalière Soudan/Sud-Soudan. & Faibles. Officiers observateurs chefs d'équipe. \\ \hline
\end{tabularx}
\end{center}

\begin{exemplebox}[Exemple chiffré : Situation 2023 (Sources : ONU DPPA / MINUSCA Force Generation)]
En 2023, le Cameroun déploie \textbf{\textasciitilde 2 500 éléments} (militaires, policiers, experts) :
\begin{itemize}
    \item \textbf{MINUSCA} : \textasciitilde 1 400 (Infanterie 2 bataillons, FPU Gendarmerie 2 unités, Hôpital Niv.2, Etat-major).
    \item \textbf{MONUSCO} : \textasciitilde 750 (Infanterie, Génie, Hôpital).
    \item \textbf{UNISFA} : \textasciitilde 200 (Infanterie, Observateurs).
    \item \textbf{BINUH / AUTRES} : \textasciitilde 150 (Policiers individuels, Experts).
\end{itemize}
\textbf{Part des femmes} : \textasciitilde 30\% (objectif ONU 2028 : 30\% militaires, 20\% policiers unitaires). Le Cameroun est \textbf{4e contributeur africain} et \textbf{12e--15e mondial}.
\end{exemplebox}

\begin{savaistu}[Le Général Martin Tumenta Chomu : un Camerounais à la tête d'une Force onusienne]
De 2015 à 2016, le \textbf{Général de division Martin Tumenta Chomu} a servi comme \textbf{Commandant de la Force (Force Commander)} de la \textbf{MINUSCA} (RCA). Il commandait alors \textasciitilde 12 000 soldats de 40 nationalités. Son mandat a été marqué par la réorganisation de la force face aux groupes armés (anti-balaka, ex-Séléka) et la protection des civils à Bangui et en province. Cette nomination illustre la \textbf{reconnaissance du professionnalisme de l'armée camerounaise} au plus haut niveau onusien. D'autres généraux (Balla Keita, Kévin Bikandi) ont tenu des postes de \emph{Deputy Force Commander} ou \emph{Sector Commanders}.
\end{savaistu}

\paragraph{Spécificités de l'apport camerounais : la « touche » nationale.}
\begin{itemize}
    \item \textbf{Gendarmerie nationale} : Déploie des \textbf{FPU (Formed Police Units)} très prisées pour le maintien de l'ordre, la protection de sites sensibles et le mentorat des polices locales (Haïti, RCA, Mali). Leur doctrine \emph{« police de proximité »} est exportée.
    \item \textbf{Police nationale} : Fournit des \textbf{IPU (Individual Police Officers)} : commissaires, officiers de police judiciaire, experts en criminalistique, chefs de composante police.
    \item \textbf{Service de Santé des Armées} : Déploie des \textbf{Hôpitaux de Niveau 2} (chirurgie d'urgence, imagerie, laboratoire) à Bangui (MINUSCA), Goma (MONUSCO), Gao (ex-MINUSMA). Référence pour les Casques bleus \emph{et} populations locales.
    \item \textbf{Génie militaire} : Déminage, ouverture de pistes, construction de camps (MONUSCO, MINUSCA).
\end{itemize}

\paragraph{Retombées pour l'outil de défense national.}
\begin{itemize}
    \item \textbf{Formation} : Création de l'\textbf{École de Maintien de la Paix (EMP)} à Yaoundé (2008), centre d'excellence régional (partenariat ONU, France, USA, Allemagne). Formation pré-déploiement obligatoire (CIMIC, DDH, Règles d'engagement, Genre/SEA).
    \item \textbf{Équipements} : Remboursement ONU (\emph{COE -- Contingent Owned Equipment}) : véhicules blindés (VBL, P4, Casspir), radios HF/VHF/SAT, matériel médical, génie. Modernisation du parc roulant et transmissions.
    \item \textbf{Expérience opérationnelle} : Combat asymétrique (IED, guérilla urbaine, protection civils), interopérabilité OTAN/ONU, gestion logistique en théâtre austère. Transférable aux théâtres nationaux (Boko Haram, NOSO).
    \item \textbf{Prestige / Postes} : Accès aux postes de direction (Force Cdr, Police Cdr, Directeurs de mission civils, Chefs de cabinet SRSG).
\end{itemize}

\paragraph{Limites et zones d'ombre.}
\begin{itemize}
    \item \textbf{Dépendance logistique} : L'ONU ne rembourse que le \emph{COE} validé. Les délais de paiement (parfois 12--18 mois) créent des trous de trésorerie pour le MINDEF. Maintenance lourde (hélicoptères, blindés) souvent externalisée ou défaillante.
    \item \textbf{Allégations d'exploitation et abus sexuels (EAS)} : Quelques cas isolés (Haïti 2010, RCA 2015, Mali 2018) ont terni l'image. Réponse : \textbf{Tolérance zéro}, cour martiale, formation pré-déploiement renforcée (module SEA), nomination d'officiers \emph{Focal Points} genre/SEA dans chaque contingent.
    \item \textbf{Coût d'opportunité} : \textasciitilde 2 500 hommes d'élite (bataillons, FPU, hôpital) projetés à l'extérieur alors que le pays fait face à \textbf{deux guerres internes simultanées} (Extrême-Nord : Boko Haram/ISWAP ; NOSO : séparatistes armés). Rotation des unités, fatigue du matériel, usure du personnel.
\end{itemize}

\begin{popfigure}
\centering
\begin{tikzpicture}[scale=0.85]
% Carte simplifiée Afrique Centrale / de l'Ouest
% Coordonnées approximatives (x,y) pour TikZ
\coordinate (CAM) at (0,0); % Yaoundé (origine relative)
\coordinate (BANGUI) at (-1.5, 2.5); % MINUSCA
\coordinate (GOMA) at (3.5, 1.5);   % MONUSCO
\coordinate (GAO) at (-4.5, 1.0);  % MINUSMA (historique)
\coordinate (ABYEI) at (2.0, 4.0); % UNISFA
\coordinate (PORT_AU_PRINCE) at (-10, -2); % Haïti (hors cadre, on met une note)
% Contour très simplifié (polygone grossier)
\draw[popline, thick, fill=popcream!30] 
(-6,-3) -- (8,-3) -- (8,6) -- (-2,7) -- (-6,4) -- cycle;
% Pays hôtes
\node[circle, fill=popblue!20, draw=popblue, thick, minimum size=6mm] at (CAM) {\small \textbf{CMR}};
\node[anchor=south] at (CAM) {\tiny Yaoundé};
% Cercles proportionnels : rayon = sqrt(effectif/pi) * facteur. Facteur ~ 0.04 pour 1500 -> 1.5cm
% MINUSCA ~1400 -> r~1.3cm
\draw[popblue, thick, fill=popblue!30] (BANGUI) circle (1.3cm);
\node[popdark, font=\bfseries\small] at (BANGUI) {MINUSCA};
\node[popdark, font=\tiny] at ($(BANGUI)+(0,-1.6)$) {\textasciitilde 1 400 pers.};
% MONUSCO ~750 -> r~0.95cm
\draw[poporange, thick, fill=poporange!30] (GOMA) circle (0.95cm);
\node[popdark, font=\bfseries\small] at (GOMA) {MONUSCO};
\node[popdark, font=\tiny] at ($(GOMA)+(0,-1.2)$) {\textasciitilde 750 pers.};
% MINUSMA ~0 (retrait 2023) mais historique ~1000 -> r~1.1cm (trait pointillé)
\draw[popgreen, thick, dashed, fill=popgreen!10] (GAO) circle (1.1cm);
\node[popdark, font=\bfseries\small] at (GAO) {MINUSMA (fin 2023)};
\node[popdark, font=\tiny] at ($(GAO)+(0,-1.4)$) {Retrait ordonné};
% UNISFA ~200 -> r~0.5cm
\draw[poppurple, thick, fill=poppurple!30] (ABYEI) circle (0.5cm);
\node[popdark, font=\bfseries\small] at (ABYEI) {UNISFA};
\node[popdark, font=\tiny] at ($(ABYEI)+(0,-0.7)$) {\textasciitilde 200 pers.};
% Flèches de déploiement depuis Yaoundé
\draw[->, popdark, thick, opacity=0.6] (CAM) -- (BANGUI);
\draw[->, popdark, thick, opacity=0.6] (CAM) -- (GOMA);
\draw[->, popdark, thick, opacity=0.6] (CAM) -- (GAO);
\draw[->, popdark, thick, opacity=0.6] (CAM) -- (ABYEI);
% Légende manuelle
\node[anchor=north west, align=left, font=\small, fill=white, fill opacity=0.9, draw=popline, rounded corners] at (-6.5, 6.5) {
\textbf{Légende :} Cercles proportionnels aux effectifs 2023.\\
\textcolor{popblue}{■} MINUSCA (RCA) \textasciitilde 1 400\\
\textcolor{poporange}{■} MONUSCO (RDC) \textasciitilde 750\\
\textcolor{popgreen}{■} MINUSMA (Mali) Retrait 2023\\
\textcolor{poppurple}{■} UNISFA (Abyéi) \textasciitilde 200\\
\textcolor{popdark}{$\leftarrow$} Flux de déploiement depuis Yaoundé
};
\end{tikzpicture}
\legende{Carte de proportionnalité : Localisation et effectifs des contingents camerounais dans les OMP majeures (Situation 2023). Les cercles sont proportionnels aux effectifs déployés. La MINUSMA est représentée en pointillés pour marquer le retrait ordonné fin 2023.}
\end{popfigure}

\begin{popfigure}
\centering
\begin{tikzpicture}
\begin{axis}[
    ybar stacked,
    height=9cm, width=\linewidth,
    bar width=12mm,
    xlabel={Année},
    ylabel={Effectifs déployés (personnel militaire + police + experts)},
    xtick={2000,2005,2010,2015,2020,2023},
    xticklabels={2000,2005,2010,2015,2020,2023},
    legend style={at={(0.5,-0.2)},anchor=north,legend columns=-1, font=\small},
    ymin=0, ymax=3000,
    major grid style={draw=popline},
    axis lines=left,
    enlarge x limits=0.15,
    tick label style={font=\small},
    label style={font=\small},
]
% Données estimées (Militaires, Policiers, Experts/Observateurs)
% 2000: 300M, 100P, 50E = 450
% 2005: 600M, 300P, 100E = 1000 (Haïti, Burundi, RDC)
% 2010: 900M, 500P, 150E = 1550 (Haïti pic, RDC, Soudan)
% 2015: 1600M, 700P, 200E = 2500 (MINUSCA + MINUSMA pic)
% 2020: 1400M, 600P, 200E = 2200 (MINUSCA stable, MINUSMA haut, MONUSCO stable)
% 2023: 1600M, 700P, 200E = 2500 (MINUSCA haut, MONUSCO, MINUSMA 0)
\addplot[popblue, fill=popblue!60] coordinates {(2000,300) (2005,600) (2010,900) (2015,1600) (2020,1400) (2023,1600)}; \addlegendentry{Militaires (Troupes)}
\addplot[poporange, fill=poporange!60] coordinates {(2000,100) (2005,300) (2010,500) (2015,700) (2020,600) (2023,700)}; \addlegendentry{Policiers (FPU/IPU)}
\addplot[popgreen, fill=popgreen!60] coordinates {(2000,50) (2005,100) (2010,150) (2015,200) (2020,200) (2023,200)}; \addlegendentry{Experts / Observateurs}
\end{axis}
\end{tikzpicture}
\legende{Évolution du nombre de personnels camerounais déployés dans les OMP (2000--2023). Barres empilées : Militaries (bleu), Policiers (orange), Experts/Observateurs (vert). Pic 2015 (MINUSCA + MINUSMA simultanées), creux relatif 2020 (COVID, drawdown MINUSMA), remontée 2023 (renforcement MINUSCA/MONUSCO). Sources : ONU DPPA, Rapports annuels MINDEF.}
\end{popfigure}

\courssub{Vote, diplomatie et candidatures : alignements et autonomie}

\paragraph{Comportement de vote à l'Assemblée générale : le « non-alignement actif ».}
Le Cameroun vote selon une grille de lecture \textbf{« intérêt national + solidarité africaine + droit international »}.
\begin{itemize}
    \item \textbf{Décolonisation / Sahara occidental} : Soutien constant au droit à l'autodétermination (RASD membre UA), aligné sur la position de l'UA.
    \item \textbf{Question de Palestine} : Vote systématique pour les résolutions pro-palestiniennes (statut d'État observateur 2012, cessez-le-feu 2023), reconnaissance de l'État de Palestine (1988). Rejet de l'annexion de Jérusalem.
    \item \textbf{Climat / Environnement} : Membre du \emph{Groupe des 77 + Chine} et du \emph{Groupe africain}. Défense de la \textbf{responsabilité commune mais différenciée (RCPD)}, financement adaptation/pertes et préjudices (COP27/28). Signature Accord de Paris 2016.
    \item \textbf{Responsabilité de Protéger (R2P) / Intervention humanitaire} : \textbf{Prudence extrême}. Le Cameroun privilégie le \textbf{consentement de l'État hôte} (Chapitre VI/VII classique) et refuse tout précédent d'intervention unilatérale (Libye 2011 : abstention/réticence ; Syrie : non-ingérence). Crains pour sa propre souveraineté (NOSO, Bakassi).
\end{itemize}

\paragraph{Candidatures aux postes de direction : le lobbying diplomatique.}
Le Cameroun transforme sa contribution en troupes en \textbf{capital politique} pour placer ses nationaux :
\begin{itemize}
    \item \textbf{SG Adjoint / Chefs d'agences} : \textbf{Philippe Douste-Blazy} (Français, mais lobbying Cameroun pour OMD), \textbf{Jean-Pierre Ondoua} (OMS), \textbf{Joseph Dion Ngute} (SG adjoint hypothétique), \textbf{Moussa Oumarou} (OIT - Directeur général adjoint 2019).
    \item \textbf{Postes politiques (DPA, DPPA, OCHA)} : Diplomates camerounais (ex. \textbf{Anatole Fabien Marie Nkou}, \textbf{Samuel Mvondo Ayolo}) aux postes de \emph{Directeur de division}, \emph{Chefs de bureau régionaux}.
    \item \textbf{Stratégie} : « Package deal » : Contribution troupes + Soutien votes (droits de l'homme, budget) + Candidatures groupées (UA) $\rightarrow$ Postes.
\end{itemize}

\courssub{Exercice résolu : Analyse d'un scénario de déploiement}

\begin{exoresolu}[Calcul de taux de contribution et analyse capacitaire]
\textbf{Énoncé.} En janvier 2024, le Conseil de sécurité autorise par la résolution 27XX une nouvelle mission multidimensionnelle en \textbf{République Centrafricaine (MINUSCA-II)} avec un plafond de \textbf{14 000 militaires} et \textbf{3 000 policiers}. Le Cameroun propose : 2 bataillons d'infanterie (850 hommes chacun), 1 FPU Gendarmerie (140 hommes), 1 Hôpital Niveau 2 (65 personnes), 10 observateurs militaires.
\begin{enumerate}
    \item Calculer l'effectif total proposé par le Cameroun.
    \item Calculer la part du Cameroun dans le plafond militaire autorisé (\%).
    \item Si l'ONU ne valide qu'un bataillon (850) au lieu de deux, quel est le \emph{taux de satisfaction} de la demande camerounaise pour l'infanterie ?
    \item Discuter : Pourquoi l'ONU peut-elle refuser le 2e bataillon malgré le besoin ?
\end{enumerate}

\tcblower
\textbf{Solution détaillée.}
\begin{enumerate}
    \item \textbf{Effectif total proposé} :
    \[ 2 \times 850 \text{ (Inf)} + 140 \text{ (FPU)} + 65 \text{ (Hôpital)} + 10 \text{ (Obs)} = 1700 + 140 + 65 + 10 = \textbf{1 915 personnes}. \]
    \item \textbf{Part dans le plafond militaire} (seuls les bataillons comptent pour le plafond militaire « troops ») :
    \[ \frac{1700}{14000} \times 100 = \textbf{12,14 \%}. \]
    Le Cameroun fournirait \textasciitilde 1/8 de la force militaire totale. Poids majeur.
    \item \textbf{Taux de satisfaction infanterie} :
    \[ \frac{\text{Validé (850)}}{\text{Demandé (1700)}} \times 100 = \textbf{50 \%}. \]
    \item \textbf{Raisons du refus (analyse capacitaire)} :
    \begin{itemize}
        \item \textbf{Capacité d'absorption logistique} : La mission n'a pas les infrastructures (camps, eau, électricité, évacuation médicale) pour absorber 1 700 soldats camerounais d'un coup.
        \item \textbf{Diversité des contributeurs} : L'ONU cherche à éviter la domination d'un seul pays (risque de « contingent unique », barrière de langue, cohérence doctrinale). Elle préfère 2 bataillons de pays différents.
        \item \textbf{Disponibilité réelle} : L'état-major onusien (FGS) vérifie la \emph{readiness} (équipements, formation, durabilité). Si le 2e bataillon n'est pas « \emph{mission ready} » (ex. manque de blindés, formation SEA incomplète), il est mis en liste d'attente.
        \item \textbf{Contraintes nationales} : L'armée camerounaise doit garder des réserves pour les théâtres intérieurs (Boko Haram, NOSO). L'État-major Yaoundé peut lui-même limiter l'offre.
    \end{itemize}
\end{enumerate}
\end{exoresolu}

\courssub{Exercices d'application}

\begin{exercice}[Exercice 1 : Vocabulaire et concepts]
Définir précisément les termes suivants en les situant dans le cadre de la participation du Cameroun à l'ONU :
\begin{enumerate}
    \item TCC / PCC (Troop / Police Contributing Country).
    \item COE (Contingent Owned Equipment) et système de remboursement « wet/dry lease ».
    \item Règle d'engagement (ROE) vs Mandat robuste (Chapitre VII).
    \item SEA (Exploitation et Abus Sexuels) : mécanisme de prévention dans les contingents camerounais.
\end{enumerate}
\end{exercice}

\begin{exercice}[Exercice 2 : Analyse de document - Discours à la 78e AGNU (Sept 2023)]
\textit{Extrait (reconstruit) :} « Monsieur le Président, le Cameroun réaffirme son attachement au multilatéralisme... Nous payons le prix du sang à Bangui, Goma, Gao... Mais nous exigeons la réforme du Conseil de sécurité : l'Afrique a besoin de deux sièges permanents avec droit de veto. L'aide au développement ne saurait remplacer la justice climatique... »
\begin{enumerate}
    \item Identifier les trois registres de l'argumentaire (sécuritaire, institutionnel, économique).
    \item Expliquer la référence au « prix du sang » : à quelles missions fait allusion le locuteur ?
    \item Quelle est la position camerounaise sur la réforme du CSNU ? La relier à la candidature du Cameroun (soutien de la France ? de la Chine ?).
\end{enumerate}
\end{exercice}

\begin{exercice}[Exercice 3 : Cas pratique - Gestion d'une allégation SEA]
\textbf{Situation.} Un soldat du bataillon camerounais (MINUSCA, Bangui) est accusé par une ONG locale d'avoir eu une relation sexuelle avec une mineure de 16 ans en échange de vivres. L'enquête de la \emph{Conduct and Discipline Unit (CDU)} de la mission est ouverte.
\begin{enumerate}
    \item Citer les \textbf{3 piliers} de la politique ONU « Tolérance Zéro » (Prévention, Application de la loi, Assistance aux victimes).
    \item Décrire la procédure \textbf{disciplinaire et pénale} côté camerounais (Code de justice militaire, Cour martiale, rapatriement).
    \item Quelles mesures de \textbf{réparation pour la victime} (fonds d'affectation spéciale, prise en charge médicale/psycho-sociale) ?
    \item En tant que commandant de contingent, quelles \textbf{mesures immédiates de prévention} prenez-vous pour le reste du bataillon ?
\end{enumerate}
\end{exercice}

\begin{corrige}[Exercice 1]
\begin{enumerate}
    \item \textbf{TCC/PCC} : Pays contributeurs de troupes / de policiers. Le Cameroun est TCC majeur (infanterie, hôpital, génie) et PCC majeur (FPU Gendarmerie, IPU Police).
    \item \textbf{COE} : Équipement détenu par le contingent. \emph{Wet lease} : l'État fournisseur assure maintenance + personnel (ex. hélicoptères, hôpital). \emph{Dry lease} : l'ONU assure la maintenance (ex. véhicules blindés, tentes). Remboursement mensuel selon tarifs ONU (manuel COE).
    \item \textbf{ROE} : Règles d'engagement (directives tactiques : quand tirer, force proportionnelle). \textbf{Mandat robuste} : Base légale (Chap VII) autorisant l'usage de la force au-delà de la légitime défense (protection civils, neutralisation groupes armés).
    \item \textbf{SEA Prévention Cameroun} : Formation pré-déploiement obligatoire (module 3j), Code de conduite signé, \emph{Focal Point} genre/SEA par unité, interdiction d'alcool/relations transactionnelles, patrouilles mixtes avec police militaire, sensibilisation continue.
\end{enumerate}
\end{corrige}

\begin{corrige}[Exercice 2]
\begin{enumerate}
    \item \textbf{Sécuritaire} : « Prix du sang » (légitimité par le sacrifice). \textbf{Institutionnel} : Réforme CSNU (sièges permanents Afrique, veto). \textbf{Économique/Climatique} : Justice climatique, dette, ODD.
    \item \textbf{Missions} : MINUSCA (RCA, \textasciitilde 40 morts), MINUSMA (Mali, \textasciitilde 30 morts), MONUSCO (RDC, \textasciitilde 25 morts), BINUH (Haïti). Total \textasciitilde 100+ morts depuis 2014.
    \item \textbf{Position} : Consensus d'Ezulwini (UA) : 2 sièges permanents + 2 non-permanents + veto. Cameroun candidat naturel (poids démographique, économique, contribution OMP). Soutien France (historique), Chine (partenaire stratégique), mais opposition rivaux régionaux (Nigéria, Afrique du Sud). Lobbying au sein du « Groupe des 10 » (C10) africain.
\end{enumerate}
\end{corrige}

\begin{corrige}[Exercice 3]
\begin{enumerate}
    \item \textbf{3 piliers} : (1) Prévention (formation, dépistage, environnement protecteur). (2) Réponse / Application de la loi (enquête CDU, suspension, rapatriement, poursuite pénale). (3) Assistance victimes (médicale, psycho-sociale, juridique, fonds réparation ONU).
    \item \textbf{Procédure Cameroun} : Suspension immédiate des fonctions. Enquête CDU (ONU) + Enquête Police Militaire / Gendarmerie (National). Si preuves : Transfert au Cameroun pour jugement devant \textbf{Cour Martiale} (compétence extraterritoriale art. 6 Code Justice Militaire). Peines : Prison, radiation, dégradation. Rapatriement « \emph{repatriation on disciplinary grounds} » (fin de solde ONU).
    \item \textbf{Réparation} : Fonds d'affectation spéciale ONU (Victims' Assistance Fund). Prise en charge médicale (hôpital niveau 2 contingent ou structure locale), soutien psycho-social (ONG partenaires), aide à la scolarisation/formation professionnelle, compensation financière (si jugement définitif).
    \item \textbf{Mesures immédiates Commandant} : (1) Conférence « \emph{Town Hall} » bataillon : rappel tolérance zéro, conséquences carrière/famille. (2) Renforcement patrouilles \emph{Force Protection} / \emph{Community Engagement} hors base. (3) Contrôle sorties (permissionnaires) : système \emph{buddy}, couvre-feu, interdiction zones à risque. (4) Vérification formation SEA de tous (rattrapage si absent). (5) Désignation officiers \emph{Focal Points} par compagnie. (6) Signalement hiérarchique (MINDEF, Ambassadeur ONU) sous 24h.
\end{enumerate}
\end{corrige}

\coursec{Coopération régionale : Union Africaine, CEMAC/CEEAC et le COPAX}

Dans la section précédente, nous avons vu que le Cameroun s'engage sur la scène universelle à travers l'ONU et ses missions de maintien de la paix. Mais avant de peser au niveau mondial, un État doit d'abord compter dans son \emph{voisinage immédiat}. C'est exactement la logique de la coopération régionale : le Cameroun, par sa position géographique charnière entre l'Afrique de l'Ouest et l'Afrique centrale, par son poids économique et démographique, joue un rôle de \cle{pivot} dans son continent et dans sa sous-région. Cette section examine successivement son rôle dans l'Union Africaine, son insertion dans la CEMAC et la CEEAC, le fonctionnement du COPAX, la gestion des crises frontalières et, enfin, les dysfonctionnements de l'intégration régionale.

\courssub{Le Cameroun et l'Union Africaine : de l'OUA à l'UA}

\textbf{Intuition.} Après les indépendances africaines (années 1960), les jeunes États se retrouvent face à un dilemme : rester isolés et faibles, ou s'unir pour peser dans les affaires du monde. Le Cameroun fait très tôt le choix de l'unité africaine.

Le \cle{25 mai 1963}, à Addis-Abeba (Éthiopie), 32 chefs d'État africains signent la Charte de l'\cle{Organisation de l'Unité Africaine} (OUA). Le Cameroun, sous la présidence d'\cle{Ahmadou Ahidjo}, est membre fondateur. Ahidjo défend une ligne de \emph{non-alignement} et de coopération : il accueille d'ailleurs à Yaoundé, dès 1961, une conférence des pays dits « de Brazzaville » puis participe aux réunions préparatoires de l'unité continentale. L'OUA repose sur deux principes cardinaux : l'\cle{intangibilité des frontières héritées de la colonisation} (\emph{uti possidetis juris}) et la \cle{non-ingérence} dans les affaires intérieures des États.

En \cle{2002}, à Durban (Afrique du Sud), l'OUA devient l'\cle{Union Africaine} (UA), organisation plus ambitieuse inspirée du modèle européen : elle dispose d'une Commission, d'un Parlement panafricain, d'une Cour africaine de justice et des droits de l'homme et d'un Conseil de paix et de sécurité. Le président \cle{Paul Biya} participe activement aux sommets de l'UA ; le Cameroun a exercé plusieurs \cle{vice-présidences} du Bureau de l'Union (notamment en 2007, 2013, 2019), sans avoir encore occupé la présidence tournante (attribuée chaque année par rotation géographique).

\begin{definition}[Union Africaine]
L'\cle{Union Africaine} (UA) est une organisation intergouvernementale créée en 2002 à Durban, succédant à l'OUA (1963). Elle regroupe les 55 États africains et vise l'unité politique, l'intégration économique (ZLECAf) et la sécurité collective du continent. Son siège est à Addis-Abeba.
\end{definition}

\begin{savaistu}[L'OAPI à Yaoundé]
Peu de Camerounais le savent : le siège de l'\cle{Organisation Africaine de la Propriété Intellectuelle} (OAPI), institution spécialisée liée au dispositif de l'UA, se trouve à \emph{Yaoundé}. Créée par l'Accord de Libreville (1962, révisé à Bangui en 1977), elle protège les brevets, marques et droits d'auteur dans 17 États africains. Yaoundé abrite aussi des représentations continentales de recherche et de coopération. La capitale camerounaise est donc discrètement une \emph{ville diplomatique africaine}.
\end{savaistu}

\courssub{La CEMAC et la CEEAC : deux organisations, une même sous-région}

\textbf{Intuition.} Imaginez un quartier où les habitants créent deux associations : l'une gère la caisse commune et le marché (l'économie), l'autre s'occupe de la sécurité et des disputes entre voisins (la politique). En Afrique centrale, ces deux associations existent : la CEMAC et la CEEAC.

La \cle{Communauté Économique et Monétaire de l'Afrique Centrale} (CEMAC), créée par le traité de N'Djamena en \cle{1994} (succédant à l'UDEAC de 1964), regroupe six États : Cameroun, Congo, Gabon, Guinée équatoriale, RCA et Tchad. Son siège est à \cle{Bangui}. Elle assure :
\begin{itemize}
\item une \cle{union monétaire} : le franc CFA de la zone (XAF), géré par la BEAC (Banque des États de l'Afrique Centrale, siège à Yaoundé) ;
\item un \cle{marché commun} : libre circulation des biens, des services, des capitaux et des personnes en principe, tarif extérieur commun (TEC) ;
\item une convergence des politiques économiques (via la Commission de la CEMAC, installée à Libreville, et le Comité de pilotage de la convergence).
\end{itemize}

La \cle{Communauté Économique des États de l'Afrique Centrale} (CEEAC), créée en \cle{1983} à Libreville et dont le siège a été transféré à \cle{Malabo} (Guinée équatoriale) en 2015, regroupe onze États (les six de la CEMAC plus l'Angola, le Burundi, la RDC, São Tomé-et-Principe et le Rwanda). Sa vocation est plus large : intégration économique, mais surtout \cle{paix, sécurité et stabilité} dans la sous-région, via le COPAX.

\begin{attention}[Ne pas confondre CEMAC et CEEAC]
Piège classique à l'examen : la \cle{CEMAC} traite d'\emph{économie et de monnaie} (6 membres, siège à Bangui), la \cle{CEEAC} traite de \emph{politique, sécurité et intégration large} (11 membres, siège à Malabo). Leurs traités sont distincts mais leurs membres se recoupent largement. La \cle{double appartenance} coûte cher aux États (contributions doublées, sommets doublés) : leur fusion, décidée en principe lors du sommet de N'Djamena en 2013, reste un « serpent de mer » de la diplomatie centrafricaine.
\end{attention}

\begin{propriete}[Poids économique du Cameroun en Afrique centrale]
Le Cameroun est le \cle{premier contributeur au budget} de la CEMAC et de la CEEAC. Son PIB représente environ \cle{50 \%} du PIB total de la zone CEMAC, et sa population (environ 28 millions d'habitants en 2024) près de la moitié de celle de la sous-région CEMAC. Le port de Douala et le port en eau profonde de Kribi servent de débouchés maritimes vitaux au Tchad et à la RCA, pays enclavés. Cette position fait du Cameroun le \emph{moteur économique} et la \cle{plateforme logistique} de l'Afrique centrale.
\end{propriete}

\begin{popfigure}
\begin{center}
\begin{tikzpicture}[scale=1, font=\small]
% Fond
\fill[popcream] (-0.5,-0.6) rectangle (12.5,7.4);
% Cameroun (polygone simplifié)
\fill[popblue!35, draw=popblue, thick] (3.2,1.2) -- (4.6,0.4) -- (5.6,1.4) -- (5.9,2.8) -- (6.6,3.6) -- (6.2,4.6) -- (6.6,5.6) -- (5.8,6.6) -- (4.6,6.9) -- (3.6,6.0) -- (3.0,4.8) -- (3.4,3.4) -- (2.8,2.4) -- cycle;
\node[popdark, font=\bfseries] at (4.7,3.9) {CAMEROUN};
% Voisins
\fill[popmuted!25, draw=popmuted] (0.2,3.4) -- (2.8,2.4) -- (3.4,3.4) -- (3.0,4.8) -- (3.6,6.0) -- (4.6,6.9) -- (3.4,7.2) -- (1.2,7.0) -- (0.2,5.6) -- cycle;
\node[popmuted] at (1.7,5.2) {NIGERIA};
\fill[popmuted!25, draw=popmuted] (4.6,6.9) -- (5.8,6.6) -- (6.6,5.6) -- (8.6,5.8) -- (9.6,6.8) -- (8.0,7.3) -- (5.0,7.3) -- cycle;
\node[popmuted] at (7.4,6.6) {TCHAD};
\fill[popmuted!25, draw=popmuted] (6.2,4.6) -- (6.6,5.6) -- (8.6,5.8) -- (9.8,4.8) -- (9.4,3.4) -- (7.6,2.8) -- (6.6,3.6) -- cycle;
\node[popmuted] at (8.1,4.3) {RCA};
\fill[popmuted!25, draw=popmuted] (5.9,2.8) -- (6.6,3.6) -- (7.6,2.8) -- (8.4,1.6) -- (7.2,0.2) -- (5.6,1.4) -- cycle;
\node[popmuted] at (7.0,1.9) {CONGO};
\fill[popmuted!25, draw=popmuted] (3.2,1.2) -- (4.6,0.4) -- (5.6,1.4) -- (4.4,-0.4) -- (2.6,-0.2) -- cycle;
\node[popmuted] at (4.1,0.3) {\scriptsize GABON / G.\'E.};
% Villes et sièges
\fill[popink] (4.5,3.3) circle (0.09); \node[popink, anchor=west, font=\scriptsize] at (4.6,3.15) {Yaound\'e (BEAC, OAPI)};
\fill[popink] (3.3,2.1) circle (0.09); \node[popink, anchor=east, font=\scriptsize] at (3.2,2.0) {Douala};
\fill[poporange] (3.6,1.5) circle (0.09); \node[poporange, anchor=west, font=\scriptsize] at (3.7,1.5) {Kribi (port)};
\fill[popink] (8.1,4.05) circle (0.08); \node[popink, anchor=south, font=\scriptsize] at (8.3,3.7) {Bangui (CEMAC)};
\fill[popink] (2.7,0.1) circle (0.08); \node[popink, anchor=north, font=\scriptsize] at (2.7,-0.15) {Malabo (CEEAC)};
% Axes routiers structurants
\draw[popgreen, very thick, ->] (4.5,3.3) -- (8.1,4.05) node[midway, above, font=\scriptsize, popgreen]{axe Yaound\'e--Bangui};
\draw[popgreen, very thick, ->] (4.5,3.3) -- (6.8,6.5) node[midway, right, font=\scriptsize, popgreen]{axe vers N'Djam\'ena};
\draw[popgreen, very thick, ->] (4.5,3.3) -- (7.0,1.6) node[midway, below, font=\scriptsize, popgreen]{axe Yaound\'e--Brazzaville};
% Zones d'insécurité
\fill[poppink!45, draw=poppink, dashed] (2.9,4.6) ellipse (0.55 and 0.4);
\node[poppink, font=\scriptsize] at (2.9,4.05) {Bakassi};
\fill[poppink!45, draw=poppink, dashed] (5.4,6.3) ellipse (0.7 and 0.45);
\node[poppink, font=\scriptsize] at (5.4,5.75) {Extr\^eme-Nord};
\fill[poppink!45, draw=poppink, dashed] (6.4,4.1) ellipse (0.55 and 0.5);
\node[poppink, font=\scriptsize] at (6.4,3.45) {Est};
% Flux de réfugiés
\draw[poppurple, very thick, ->, dashed] (8.6,4.3) -- (6.9,4.15);
\draw[poppurple, very thick, ->, dashed] (2.6,5.6) -- (4.6,6.2);
% Nord
\draw[->, popdark, thick] (11.9,5.9) -- (11.9,6.9); \node[popdark] at (11.9,7.15) {N};
% Légende graphique
\node[anchor=west, font=\scriptsize] at (9.9,2.6) {\textcolor{popgreen}{\rule{0.5cm}{2pt}} axes routiers};
\node[anchor=west, font=\scriptsize] at (9.9,2.1) {\textcolor{poppurple}{\rule{0.5cm}{2pt}} flux de r\'efugi\'es};
\node[anchor=west, font=\scriptsize] at (9.9,1.6) {\textcolor{poppink}{\rule{0.4cm}{0.4cm}} zones d'ins\'ecurit\'e};
\end{tikzpicture}
\end{center}
\legende{Croquis schématique : le Cameroun, pivot de l'Afrique centrale. 1. Axes routiers structurants vers Bangui, N'Djamena et Brazzaville ; 2. flux de réfugiés centrafricains (Est) et nigérians (Extrême-Nord) ; 3. zones d'insécurité (Bakassi, Extrême-Nord, Est) ; 4. sièges d'institutions : Bangui (CEMAC), Malabo (CEEAC), Yaoundé (BEAC, OAPI). Échelle indicative, contours simplifiés.}
\end{popfigure}

\courssub{Dossier 4 : le COPAX, outil de paix de l'Afrique centrale}

\textbf{Intuition.} Quand une querelle éclate entre voisins, il vaut mieux qu'un comité du quartier intervienne vite, avant que le feu ne prenne. Le COPAX est ce « comité de quartier » de l'Afrique centrale : il prévient, alerte et, si nécessaire, déploie une force.

\begin{definition}[COPAX]
Le \cle{COPAX} (Conseil de Paix et de Sécurité de l'Afrique Centrale) est le mécanisme de \emph{prévention, de gestion et de règlement des conflits} en Afrique centrale, institué par le \cle{Protocole de 2000} (Protocole relatif à l'établissement du COPAX, adopté à N'Djamena dans le cadre de la CEEAC). Il repose sur deux instruments opérationnels principaux : le \cle{MARAC} (Mécanisme d'Alerte Rapide de l'Afrique Centrale), chargé de l'observation et de l'alerte précoce, et la \cle{FOMAC} (Force Multinationale de l'Afrique Centrale), chargée des opérations de maintien et de consolidation de la paix.
\end{definition}

Le COPAX s'articule en organes hiérarchisés : le \cle{Sommet des Chefs d'État} (organe suprême de décision), le \cle{Conseil des Ministres} (préparation et suivi des décisions), le \cle{Comité des Chefs d'État-Major} (planification militaire), le MARAC (renseignement et alerte) et la FOMAC (déploiement sur le terrain).

\begin{popfigure}
\begin{center}
\begin{tikzpicture}[font=\small,
  box/.style={draw=popblue, fill=popblueL, rounded corners=2pt, align=center, minimum height=0.9cm, text width=4.2cm},
  mil/.style={draw=poporange, fill=poporangeL, rounded corners=2pt, align=center, minimum height=0.9cm, text width=4.2cm},
  op/.style={draw=popgreen, fill=popgreenL, rounded corners=2pt, align=center, minimum height=0.9cm, text width=4.2cm}]
\node[box, fill=popblue!40, text width=6.5cm] (sommet) at (0,4.2) {\textbf{Sommet des Chefs d'\'Etat et de Gouvernement}\\ \scriptsize organe supr\^eme de d\'ecision};
\node[box, align=center] (ministres) at (-3.2,2.2){\textbf{Conseil des Ministres}\\ \scriptsize pr\'epare et suit les d\'ecisions};
\node[mil, align=center] (etatmajor) at (3.2,2.2){\textbf{Comit\'e des Chefs\\ d'\'Etat-Major}\\ \scriptsize planification militaire};
\node[op, align=center] (marac) at (-3.2,0){\textbf{MARAC}\\ \scriptsize alerte pr\'ecoce, observation};
\node[op, align=center] (fomac) at (3.2,0){\textbf{FOMAC}\\ \scriptsize force d'interposition et de paix};
\draw[->, popblue, thick] (sommet) -- (ministres);
\draw[->, popblue, thick] (sommet) -- (etatmajor);
\draw[->, popgreen, thick] (ministres) -- (marac);
\draw[->, poporange, thick] (etatmajor) -- (fomac);
\draw[<->, popmuted, dashed] (marac) -- (fomac) node[midway, below, font=\scriptsize, popmuted]{coordination};
\end{tikzpicture}
\end{center}
\legende{Architecture du COPAX. 1. Le Sommet des Chefs d'État décide ; 2. le Conseil des Ministres et le Comité des Chefs d'État-Major exécutent ; 3. le MARAC alerte et observe ; 4. la FOMAC intervient sur le terrain (RCA : FOMUC 2002, MICOPAX 2008, relais MINUSCA 2014 ; Tchad 2008).}
\end{popfigure}

Le COPAX est intervenu concrètement à plusieurs reprises : en \cle{RCA} (mission \cle{FOMUC} de la CEMAC dès 2002, puis \cle{MICOPAX} de la CEEAC/COPAX à partir de 2008, relais de la \cle{MINUSCA} onusienne en 2014), au \cle{Tchad} lors des tensions de 2008 (opération FOMAC). Le Cameroun y contribue par des contingents militaires et policiers, des financements et l'accueil de réunions de crise (sommet de Yaoundé 2013 sur la RCA).

\begin{methode}[Analyser un conflit régional : le triangle « Acteurs -- Ressources -- Institutions »]
Pour analyser rigoureusement un conflit en Afrique centrale (dissertation ou étude de cas), procédez en trois étapes :
\begin{enumerate}
\item \textbf{Identifier les acteurs} : États, groupes armés, populations, puissances extérieures, organisations régionales.
\item \textbf{Repérer les ressources en jeu} : terres, eau, minerais (diamants, or, pétrole), routes commerciales, position frontalière stratégique.
\item \textbf{Évaluer les institutions de médiation} : COPAX, CEEAC, UA, ONU, commissions mixtes bilatérales — leurs moyens, leurs succès, leurs limites.
\end{enumerate}
Concluez en jugeant l'efficacité de la médiation et en proposant des pistes de sortie de crise (sécurisation, partage des richesses, réforme institutionnelle).
\end{methode}

\begin{exoresolu}[Application de la méthode : la crise centrafricaine de 2013-2014]
Énoncé. En mars 2013, la coalition rebelle Séléka renverse le président centrafricain François Bozizé ; des milices anti-balaka ripostent, plongeant le pays dans une guerre civile. Analysez cette crise à l'aide du triangle « Acteurs -- Ressources -- Institutions » et précisez le rôle du Cameroun.
\tcblower
Solution détaillée.
\textbf{1. Acteurs.} D'un côté, la Séléka (rebelles majoritairement musulmans, hétéroclites) ; de l'autre, les milices anti-balaka (autodéfense à dominante chrétienne/animiste) ; l'État centrafricain effondré (armée inexistante, administration absente) ; les États voisins (Tchad, Cameroun, Congo, RDC) ; les organisations (CEEAC/COPAX via MICOPAX, UA via MISCA, ONU via MINUSCA, France via Opération Sangaris).
\textbf{2. Ressources.} Le contrôle du pouvoir à Bangui, les gisements de diamants et d'or (financement des groupes armés), les axes commerciaux vers Douala (voie d'approvisionnement vitale de la RCA enclavée : corridor Douala--Bangui), et les dynamiques identitaires instrumentalisées.
\textbf{3. Institutions de médiation.} La CEEAC déploie la MICOPAX (force du COPAX), relayée en 2013 par la MISCA (UA) puis en 2014 par la MINUSCA (ONU) ; le président congolais Denis Sassou Nguesso (médiateur CEEAC) et les sommets de N'Djamena organisent la transition politique (Catherine Samba-Panza).
\textbf{Rôle du Cameroun.} Il fournit des contingents à la MICOPAX puis à la MINUSCA, sécurise le corridor Douala--Bangui (approvisionnement en carburant, vivres, médicaments), accueille plus de 300 000 réfugiés dans l'Est et l'Adamaoua (camps de Gado, Borgop, etc.) et abrite des réunions de médiation.
\textbf{Conclusion.} La médiation régionale a évité l'effondrement total mais a montré ses limites (moyens insuffisants de la FOMAC/MICOPAX, mandat robuste manquant), d'où la nécessité du relais onusien (MINUSCA, Chapitre VII). Le triangle montre que la paix durable exige la combinaison des trois dimensions : désarmer les acteurs (DDR), partager équitablement les ressources, renforcer les institutions de l'État centrafricain.
\end{exoresolu}

\courssub{La gestion des crises frontalières}

Le Cameroun partage des frontières avec six États ; plusieurs ont donné lieu à des contentieux juridiques ou à des crises humanitaires et sécuritaires.

\textbf{Le différend de Bakassi (Cameroun--Nigeria).} Cette presqu'île riche en pétrole et en ressources halieutiques fut disputée entre les deux pays, avec des accrochages armés dans les années 1980-1990 (lac Tchad, Bakassi). Le Nigeria accepta de saisir la \cle{Cour Internationale de Justice} (CIJ), qui, par un arrêt du \cle{10 octobre 2002}, attribua la souveraineté de Bakassi au Cameroun. L'\cle{Accord de Greentree} (juin 2006), sous l'égide du Secrétaire général de l'ONU Kofi Annan, organisa le transfert pacifique de l'administration et le retrait des troupes nigérianes, supervisé par la \cle{Commission mixte Cameroun--Nigeria} (CMCN) présidée par un représentant du Secrétaire général. Le retrait nigérian de la péninsule fut effectif en \cle{août 2008} ; la Commission mixte a achevé ses travaux de démarcation terrestre et maritime en \cle{2013} : c'est un \emph{modèle africain de règlement pacifique} d'un conflit frontalier par la voie judiciaire.

\textbf{Le bassin du Lac Tchad.} Le Cameroun est membre de la \cle{Commission du Bassin du Lac Tchad} (CBLT, créée en 1964, siège à N'Djamena), qui gère la ressource en eau et le développement durable du bassin. Mais la région de l'Extrême-Nord subit depuis \cle{2014} les attaques de la secte \cle{Boko Haram} (et de sa faction ISWAP), qui provoquent déplacements internes (plus de 300 000 PDI) et afflux de réfugiés nigérians (camp de Minawao, géré par le HCR).

\textbf{La frontière avec la RCA.} À l'Est, l'insécurité chronique centrafricaine (groupes armés, coupeurs de route, \emph{zaraguinas}) fragilise la frontière poreuse et provoque l'arrivée massive et continue de réfugiés (vagues de 2005, 2013-2014, 2020-2021).

\begin{exemplebox}[Le Cameroun, premier pays d'asile de la sous-région]
Selon le HCR (données 2024), le Cameroun accueille environ \cle{480 000 réfugiés et demandeurs d'asile} : une majorité de \emph{Centrafricains} (environ 350 000) dans les régions de l'Est, de l'Adamaoua et du Nord, et des \emph{Nigérians} (environ 120 000) dans l'Extrême-Nord fuyant Boko Haram (camp de Minawao et sites spontanés). Le Cameroun est ainsi le \cle{premier pays d'asile} d'Afrique centrale et l'un des premiers du continent. Cet accueil, conforme au droit international (Convention de Genève de 1951 sur le statut des réfugiés, Convention de l'OUA de 1969), illustre la solidarité régionale mais pèse lourdement sur les services sociaux (écoles, hôpitaux, sécurité alimentaire, eau/assainissement) des zones d'accueil, parmi les plus pauvres du pays, et crée parfois des tensions locales.
\end{exemplebox}

\courssub{Les projets structurants de l'intégration sous-régionale}

L'intégration ne se décrète pas seulement dans les traités : elle se construit avec du bitume, de l'électricité et des ports.

\begin{itemize}
\item \textbf{Routes (corridors multimodaux)} : le corridor \cle{Douala--Bangui} (route N1/N3, ~1000 km) est le poumon commercial de la RCA (90 \% de ses échanges) ; l'axe \cle{Douala--N'Djamena} (via Garoua, Maroua, Kousseri, ~1800 km) alimente le Tchad ; le corridor \cle{Yaoundé--Brazzaville} (via Sangmélima, Ouesso) et l'axe vers \cle{Libreville} désenclave le nord du Congo et le Gabon. Ces axes sont au cœur du Programme PIDA (Programme pour le développement des infrastructures en Afrique) de l'UA.
\item \textbf{Énergie} : le projet d'\cle{interconnexion électrique Cameroun--Tchad} (ligne haute tension 225 kV), adossé aux barrages de \cle{Lom Pangar} (régulateur du Sanaga, mis en service 2017) et de \cle{Nachtigal} (420 MW, en construction, mise en service prévue 2024-2025), doit permettre au Cameroun d'exporter de l'électricité excédentaire vers le Tchad (et à terme vers la RCA). Le projet Chollet (sur la Dja, frontière Congo/Cameroun) concerne l'interconnexion Cameroun--Congo.
\item \textbf{Port de Kribi} : port en eau profonde (phase 1 inaugurée en 2018), il est conçu comme un \cle{hub sous-régional} capable d'accueillir les grands porte-conteneurs (Post-Panamax) et de servir le Tchad et la RCA via le corridor routier/ferroviaire projeté (chemin de fer Kribi--Mbalam--Ngaoundéré), soulageant Douala (port d'estuaire, tirant d'eau limité).
\end{itemize}

Ces infrastructures font du Cameroun la \emph{plateforme logistique} de l'Afrique centrale : elles concrétisent son rôle de pivot économique.

\courssub{Les dysfonctionnements de l'intégration régionale}

Malgré ces avancées, l'intégration centrafricaine souffre de faiblesses réelles qu'il faut savoir critiquer à l'examen.

\begin{enumerate}
\item \textbf{Libre circulation entravée.} En théorie, les citoyens de la CEMAC circulent librement avec la carte d'identité biométrique CEMAC ; en pratique, les \cle{contrôles routiers} multiples (police, gendarmerie, douane, eaux et forêts), les tracasseries administratives et le racket aux frontières freinent les échanges, découragent les commerçants et renchérissent le coût du transport.
\item \textbf{Double appartenance coûteuse et redondante.} Les États financent à la fois la CEMAC et la CEEAC, aux chevauchements nombreux (membres, objectifs, secrétariats) ; la fusion décidée en 2013 n'aboutit pas (disputes sur le siège, la répartition des postes, l'acquis communautaire), ce qui dilue les moyens et la lisibilité pour les citoyens.
\item \textbf{Inégalités de développement et faible commerce intra-zone.} Le Cameroun concentre environ la moitié du PIB de la zone CEMAC ; les échanges intra-communautaires restent très faibles (\textbf{moins de 5 \%} du commerce total des États membres), car les économies sont structurellement similaires (exportatrices de matières premières : pétrole, bois, cacao, coton, minerais) et s'exportent surtout vers l'Europe, la Chine et l'Inde, non vers le voisin.
\item \textbf{Insécurité persistante et dépendance extérieure.} Boko Haram/ISWAP dans le bassin du Lac Tchad, les crises centrafricaines récurrentes et la piraterie dans le golfe de Guinée montrent que la FOMAC manque de moyens propres (logistique, renseignement, financement) et dépend largement des financements extérieurs (UE via la Facilité pour la paix, ONU, partenaires bilatéraux) pour déployer ses opérations.
\end{enumerate}

\begin{aretenir}[L'essentiel de la section]
\begin{itemize}
\item Le Cameroun est membre fondateur de l'OUA (1963) et membre actif de l'UA (2002), où il a exercé des vice-présidences (2007, 2013, 2019).
\item La \cle{CEMAC} (6 États, siège à Bangui, BEAC à Yaoundé) gère l'économie et la monnaie (XAF) ; la \cle{CEEAC} (11 États, siège à Malabo) gère la politique, la sécurité (COPAX) et l'intégration large. Le Cameroun est le premier contributeur des deux organisations et pèse environ 50 \% du PIB de la zone CEMAC.
\item Le \cle{COPAX} (Protocole de 2000) prévient et règle les conflits grâce au MARAC (alerte) et à la FOMAC (force d'intervention : FOMUC/MICOPAX en RCA 2002/2008, Tchad 2008).
\item Le règlement pacifique de Bakassi (CIJ 2002, Greentree 2006, retrait 2008, fin travaux 2013) est un modèle ; le Cameroun accueille environ 480 000 réfugiés (HCR 2024), premier pays d'asile de la sous-région.
\item Les projets structurants (corridors Douala--Bangui, Douala--N'Djamena, Lom Pangar, Nachtigal, port de Kribi) font du Cameroun le hub de l'Afrique centrale, mais la libre circulation entravée, la double appartenance, les inégalités et l'insécurité freinent l'intégration réelle.
\end{itemize}
\end{aretenir}

\begin{exercice}[Pour s'entraîner]
\begin{enumerate}
\item Définis le COPAX et présente ses deux instruments opérationnels.
\item Construis un tableau comparatif CEMAC / CEEAC (date de création, siège, nombre de membres, domaine de compétence principal, monnaie).
\item Sujet de réflexion : « Le Cameroun, pivot de l'intégration en Afrique centrale : force ou fardeau ? » Rédige un développement argumenté en mobilisant au moins trois exemples précis étudiés dans cette section.
\end{enumerate}
\end{exercice}

\begin{corrige}[Exercices de la section]
\textbf{1.} Le COPAX (Conseil de Paix et de Sécurité de l'Afrique Centrale) est le mécanisme régional de prévention, de gestion et de règlement des conflits, institué par le Protocole de 2000 (CEEAC). Ses deux instruments opérationnels sont : le \textbf{MARAC} (Mécanisme d'Alerte Rapide de l'Afrique Centrale), chargé de l'observation, de l'analyse et de l'alerte précoce sur les crises potentielles ; et la \textbf{FOMAC} (Force Multinationale de l'Afrique Centrale), force de déploiement rapide chargée des opérations de maintien de la paix (ex. : RCA 2002 FOMUC, 2008 MICOPAX, Tchad 2008).

\textbf{2.}
\begin{center}
\begin{tabularx}{\linewidth}{|l|X|X|}\hline
\rowcolor{popblueL} \textbf{Critère} & \textbf{CEMAC} & \textbf{CEEAC} \\ \hline
Création & 1994 (Traité de N'Djamena) & 1983 (Traité de Libreville) \\ \hline
Siège institutionnel & Bangui (RCA) & Malabo (Guinée équatoriale) \\ \hline
Siège organes techniques & Yaoundé (BEAC), Libreville (Commission) & -- \\ \hline
Nombre de membres & 6 États & 11 États \\ \hline
Domaine principal & Union économique et monétaire (XAF, marché commun) & Paix, sécurité, stabilité, intégration globale \\ \hline
Monnaie & Franc CFA (XAF) géré par la BEAC & Aucune (diverses monnaies nationales) \\ \hline
\end{tabularx}
\end{center}

\textbf{3.} \textbf{Introduction :} Le Cameroun, par sa géographie (charnière Afrique de l'Ouest/Afrique centrale), son économie (1er PIB CEMAC) et son histoire, s'impose comme le \cle{pivot} de l'intégration centrafricaine. Cette position est-elle un atout (force) ou une charge (fardeau) ?

\textbf{I. Une force : le moteur de l'intégration}
\begin{itemize}
\item \textbf{Poids économique et financier :} 1er contributeur aux budgets CEMAC/CEEAC, ~50 \% du PIB CEMAC. Le port de Douala (1er port bois d'Afrique) et Kribi (hub futur) assurent la survie économique du Tchad et de la RCA (enclavement).
\item \textbf{Leadership diplomatique et sécuritaire :} Rôle clé dans le COPAX (troupes, financement, médiation Yaoundé 2013). Règlement exemplaire de Bakassi (CIJ, Greentree) : modèle de droit international.
\item \textbf{Infrastructures structurantes :} Corridors Douala--Bangui, Douala--N'Djamena ; barrages Lom Pangar/Nachtigal (export électricité) ; projet ferroviaire Kribi--Ngaoundéré. Le Cameroun \emph{construit} l'intégration physique.
\end{itemize}

\textbf{II. Un fardeau : les coûts du leadership}
\begin{itemize}
\item \textbf{Charge humanitaire et sécuritaire :} 480 000 réfugiés (HCR 2024) à charge sur des régions pauvres (Est, Extrême-Nord, Adamaoua). Importation de l'insécurité (Boko Haram, groupes armés RCA) : coût militaire, vies humaines, déplacés internes (Extrême-Nord).
\item \textbf{Coût institutionnel :} Double appartenance CEMAC/CEEAC (budgétisation lourde, sommets parallèles). Fusion bloquée depuis 2013.
\item \textbf{Asymétrie économique :} Faiblesse des échanges intra-zone (< 5 \%) : le Cameroun exporte peu vers ses partenaires CEMAC/CEEAC (produits similaires), ses clients sont hors zone. L'intégration ne profite pas encore pleinement au tissu productif local (PME).
\end{itemize}

\textbf{Conclusion nuancée :} Le statut de pivot est une \emph{force potentielle} (actif géostratégique, diplomatique, infrastructurel) qui devient \emph{fardeau réel} si l'intégration reste formelle. Transformer la contrainte en opportunité exige : 1) la finalisation de la fusion CEMAC/CEEAC pour rationaliser ; 2) la levée effective des barrières routières (libre circulation réelle) ; 3) la diversification des échanges intra-zone (transformation locale, ZLECAf) ; 4) un partage équitable des coûts de la sécurité (fonds de solidarité régional). Le Cameroun a intérêt à une Afrique centrale intégrée et stable : c'est sa profondeur stratégique.
\end{corrige}

\coursec{Partenariats historiques et institutionnels : UE, Francophonie, Commonwealth, OCI}

La section précédente a mis en lumière l'ancrage régional du Cameroun au sein de l'Union Africaine, de la CEMAC/CEEAC et du mécanisme de sécurité du COPAX. Cette dynamique sous-régionale ne saurait être isolée : elle s'inscrit dans une architecture plus large où le Cameroun active des leviers diplomatiques, économiques et culturels multiples. La singularité du Cameroun réside dans son \cle{triple héritage} (francophone, anglophone, arabo-islamique) qui lui confère une capacité rare à dialoguer simultanément dans des cercles institutionnels distincts. Cette section analyse comment le pays gère ces appartenances croisées --- Union Européenne, Francophonie, Commonwealth, Organisation de la Coopération Islamique --- pour diversifier ses partenariats, mobiliser des ressources et porter sa voix sur la scène internationale.

\courssub{Le Cameroun et l'Union Européenne : de Yaoundé à l'Accord de Samoa}

\textbf{Genèse et cadre juridique.}\par
La relation Cameroun--UE puise sa source dans les \cle{Conventions de Yaoundé} (1963, 1969) signées dès l'indépendance, puis dans les \cle{Conventions de Lomé} (I à IV, 1975--2000) qui ont structuré le groupe \cle{ACP} (Afrique, Caraïbes, Pacifique). Le tournant majeur est l'\cle{Accord de Cotonou} (2000, révisé en 2005 et 2010) qui a introduit le principe de \cle{partenariat politique}, la \cle{conditionnalité} (respect des droits de l'homme, démocratie, bonne gouvernance) et la programmation par \cle{Fonds Européen de Développement (FED)}. Depuis novembre 2023, l'\cle{Accord de Samoa} (nouvel accord de partenariat OEACP--UE pour 20 ans) succède à Cotonou, renforçant la dimension géopolitique (migrations, climat, sécurité) et remplaçant le FED par l'instrument \cle{NDICI--Global Europe} (instrument unique de financement de l'action extérieure de l'UE).

\begin{definition}[APD --- Aide Publique au Développement]
L'APD désigne les ressources fournies par les pays membres du \cle{CAD/OCDE} (Comité d'Aide au Développement) aux pays en développement, sous forme de \textbf{dons} ou de \textbf{prêts à caractère concessionnel} (taux d'intérêt bas, longue maturité, période de grâce). Pour être comptabilisée comme APD, une transaction doit présenter un \textbf{élément de don (grant element) supérieur à 25\,\%} (calculé à un taux d'actualisation de 10\,\%). L'APD exclut les prêts aux conditions du marché, les investissements directs étrangers (IDE) et l'aide militaire.
\end{definition}

\textbf{Le 11\textsuperscript{e} FED (2014--2020) : enveloppe et secteurs.}\par
L'enveloppe indicative nationale du Cameroun s'élevait à environ \textbf{287 millions d'euros} (dons), structurée autour de trois secteurs de concentration :
\begin{enumerate}
    \item \textbf{Gouvernance économique et financière} : réforme des finances publiques, transparence budgétaire, lutte contre la corruption.
    \item \textbf{Développement rural durable} : chaînes de valeur agricoles, sécurité alimentaire, résilience climatique dans le Nord.
    \item \textbf{Infrastructures d'appui à l'intégration régionale} : corridors routiers (par exemple l'axe Bamenda--Mamfe--Ekok/Mfum vers le Nigeria), énergie.
\end{enumerate}
Un \cle{programme d'appui budgétaire} (tranche fixe et tranches variables) conditionnait le décaissement à des indicateurs de performance (recettes fiscales, dépenses sociales, audit de la dette).

\begin{exemplebox}[Projet phare : l'appui à la gouvernance économique et financière]
Les programmes successifs d'appui à la gouvernance économique et financière, financés sur les 10\textsuperscript{e} et 11\textsuperscript{e} FED, ont soutenu la modernisation de la gestion des finances publiques : informatisation de la chaîne de la dépense, généralisation de la comptabilité d'engagement et renforcement des capacités de la Cour des Comptes et de la CONSUPE. Résultat observable : amélioration de la crédibilité budgétaire et de la transparence dans les évaluations internationales de la gestion des finances publiques (cadre PEFA).
\end{exemplebox}

\begin{methode}[Lire un Cadre de Partenariat Pays (CPP) UE--Cameroun]
Le CPP (par exemple 2021--2027 sous NDICI) est la boussole de la coopération. Voici comment l'analyser en 4 étapes :
\begin{enumerate}
    \item \textbf{Identifier les piliers stratégiques} (ex. : 1. Paix, sécurité, stabilité ; 2. Croissance verte et emploi ; 3. Gouvernance démocratique et droits humains).
    \item \textbf{Décortiquer la matrice de résultats} : pour chaque pilier, repérer les \emph{indicateurs de résultat} (ex. « taux d'électrification rurale »), leur \emph{valeur de base} (baseline), leur \emph{cible} (target 2027) et la \emph{source de vérification} (ex. rapport sectoriel, enquête ECAM de l'INS).
    \item \textbf{Analyser la matrice de conditionnalité} : distinguer les \emph{conditions préalables} (ex. adoption d'un texte d'application d'une loi de décentralisation) des \emph{conditions de décaissement} (tranches variables budgétaires). Évaluer leur faisabilité politique.
    \item \textbf{Vérifier l'alignement} : le CPP reprend-il les priorités de la \cle{SND30} (Stratégie Nationale de Développement 2020--2030) ? Y a-t-il duplication avec d'autres PTF (Partenaires Techniques et Financiers) ?
\end{enumerate}
\end{methode}

\begin{popfigure}
\begin{tikzpicture}[scale=0.95, every node/.style={font=\small}]
    \colorlet{ue}{popblue!70}
    \colorlet{fr}{poporange!70}
    \colorlet{cw}{popgreen!70}
    \colorlet{oci}{poppurple!70}
    \colorlet{ua}{poppink!70}
    \colorlet{cemac}{popgold!70}
    \begin{scope}[fill opacity=0.35]
        \fill[ue]   (0, 1.2) circle (2.2);
        \fill[fr]   (1.8, 0) circle (2.2);
        \fill[cw]   (-1.8, 0) circle (2.2);
        \fill[oci]  (0, -1.5) circle (2.0);
        \fill[ua]   (2.5, 2.0) circle (1.6);
        \fill[cemac](-2.5, 2.0) circle (1.6);
    \end{scope}
    \node[ue, font=\bfseries] at (0, 3.5) {UE / OEACP (Samoa)};
    \node[fr, font=\bfseries] at (4.0, 0) {Francophonie (OIF)};
    \node[cw, font=\bfseries] at (-4.0, 0) {Commonwealth};
    \node[oci, font=\bfseries] at (0, -3.6) {OCI};
    \node[ua, font=\bfseries] at (3.6, 3.4) {Union Africaine};
    \node[cemac, font=\bfseries] at (-3.6, 3.4) {CEMAC / CEEAC};
    \node[align=center, font=\tiny] at (0.9, 1.0) {Diplomatie\\ multilatérale};
    \node[align=center, font=\tiny] at (-0.9, 1.0) {Valeurs\\ démocratiques};
    \node[align=center, font=\tiny] at (1.2, -0.6) {Éducation /\\ Formation pro.};
    \node[align=center, font=\tiny] at (-1.2, -0.6) {Justice /\\ Bilinguisme};
    \node[align=center, font=\tiny] at (0, -1.6) {Financement\\ (BID + UE/BAD)};
    \node[align=center, font=\tiny] at (1.6, 2.2) {Paix /\\ Sécurité};
    \node[align=center, font=\tiny] at (-1.6, 2.2) {Intégration\\ régionale};
\end{tikzpicture}
\legende{Diagramme schématique des appartenances institutionnelles du Cameroun et des leviers diplomatiques aux intersections. La zone centrale (chevauchement UE--Francophonie--Commonwealth) concentre les leviers de gouvernance, d'éducation et de justice. L'OCI ajoute le levier financier (BID) et politique. L'UA et la CEMAC/CEEAC ancrent l'action régionale. Schéma simplifié, sans échelle.}
\end{popfigure}

\courssub{La Francophonie : espace linguistique, éducatif et politique}

\textbf{Positionnement institutionnel.}\par
Le Cameroun, membre de l'\cle{ACCT} (Agence de Coopération Culturelle et Technique, créée en 1970 à Niamey) devenue \cle{OIF} (Organisation Internationale de la Francophonie, dénomination adoptée en 2005), occupe une place singulière : pays officiellement bilingue (français--anglais), il est l'un des États où le français est langue officielle majoritairement parlée en Afrique centrale. L'OIF opère via quatre opérateurs directs : l'\cle{Université Senghor} (Alexandrie), \cle{TV5Monde}, l'\cle{Association Internationale des Maires Francophones (AIMF)} et l'\cle{AUF} (Agence Universitaire de la Francophonie), très présente dans les universités camerounaises (campus numériques, appui à la recherche).

\begin{savaistu}[Le Cameroun et les grands rendez-vous francophones]
Le Cameroun participe à tous les Sommets de la Francophonie et y défend la prévention des conflits, la diversité culturelle et la place de la jeunesse. Yaoundé s'était vu attribuer l'organisation des \textbf{Jeux de la Francophonie} prévus en 2021 ; reportés en raison de la pandémie de Covid-19 puis de retards d'infrastructures, ils ont finalement été réattribués à Kinshasa (RDC), qui les a organisés en 2023. Cet épisode illustre les exigences logistiques des grands événements internationaux et la nécessité d'une coordination interministérielle rigoureuse (jeunesse et sports, travaux publics, enseignement).
\end{savaistu}

\textbf{Coopération éducative et formation professionnelle.}\par
L'OIF appuie le \cle{bilinguisme} et l'enseignement des langues, la \cle{formation professionnelle} (appuis aux réformes de l'enseignement technique et à l'insertion des jeunes) et la \cle{recherche} (réseaux universitaires, bourses de mobilité). Pour la filière technique, ces appuis se traduisent par des équipements d'ateliers, la formation des enseignants et l'alignement des diplômes sur les besoins des entreprises.

\courssub{Le Commonwealth : adhésion tardive, apports concrets et tensions internes}

\textbf{Contexte et adhésion (1995).}\par
L'entrée du Cameroun dans le \cle{Commonwealth} (Sommet d'Auckland, 1995) valorise l'héritage anglophone (ex-Southern Cameroons sous administration britannique) et ouvre l'accès aux réseaux \cle{anglophones} (justice, éducation, administration). Le Cameroun est, avec le Mozambique (1995) et le Rwanda (2009), l'un des rares membres \textbf{sans avoir été une colonie britannique au sens strict} : les régions du Nord-Ouest et du Sud-Ouest étaient des territoires sous mandat de la SDN puis sous tutelle de l'ONU, administrés par la Grande-Bretagne.

\textbf{Apports mesurables.}
\begin{itemize}
    \item \textbf{Bourses :} le \cle{Commonwealth Scholarship and Fellowship Plan (CSFP)} offre chaque année des bourses de masters et doctorats (Royaume-Uni, Canada, Afrique du Sud, Malaisie, etc.) aux étudiants et cadres camerounais.
    \item \textbf{Justice et administration :} appui technique à l'harmonisation du \cle{droit coutumier anglo-saxon (Common Law)} dans les régions du Nord-Ouest et du Sud-Ouest (formation des magistrats, traduction des textes juridiques).
    \item \textbf{Éducation :} le \cle{Commonwealth of Learning (COL)} soutient l'éducation ouverte et à distance (ressources éducatives libres), utile pour les zones en crise.
    \item \textbf{Diplomatie :} participation aux \cle{CHOGM} (Sommets des Chefs de Gouvernement du Commonwealth) --- plateforme de plaidoyer sur le climat (vulnérabilité du bassin du Congo) et la santé.
\end{itemize}

\begin{attention}[Piège : confusion entre complémentarité et substitution]
\textbf{Erreur fréquente :} « L'adhésion au Commonwealth remplace la Francophonie. »
\textbf{Réalité :} le Cameroun est membre \emph{à part entière} des \textbf{deux organisations} simultanément depuis 1995. C'est un \textbf{atout de diversification} (double accès aux bourses, aux forums, aux financements) et non un choix exclusif. La crise anglophone (depuis 2016) teste cette double appartenance : le Commonwealth (via le \cle{CMAG} --- Commonwealth Ministerial Action Group) suit la situation des droits humains, tandis que l'OIF propose son expertise en médiation. La diplomatie camerounaise doit gérer ces deux regards croisés sans les opposer.
\end{attention}

\courssub{L'Organisation de la Coopération Islamique (OCI) : le levier sahelo-saharien}

\textbf{Adhésion et identité (1974).}\par
L'adhésion du Cameroun à la \cle{Conférence Islamique} (devenue OCI en 2011) ancre le pays dans sa dimension \cle{nord/sahélienne} (régions de l'Adamaoua, du Nord et de l'Extrême-Nord, à forte population musulmane) et lui donne une tribune au sein du \cle{monde musulman} (57 États membres). Le Cameroun y défend la \cle{cause palestinienne} (alignement sur le consensus OCI/ONU) et la lutte contre l'islamophobie, tout en affirmant sa laïcité constitutionnelle.

\textbf{Levier financier : la Banque Islamique de Développement (BID / IsDB).}\par
C'est le bras opérationnel majeur de l'OCI. La BID finance au Cameroun via des \cle{prêts concessionnels} (modes \emph{Istisna'a}, \emph{Ijara}/leasing, \emph{Murabaha}) et des \cle{dons} :
\begin{itemize}
    \item \textbf{Infrastructures routières :} participation au financement d'axes structurants du Grand Nord (corridors vers le Tchad et le Nigeria).
    \item \textbf{Hydraulique et énergie :} cofinancement du barrage-réservoir de \cle{Lom Pangar} (région de l'Est) aux côtés de la Banque Mondiale, de la BAD, de l'AFD et de l'État camerounais ; projets d'adduction d'eau potable.
    \item \textbf{Secteur social :} appui aux infrastructures sanitaires et à la formation technique dans les régions septentrionales.
\end{itemize}
La \cle{finance islamique} (émissions de \emph{Sukuk} souverains) est à l'étude pour diversifier les sources de financement de la dette publique.

\courssub{Articulation des trois cercles : la « diplomatie du triple héritage »}

Le Cameroun active une \cle{diplomatie du triple héritage} (français, anglais, arabe/islamique) qui structure sa politique étrangère :
\begin{center}
\begin{tabularx}{\linewidth}{|l|X|X|X|}\hline
\rowcolor{popblueL}\textbf{Cercle} & \textbf{Langue de travail} & \textbf{Capital diplomatique} & \textbf{Ressources mobilisables} \\ \hline
Francophonie / UE & Français & Réseaux institutionnels (OIF, AUF), APD en dons (FED/NDICI), expertise administrative & Gouvernance, éducation, monde rural, infrastructures \\ \hline
Commonwealth & Anglais & Réseaux juridiques (Common Law), bourses CSFP, diaspora anglophone, CHOGM & Justice, enseignement supérieur, plaidoyer climatique \\ \hline
OCI / monde arabe & Arabe / Français & Solidarité politique (Palestine, Sahel), finance islamique (BID, Sukuk), marchés halal & Infrastructures lourdes, énergie, agriculture, finance \\ \hline
\end{tabularx}
\end{center}

\begin{aretenir}[Complexité institutionnelle]
Cette triple appartenance génère des \textbf{coûts de transaction} élevés : rapports distincts, missions d'évaluation parallèles, conditionnalités parfois contradictoires (ex. conditionnalité politique de l'UE vs principe de non-ingérence mis en avant à l'OCI). L'enjeu pour l'administration camerounaise (MINREX, MINEPAT) est l'\textbf{harmonisation des cadres de résultats} (alignement sur la SND30) et la \textbf{mutualisation des évaluations} (revues conjointes des bailleurs).
\end{aretenir}

\courssub{Dossier : évaluation de l'Aide Publique au Développement (APD) reçue via ces canaux}

\textbf{Structure de l'APD nette (ordres de grandeur OCDE/CAD, 2022--2023).}\par
L'APD reçue par le Cameroun se répartit entre \textbf{dons} (majoritaires pour le social, la gouvernance, l'humanitaire) et \textbf{prêts concessionnels} (majoritaires pour les infrastructures lourdes). Avec la graduation du Cameroun dans la catégorie des pays à revenu intermédiaire, la part des prêts concessionnels tend à augmenter par rapport aux dons.

\begin{popfigure}
\begin{center}
\begin{tikzpicture}
\begin{axis}[
    ybar,
    width=0.95\linewidth,
    height=7cm,
    bar width=10pt,
    symbolic x coords={UE,BAD,BM,ONU,BID,France,Chine,USA/Japon},
    xtick=data,
    x tick label style={rotate=30, anchor=east, font=\tiny},
    ymin=0, ymax=30,
    ylabel={Part (\%)},
    tick label style={font=\tiny},
    title={Répartition estimée de l'APD nette reçue par bailleur (2022--2023)},
    title style={font=\small\bfseries}
]
\addplot[fill=popblue!50, draw=popblue] coordinates {(UE,18) (BAD,14) (BM,12) (ONU,6) (BID,5) (France,15) (Chine,10) (USA/Japon,8)};
\end{axis}
\end{tikzpicture}
\end{center}
\legende{Diagramme en barres de la structure estimée de l'APD nette par bailleur principal (parts en \,\% d'un total de l'ordre de 600 à 700 millions de dollars par an). Les \textbf{multilatéraux} (UE, BAD, BM, ONU, BID) représentent environ 55\,\% : dons (UE, ONU) et prêts concessionnels (BAD, BM/IDA, BID). Les \textbf{bilatéraux} (France, Chine, USA, Japon) représentent environ 45\,\% : France (dons et prêts), Chine (prêts concessionnels et crédits à l'exportation), USA (dons santé et humanitaire), Japon (dons et prêts infrastructures). Source : OCDE/CAD et rapports annuels des bailleurs ; valeurs arrondies et indicatives.}
\end{popfigure}

\textbf{Conditionnalités et efficacité.}
\begin{itemize}
    \item \textbf{Conditionnalité politique (UE, Commonwealth) :} liée aux droits de l'homme, aux élections, à la gouvernance. Elle peut entraîner la suspension d'un appui budgétaire (ex. reports de tranches variables du 11\textsuperscript{e} FED en cas de retard des réformes des finances publiques).
    \item \textbf{Conditionnalité sectorielle (BM, BAD, BID) :} réformes structurelles (subventions carburant, gouvernance des entreprises publiques, climat des affaires).
    \item \textbf{Alignement et harmonisation (Déclaration de Paris 2005, Accra 2008, Busan 2011) :} progrès via le \cle{Cadre de Partenariat Pays (CPP)} de l'UE et le \cle{Cadre de Partenariat Pays (CPF)} de la Banque Mondiale. Persistent : lenteur des procédures de décaissement, fragmentation des projets (multiplication de petits projets), faiblesse de l'appropriation nationale (unités d'exécution de projets parallèles aux administrations).
\end{itemize}

\begin{exoresolu}[Analyse de la concessionnalité d'un prêt BID vs prêt commercial]
\textbf{Énoncé :} le Cameroun négocie un financement de 50 milliards de FCFA pour un barrage.
\begin{itemize}
    \item Option A (BID) : prêt de type \emph{Istisna'a}, taux de service 2\,\%/an, maturité 25 ans, différé d'amortissement 5 ans.
    \item Option B (banque commerciale) : prêt au taux de 6,5\,\%/an, maturité 12 ans, différé 2 ans.
\end{itemize}
Estimer l'élément de don (Grant Element) de chaque option (taux d'actualisation OCDE = 10\,\%) et déterminer laquelle est comptabilisable en APD.
\tcblower
\textbf{Solution détaillée :}
L'élément de don mesure l'écart entre la valeur nominale du prêt et la valeur actualisée du service de la dette :
\[ GE = 1 - \frac{VAN(\text{service de la dette})}{\text{Montant nominal}}. \]
Pour un prêt remboursé par annuités constantes après un différé $G$, on utilise l'approximation :
\[ VAN \approx \frac{r}{q}\left[1-(1+q)^{-(N-G)}\right](1+q)^{-G}\times M, \]
où $r$ est le taux nominal, $q=0{,}10$ le taux d'actualisation, $N$ la maturité, $G$ le différé, $M=50$ milliards.

\textbf{Option A (BID) :} $r=0{,}02$, $N=25$, $G=5$.
\begin{align*}
VAN_A &\approx \frac{0{,}02}{0{,}10}\left[1-1{,}10^{-20}\right]\times 1{,}10^{-5}\times 50 \\
&\approx 0{,}2 \times 0{,}851 \times 0{,}621 \times 50 \approx 5{,}3 \text{ milliards.}
\end{align*}
\[ GE_A = 1-\frac{5{,}3}{50} \approx 0{,}894 \approx 89\,\% > 25\,\%. \]
\textbf{Conclusion :} le prêt BID est comptabilisable en APD (concessionnalité très élevée, proche du don).

\textbf{Option B (commerciale) :} $r=0{,}065$, $N=12$, $G=2$.
\begin{align*}
VAN_B &\approx \frac{0{,}065}{0{,}10}\left[1-1{,}10^{-10}\right]\times 1{,}10^{-2}\times 50 \\
&\approx 0{,}65 \times 0{,}614 \times 0{,}826 \times 50 \approx 16{,}5 \text{ milliards.}
\end{align*}
\[ GE_B = 1-\frac{16{,}5}{50} \approx 0{,}67 = 67\,\% > 25\,\%. \]
\textbf{Conclusion :} ce prêt serait aussi comptabilisable en APD, ce qui montre la limite de l'exemple : à 6,5\,\% sur 12 ans, le prêt reste fortement concessionnel par rapport au taux d'actualisation de 10\,\%. Si le taux commercial réel était de 10\,\% (conditions de marché), alors $GE \approx 0\,\%$ et le prêt ne serait \textbf{pas} de l'APD. La BID offre donc une concessionnalité \emph{bien supérieure}, libérant de l'espace budgétaire pour les dépenses sociales.
\end{exoresolu}

\begin{exercice}[Analyse comparative de deux cadres de partenariat]
Consultez les extraits suivants de deux cadres de partenariat récents :
\begin{enumerate}
    \item \textbf{CPP UE--Cameroun 2021--2027 (pilier Gouvernance) :} indicateur : « taux d'exécution du budget d'investissement public supérieur à 85\,\% » ; condition de décaissement d'une tranche variable : « adoption des textes d'application de la loi sur la décentralisation (transferts de compétences et de ressources) ».
    \item \textbf{CPF Banque Mondiale 2022--2026 (objectif climat des affaires) :} indicateur : « nombre de jours pour créer une entreprise inférieur à 5 jours » ; jalon de décaissement (DLI) : « déploiement du guichet unique dématérialisé dans les chefs-lieux de région ».
\end{enumerate}
\textbf{Questions :}
\begin{enumerate}
    \item Identifiez le type de conditionnalité (politique ou sectorielle/technique) pour chaque jalon.
    \item Expliquez pourquoi le non-respect du jalon UE (décentralisation) a un impact diplomatique plus large que le retard du guichet unique de la BM.
    \item Proposez une action concrète du MINEPAT pour mutualiser le suivi de ces deux indicateurs.
\end{enumerate}
\end{exercice}

\begin{corrige}[Exercice n]
\begin{enumerate}
    \item \textbf{UE :} conditionnalité \textbf{politique et institutionnelle} (réforme structurelle de l'État, transfert de pouvoirs vers les collectivités territoriales). \textbf{BM :} conditionnalité \textbf{sectorielle et technique} (amélioration du climat des affaires, simplification administrative).
    \item Le jalon sur la décentralisation engage la \textbf{crédibilité démocratique} du pays vis-à-vis de l'UE : les accords de partenariat ACP--UE (Cotonou puis Samoa) font des droits de l'homme, de la démocratie et de l'État de droit des « éléments essentiels ». Leur non-respect peut ouvrir des \textbf{consultations politiques} (procédure de l'article 96 de Cotonou, reprise dans l'Accord de Samoa) voire une suspension de la coopération. Le guichet unique est une réforme technique : son retard pénalise le décaissement du projet concerné mais ne remet pas en cause la relation politique globale.
    \item \textbf{Action MINEPAT :} créer un \textbf{comité de pilotage unique « Réformes structurelles »} (coprésidé par le MINEPAT, le MINATD et le MINCOMMERCE) doté d'une \textbf{matrice de suivi harmonisée} reliant : (1) le taux d'exécution du budget d'investissement public (indicateur UE) ; (2) les dépenses d'investissement transférées aux régions (décentralisation) ; (3) les délais de création d'entreprise au niveau régional (guichet unique). Un tableau de bord trimestriel unique remplace les rapports séparés adressés à chaque bailleur.
\end{enumerate}
\end{corrige}

\coursec{Grands partenaires bilatéraux : France, Royaume-Uni, Chine, Japon, USA, Brésil}

Après avoir étudié les cadres \emph{institutionnels} et multilatéraux (Union européenne, Francophonie, Commonwealth, Organisation de la Coopération Islamique) dans lesquels le Cameroun s'inscrit, il faut maintenant descendre au niveau le plus concret de la diplomatie : les relations \cle{bilatérales}, c'est-à-dire les liens directs entre le Cameroun et un autre État. C'est à ce niveau que se signent les prêts, que se construisent les barrages, que se forment les soldats et que s'installent les entreprises. Comprendre ces partenariats, c'est comprendre qui finance le développement du Cameroun, qui achète ses matières premières, et à quelles conditions.

\begin{definition}[Coopération Nord-Sud et coopération Sud-Sud]
On distingue deux grands types de coopération bilatérale :
\begin{itemize}
\item la \cle{coopération Nord-Sud} : elle lie un pays développé (le « Nord ») à un pays en développement (le « Sud »). Exemples : Cameroun--France, Cameroun--USA, Cameroun--Japon. Elle s'accompagne souvent d'aide publique au développement (APD) et de \cle{conditionnalités} (exigences de bonne gouvernance, réformes économiques) ;
\item la \cle{coopération Sud-Sud} : elle lie deux pays en développement ou émergents, sur une base présentée comme un partenariat « gagnant-gagnant », sans leçons politiques. Exemples : Cameroun--Chine, Cameroun--Brésil.
\end{itemize}
\end{definition}

\courssub{La France : le premier partenaire historique}

La France est le partenaire le plus ancien et le plus complet du Cameroun : histoire (administration sous mandat puis tutelle), langue, armée, entreprises, communauté expatriée. Cette relation est à la fois très dense et très débattue.

\textbf{Les piliers de la relation franco-camerounaise}\par
\begin{itemize}
\item \textbf{Défense} : des accords de défense lient les deux pays depuis 1974 (renouvelés et actualisés, notamment en 2019). Ils prévoient coopération militaire, formation des officiers et assistance technique.
\item \textbf{Aide publique au développement} : l'\cle{Agence française de développement} (AFD) finance routes, eau potable, éducation et appui budgétaire. La France est l'un des premiers bailleurs de fonds du Cameroun.
\item \textbf{Investissements} : de grandes entreprises françaises sont présentes dans des secteurs stratégiques : Bolloré (port de Douala, chemin de fer Camrail), TotalEnergies (hydrocarbures, distribution), Orange (télécommunications), Vinci (construction).
\item \textbf{Présence humaine} : environ \num{20000} ressortissants français vivent au Cameroun, et une très importante diaspora camerounaise vit en France.
\end{itemize}

\textbf{Les contentieux}\par
La relation n'est pas sans tensions. On parle de \cle{Françafrique} pour désigner les réseaux d'influence politique et économique hérités de la colonisation. Les critiques portent sur : le poids de la dette, l'\emph{ingérence perçue} dans la vie politique interne, et la position dominante des entreprises françaises dans les secteurs rentiers (port, rail, énergie). Le débat sur la « souveraineté économique » du Cameroun passe largement par la relation avec Paris.

\begin{attention}[Piège de rédaction au Bac]
Ne réduis jamais la relation Cameroun--France à un jugement moral (« la France exploite » ou « la France aide »). Une bonne copie présente \emph{les deux faces} : apports réels (financements, formation, sécurité) \emph{et} limites réelles (dépendance, ingérence perçue, dette). C'est cette nuance argumentée qui est évaluée.
\end{attention}

\courssub{Le Royaume-Uni (Grande-Bretagne) : le partenaire discret mais efficace}

Héritier de l'administration des régions anglophones (Nord-Ouest et Sud-Ouest), le Royaume-Uni entretient une relation moins visible que la France, mais réelle et ciblée.

\begin{itemize}
\item \textbf{Investissements} : Shell (lubrifiants, distribution, anciennement amont pétrolier), G4S (sécurité privée, désormais Allied Universal), British American Tobacco (industrie, coté à la bourse de Douala).
\item \textbf{Formation des élites} : les bourses \cle{Chevening} permettent chaque année à de jeunes cadres camerounais de faire des masters dans les universités britanniques.
\item \textbf{Appui sécuritaire} : le Royaume-Uni participe à la formation du \cle{BIR} (Bataillon d'intervention rapide), unité d'élite engagée contre Boko Haram dans l'Extrême-Nord.
\end{itemize}

C'est une relation de « niche » : peu de grands discours, mais des leviers précis (langue anglaise, formation, sécurité, Commonwealth étudié à la section précédente).

\courssub{La Chine : premier partenaire commercial et financier d'infrastructures}

Depuis les années 2000, la Chine est devenue le \cle{premier partenaire commercial} du Cameroun et son premier financeur d'infrastructures. C'est le cas d'école de la coopération Sud-Sud.

\textbf{Les réalisations emblématiques}\par
\begin{itemize}
\item \textbf{Barrages hydroélectriques} : Memve'ele (Sud) et Lom Pangar (Est), qui augmentent fortement la production d'électricité ;
\item \textbf{Routes} : l'autoroute Kribi--Lolabé (liaison avec le port en eau profonde de Kribi) et la route Yaoundé--Nsimalen (accès aéroport) ;
\item \textbf{Équipements publics} : le stade de football d'Olembe (Yaoundé), l'Hôpital de gynéco-obstétrique et de pédiatrie de Yaoundé.
\end{itemize}

\textbf{Le modèle chinois}\par
La Chine ne fait pas de « dons » au sens classique : elle accorde des \cle{prêts concessionnels}, souvent via sa banque d'État \cle{Exim Bank}, remboursables sur longue période, parfois garantis par les ressources naturelles du pays (coton, bois, minerais) ou par des comptes séquestres (\emph{escrow}) sur les recettes futures. Les travaux sont généralement exécutés par des entreprises chinoises.

\begin{methode}[Analyser un accord de prêt chinois (Exim Bank)]
Face à un document présentant un prêt chinois, voici la démarche d'analyse en quatre étapes :
\begin{enumerate}
\item \textbf{Identifier le coût} : le taux d'intérêt (souvent 2 à 3 \%, donc « concessionnel », c'est-à-dire inférieur au taux du marché) ;
\item \textbf{Identifier la durée} : la maturité (durée totale de remboursement, souvent 20 ans) et le \cle{délai de grâce} (5 à 7 ans pendant lesquels on ne rembourse pas le capital, le temps que l'ouvrage produise des revenus) ;
\item \textbf{Repérer les garanties} : clause de garantie sur les ressources naturelles ou compte séquestre (\emph{escrow}) où transitent les recettes (ex. recettes du port ou de l'électricité) ;
\item \textbf{Conclure} : le prêt est-il soutenable ? L'ouvrage financé génère-t-il assez de revenus pour rembourser ? Qui réalise les travaux (main-d'œuvre locale ou importée) ?
\end{enumerate}
\end{methode}

\begin{exemplebox}[Le poids économique de la Chine au Cameroun]
Le commerce Cameroun--Chine atteint environ \num{2500} milliards de FCFA par an, avec un fort déséquilibre : les \emph{importations} camerounaises (machines, textiles, produits manufacturés) dépassent largement les \emph{exportations} (pétrole brut, bois, coton). Le stock d'investissements chinois au Cameroun dépasse \num{3000} milliards de FCFA.
\end{exemplebox}

\begin{attention}[La « dette chinoise » : un chiffre souvent mal compris]
Dans le débat public, on entend que « la Chine endette le Cameroun ». Il faut être précis : la dette envers la Chine représente environ 15 à 18 \% du \emph{PIB total}, mais environ 40 à 50 \% de la \emph{dette extérieure publique}. Les deux affirmations sont vraies mais ne mesurent pas la même chose. À l'examen, cite toujours la \emph{référence} (PIB ou dette extérieure) pour éviter l'exagération.
\end{attention}

\courssub{Le Japon : le partenaire de la « qualité »}

Le Japon mise sur la qualité technique plutôt que sur le volume. Son agence de coopération, la \cle{JICA}, finance et accompagne des projets ciblés, souvent en associant expertise japonaise et financement camerounais ou international : c'est l'approche dite « main dans la main » (technique + financement).

\begin{itemize}
\item \textbf{Infrastructures} : le deuxième pont sur le Wouri à Douala (pont de Bonabéri), ouvrage majeur pour désengorger la capitale économique ;
\item \textbf{Agriculture} : appui à la riziculture, notamment à la SEMRY (Yagoua, Extrême-Nord) ;
\item \textbf{Formation} : bourses et stages au Japon pour les cadres africains ;
\item \textbf{Santé} : équipements hospitaliers et appui aux programmes de santé publique.
\end{itemize}

\begin{savaistu}[Le programme ABE Initiative]
L'\cle{ABE Initiative} (\emph{African Business Education Initiative for Youth}), lancée par le Japon en 2014, a permis à environ 150 cadres camerounais de suivre un master et un stage en entreprise au Japon. Objectif : former une génération de gestionnaires capables de devenir des « ponts » économiques entre le Japon et l'Afrique.
\end{savaistu}

\courssub{Les États-Unis : sécurité, santé et gouvernance}

La coopération américaine se concentre sur trois domaines :
\begin{itemize}
\item \textbf{Sécurité} : formation du BIR, coopération en matière de renseignement et d'équipement (dont drones de surveillance) dans la lutte contre Boko Haram dans l'Extrême-Nord ;
\item \textbf{Santé} : le programme \cle{PEPFAR} (lutte contre le VIH/sida) et la \cle{PMI} (\emph{President's Malaria Initiative}, lutte contre le paludisme) financent traitements, moustiquaires et dépistage ;
\item \textbf{Gouvernance et économie} : l'agence \cle{USAID} appuie la société civile et la décentralisation ; des entreprises américaines investissent dans l'énergie (gaz de Kribi via Golar LNG/Perenco) et le pétrole.
\end{itemize}

Les États-Unis offrent aussi un accès privilégié à leur marché grâce à l'\cle{AGOA} (\emph{African Growth and Opportunity Act}), loi qui permet à certains produits africains d'entrer aux USA sans droits de douane (\emph{duty-free}).

\begin{attention}[Pourquoi l'AGOA profite peu au Cameroun]
L'AGOA est une opportunité théorique : en pratique, les exportations camerounaises vers les USA, hors pétrole, sont négligeables. La cause n'est pas le refus américain mais la \emph{faible capacité productive locale} : peu d'entreprises camerounaises produisent en quantité et en qualité suffisantes pour le marché américain. Leçon : un avantage commercial ne vaut que si l'appareil productif peut l'exploiter.
\end{attention}

\courssub{Le Brésil : la coopération Sud-Sud agricole}

Le Brésil, grande puissance agricole tropicale, partage avec le Cameroun un climat et des défis proches : c'est le partenaire naturel de la coopération Sud-Sud. Son agence de recherche agricole \cle{Embrapa}, mondialement reconnue, coopère avec le Cameroun sur l'agriculture tropicale : amélioration des rendements du coton et du manioc, techniques de culture, bioéthanol (carburant produit à partir de la canne à sucre ou du manioc). Des accords techniques portent aussi sur la santé et la formation. L'intérêt de ce partenariat : des technologies \emph{adaptées aux tropiques}, moins coûteuses que les technologies occidentales.

\begin{popfigure}
\begin{center}
\begin{tikzpicture}[font=\small]
% Cameroun au centre
\node[draw=popdark,fill=popgoldL,rounded corners,very thick,minimum width=3.2cm,minimum height=1.4cm,align=center] (cam) at (0,0) {\textbf{CAMEROUN}\\ \footnotesize Yaoundé};
% Partenaires
\node[draw=popblue,fill=popblueL,rounded corners,align=center] (fr) at (-5.5,3) {\textbf{France}\\ \footnotesize APD, défense,\\ \footnotesize pétrole, port, rail};
\node[draw=popblue,fill=popblueL,rounded corners,align=center] (uk) at (-6,0) {\textbf{Royaume-Uni}\\ \footnotesize Bourses Chevening,\\ \footnotesize formation BIR};
\node[draw=poppink,fill=poppinkL,rounded corners,align=center] (cn) at (0,4) {\textbf{Chine}\\ \footnotesize Barrages Memve'ele,\\ \footnotesize Lom Pangar, routes,\\ \footnotesize stade Olembe};
\node[draw=poporange,fill=poporangeL,rounded corners,align=center] (jp) at (5.8,3) {\textbf{Japon}\\ \footnotesize Pont sur le Wouri,\\ \footnotesize riz SEMRY, ABE};
\node[draw=popgreen,fill=popgreenL,rounded corners,align=center] (us) at (6,0) {\textbf{USA}\\ \footnotesize Santé PEPFAR/PMI,\\ \footnotesize sécurité, AGOA};
\node[draw=popgreen,fill=popgreenL,rounded corners,align=center] (br) at (0,-3.4) {\textbf{Brésil}\\ \footnotesize Embrapa : coton,\\ \footnotesize manioc, bioéthanol};
% Flux : épaisseur proportionnelle
\draw[-{Stealth},line width=3pt,popblue] (fr) -- (cam);
\draw[-{Stealth},line width=1.2pt,popblue] (uk) -- (cam);
\draw[-{Stealth},line width=4.5pt,poppink] (cn) -- (cam);
\draw[-{Stealth},line width=1.8pt,poporange] (jp) -- (cam);
\draw[-{Stealth},line width=2.2pt,popgreen] (us) -- (cam);
\draw[-{Stealth},line width=1pt,popgreen] (br) -- (cam);
\end{tikzpicture}
\end{center}
\legende{Carte de flux schématique : les partenaires stratégiques du Cameroun. L'épaisseur des flèches est proportionnelle à l'intensité globale des flux (APD, IDE, commerce) : la Chine et la France dominent, le Brésil reste modeste mais ciblé. Les mentions indiquent les projets ou secteurs phares de chaque partenaire.}
\end{popfigure}

\begin{popfigure}
\begin{center}
\begin{tikzpicture}[scale=1.05,font=\footnotesize]
% Grille radar 5 axes, 4 niveaux
\foreach \r in {1,2,3,4}{
  \draw[popline] (90:\r) -- (162:\r) -- (234:\r) -- (306:\r) -- (18:\r) -- cycle;
}
\foreach \a/\t in {90/{Volume APD},162/{Volume commerce},234/{Transfert technologique},306/{Alignement politique},18/{Conditionnalité}}{
  \draw[popmuted] (0,0) -- (\a:4.3);
  \node at (\a:4.9) {\t};
}
% France
\draw[very thick,popblue,fill=popblue!20] (90:3.5) -- (162:2.5) -- (234:2) -- (306:3.5) -- (18:3.5) -- cycle;
% Chine
\draw[very thick,poppink,fill=poppink!15] (90:2) -- (162:4) -- (234:3.5) -- (306:2.5) -- (18:1) -- cycle;
% USA
\draw[very thick,popgreen] (90:2.5) -- (162:1.5) -- (234:2.5) -- (306:2.5) -- (18:3) -- cycle;
% Légende
\node[popblue] at (-3.6,-4.4) {\textbf{--- France}};
\node[poppink] at (0,-4.4) {\textbf{--- Chine}};
\node[popgreen] at (3.4,-4.4) {\textbf{--- USA}};
\end{tikzpicture}
\end{center}
\legende{Graphique radar : profils comparés de coopération de trois partenaires (échelle indicative de 1 à 4). La France domine en APD et en alignement politique mais avec de fortes conditionnalités ; la Chine domine en commerce et transferts d'infrastructures avec peu de conditionnalités politiques ; les USA se distinguent par la santé et la sécurité.}
\end{popfigure}

\begin{exoresolu}[Comparer deux partenariats]
\textbf{Énoncé.} Un élève affirme : « La Chine et la France font exactement la même chose au Cameroun, seul le drapeau change. » À l'aide du cours, rédigez un paragraphe argumenté (8 à 12 lignes) pour discuter cette affirmation.
\tcblower
\textbf{Solution détaillée.} L'affirmation est excessive car les deux partenariats diffèrent par leur nature, leurs instruments et leurs conditions. \textbf{Nature} : la relation avec la France est une coopération Nord-Sud, héritée de l'histoire (accords de défense de 1974, AFD, entreprises comme Bolloré ou TotalEnergies) ; la relation avec la Chine est une coopération Sud-Sud récente, fondée sur un échange présenté comme « gagnant-gagnant ». \textbf{Instruments} : la France privilégie l'aide publique au développement et les investissements privés, tandis que la Chine privilégie les prêts concessionnels d'Exim Bank finançant de grandes infrastructures (barrages de Memve'ele et Lom Pangar, autoroute Kribi--Lolabé), souvent construites par des entreprises chinoises et garanties par les ressources naturelles. \textbf{Conditions} : la coopération française s'accompagne de conditionnalités de gouvernance et entretient le débat sur la Françafrique ; la coopération chinoise affiche le principe de non-ingérence mais fait peser un risque d'endettement (environ 40 à 50 \% de la dette extérieure publique). \textbf{Conclusion} : les deux partenaires financent le développement du Cameroun, mais selon des logiques différentes ; l'affirmation de l'élève ignore ces différences de nature, d'instruments et de conditions.
\end{exoresolu}

\begin{aretenir}[L'essentiel de la section]
\begin{itemize}
\item La \textbf{France} est le premier partenaire historique (défense, AFD, grandes entreprises), mais la relation est marquée par le contentieux de la Françafrique ;
\item Le \textbf{Royaume-Uni} joue la carte de la discrétion : bourses Chevening, formation du BIR, investissements ciblés ;
\item La \textbf{Chine} est le premier partenaire commercial et le premier financeur d'infrastructures, via les prêts concessionnels d'Exim Bank ; sa part dans la dette extérieure (40 à 50 \%) doit être citée avec sa référence exacte ;
\item Le \textbf{Japon} mise sur la qualité (JICA, pont sur le Wouri, ABE Initiative) ;
\item Les \textbf{USA} dominent dans la sécurité et la santé (PEPFAR, PMI) ; l'AGOA reste sous-exploitée faute de capacité productive ;
\item Le \textbf{Brésil} illustre la coopération Sud-Sud agricole (Embrapa, coton, manioc, bioéthanol).
\end{itemize}
\end{aretenir}

\begin{exercice}[Pour s'entraîner]
\begin{enumerate}
\item Définis la coopération Sud-Sud et donne deux exemples tirés du cours.
\item Le Cameroun contracte un prêt Exim Bank de 100 milliards de FCFA au taux de 2 \%, sur 20 ans, avec 6 ans de délai de grâce, garanti par les recettes d'un barrage. Analyse cet accord en suivant la méthode en quatre étapes.
\item Construis un tableau comparatif France / Chine / USA avec quatre lignes : domaines privilégiés, instruments, conditionnalités, limites.
\end{enumerate}
\end{exercice}

\begin{corrige}[Exercices de la section]
\begin{enumerate}
\item La coopération Sud-Sud est une coopération entre pays en développement ou émergents, présentée comme un partenariat « gagnant-gagnant » sans conditionnalités politiques. Exemples : Cameroun--Chine (infrastructures contre ressources) et Cameroun--Brésil (agriculture tropicale avec Embrapa).
\item \textbf{Coût} : 2 \%, taux concessionnel (inférieur au marché) ; \textbf{durée} : maturité de 20 ans avec 6 ans de délai de grâce, ce qui laisse le temps au barrage de produire de l'électricité et des revenus ; \textbf{garantie} : les recettes du barrage servent de sûreté, ce qui expose le pays en cas de recettes insuffisantes ; \textbf{conclusion} : prêt soutenable seulement si le barrage génère des revenus suffisants et si les travaux bénéficient aussi à l'économie locale.
\item Tableau attendu : France (APD, défense, entreprises / AFD, IDE / conditionnalités de gouvernance / Françafrique, dette) ; Chine (infrastructures, commerce / prêts Exim Bank / non-ingérence affichée / endettement, main-d'œuvre importée) ; USA (sécurité, santé / PEPFAR, PMI, USAID, AGOA / exigences de gouvernance / AGOA sous-exploité).
\end{enumerate}
\end{corrige}

\coursec{L'aide au développement : acteurs, instruments, efficacité et conditionnalités (Dossier 3)}

Après avoir analysé les relations bilatérales du Cameroun avec ses grands partenaires historiques et émergents (France, Royaume-Uni, Chine, Japon, USA, Brésil), il convient d'examiner l'architecture globale du financement du développement. L'aide publique au développement (APD) ne constitue qu'une pièce — bien que centrale — d'un ensemble plus vaste incluant les investissements privés, les transferts de la diaspora et la mobilisation des ressources internes. Cette section décrypte la typologie des flux, les instruments de mise en œuvre, le cadre des conditionnalités, les principes internationaux d'efficacité et l'émergence de nouveaux acteurs, pour poser \emph{in fine} la question de la souveraineté du développement camerounais.

\courssub{Typologie des flux financiers au développement}

L'aide au développement ne se résume pas aux dons. Elle s'inscrit dans un continuum de ressources externes aux caractéristiques financières, juridiques et politiques distinctes. La distinction fondamentale oppose les flux \textbf{concessionnels} (intégrant un don implicite via un taux d'intérêt inférieur à celui du marché ou une durée de remboursement longue) aux flux \textbf{non concessionnels} (conditions de marché).

\begin{definition}[Aide Publique au Développement (APD)]
Flux financiers officiels (États, collectivités, institutions multilatérales) vers les pays en développement (liste du CAD/OCDE), dont l'objectif principal est la promotion du développement économique et du bien-être social. Elle comprend des \textbf{dons} (subventions sans remboursement) et des \textbf{prêts concessionnels} (contenant un élément de don d'au moins 25\,\% : taux d'intérêt faible, différés d'amortissement longs).
\end{definition}

\begin{definition}[Appui Budgétaire Sectoriel (ABS)]
Transfert de ressources financières au Trésor public d'un pays partenaire, non affecté à un projet précis, mais conditionné à la mise en œuvre de réformes sectorielles et à l'atteinte d'indicateurs de performance (exemples : taux d'achèvement du primaire, taux de vaccination, taux d'exécution des budgets décentralisés). L'ABS vise le renforcement des systèmes publics (budget, comptabilité, passation des marchés) plutôt que la création de structures parallèles.
\end{definition}

\begin{itemize}
    \item \textbf{APD :} sources bilatérales (France/AFD, Allemagne/GIZ et KfW, Japon/JICA, USA/USAID, Chine) et multilatérales (Banque mondiale/IDA, Banque africaine de développement/FAD, Union européenne, FMI, agences des Nations unies).
    \item \textbf{Aide non concessionnelle (prêts commerciaux / Eurobonds) :} emprunts sur les marchés financiers internationaux (Eurobonds du Cameroun en 2015, puis opérations de refinancement ultérieures) ou prêts bancaires syndiqués. Taux de marché, maturités plus courtes. Utilisés pour financer de grandes infrastructures (barrages, routes, ports) ou refinancer la dette. Risques : alourdissement du service de la dette, exposition au risque de change.
    \item \textbf{Investissements Directs Étrangers (IDE) :} apport en capital durable (au moins 10\,\% des droits de vote) par une entreprise étrangère (création de filiale, fusion-acquisition, réinvestissement des bénéfices). Moteur de transfert de technologie, d'emplois et d'accès aux marchés. Au Cameroun : secteur pétrolier (historique), télécommunications (Orange, MTN), agro-industrie (CDC, HEVECAM, SOCAPALM), énergie (centrale hydroélectrique de Nachtigal), mines. Flux volatils, sensibles au climat des affaires et au rapatriement des bénéfices.
    \item \textbf{Transferts de la diaspora :} envois de fonds (\emph{remittances}) des Camerounais de l'étranger (France, USA, Allemagne, Belgique, pays du Golfe, Afrique du Sud). Estimation Banque mondiale : de l'ordre de \textbf{300 milliards de FCFA par an}, soit l'une des premières sources de devises du pays. Ils financent la consommation, le logement, l'éducation, la santé et les PME familiales. Ils sont stables et contre-cycliques (ils augmentent en cas de crise au pays).
\end{itemize}

\begin{propriete}[Stabilité des transferts de la diaspora]
Les transferts de la diaspora sont plus stables et contre-cycliques que l'APD ou les IDE. En période de choc (crise économique, catastrophe naturelle, conflit, pandémie de COVID-19), l'APD peut être gelée ou redirigée et les IDE se contractent (aversion au risque), alors que la diaspora augmente ses envois pour soutenir les familles restées au pays. Ces transferts constituent un filet de sécurité sociale privé massif et non conditionnel.
\end{propriete}

\begin{exemplebox}[Structure des flux externes nets reçus par le Cameroun (ordres de grandeur annuels moyens récents, milliards de FCFA)]
\begin{center}
\begin{tabularx}{\linewidth}{|l|X|X|}\hline
\textbf{Type de flux} & \textbf{Montant estimé} & \textbf{Caractère} \\ \hline
APD nette (dons + prêts concessionnels) & 400 -- 500 & Concessionnelle, conditionnelle, sectorielle \\ \hline
Transferts de la diaspora & environ 300 & Privés, stables, contre-cycliques, non conditionnels \\ \hline
Recettes pétrolières (export net) & 500 -- 800 & Propres, volatiles (prix du baril), exhaustibles \\ \hline
IDE (flux nets entrants) & 200 -- 400 & Privés, lucratifs, sectoriellement ciblés \\ \hline
Emprunts commerciaux (Eurobonds, banques) & Variable (gros montants ponctuels) & Non concessionnels, remboursables, risqués \\ \hline
\end{tabularx}
\end{center}
\legende{Sources : BEAC, Banque mondiale, OCDE/CAD, MINFI. Ordres de grandeur annuels moyens récents, donnés à titre indicatif.}
\end{exemplebox}

\courssub{Alignement sur les agendas globaux et nationaux : ODD 2030 et SND30}

L'efficacité de l'aide suppose qu'elle ne soit pas un « projet de donateur » greffé artificiellement, mais qu'elle soutienne la \textbf{stratégie nationale de développement}. Le Cameroun a aligné sa planification sur deux cadres majeurs :

\begin{itemize}
    \item \textbf{Agenda 2030 / ODD (Objectifs de Développement Durable) :} 17 objectifs et 169 cibles adoptés par l'ONU en 2015. Cadre universel de référence pour tous les partenaires. Le Cameroun présente des rapports nationaux volontaires au Forum politique de haut niveau des Nations unies.
    \item \textbf{SND30 (Stratégie Nationale de Développement 2020-2030) :} référentiel unique de la politique de développement du Cameroun. Vision : faire du Cameroun « une nation émergente, démocratique et unie dans sa diversité ». Quatre piliers : (1) transformation structurelle de l'économie, (2) développement du capital humain et bien-être, (3) promotion de l'emploi et de l'insertion économique, (4) gouvernance, décentralisation et gestion stratégique de l'État. Elle est déclinée en plans triennaux et en budgets annuels.
\end{itemize}

L'alignement se traduit notamment par les \textbf{Appuis Budgétaires Sectoriels (ABS)} : l'Union européenne (éducation, gouvernance, décentralisation), la Banque mondiale (santé, protection sociale, éducation), la BAfD (eau, énergie, transport) et l'AFD (éducation, formation professionnelle, développement urbain, climat) transfèrent des ressources au Trésor contre des \textbf{indicateurs de performance} liés au décaissement.
\emph{Exemples :} un appui au secteur de l'éducation peut conditionner le décaissement au taux d'achèvement du primaire, au ratio élèves/enseignant qualifié ou à la disponibilité des manuels scolaires ; un appui à la santé peut viser le taux de vaccination ou la proportion d'accouchements assistés.

\begin{attention}[L'aide « liée » : un coût caché majeur]
Une part de l'APD bilatérale reste « liée » : le pays bénéficiaire est obligé d'acheter des biens, services, expertises ou équipements \textbf{exclusivement dans le pays donateur}. Cela renchérit le coût réel des projets (surcoût estimé de 15 à 30\,\% selon l'OCDE : prix plus élevés, absence de concurrence, inadaptation au contexte), réduit l'effet de levier sur l'économie locale (pas de commandes aux PME camerounaises) et affaiblit l'appropriation. Le Cameroun milite, avec le groupe ACP, pour le « déliement » de l'aide dans les négociations internationales (accords ACP-UE, partenariat UE-Afrique).
\end{attention}

\courssub{Instruments de mise en œuvre : du projet à l'appui budgétaire}

Le choix de l'instrument conditionne la gouvernance, la durabilité et l'appropriation nationale de l'aide.

\begin{center}
\begin{tabularx}{\linewidth}{|l|X|X|X|}\hline
\textbf{Instrument} & \textbf{Description} & \textbf{Gestion} & \textbf{Avantages / Limites} \\ \hline
\textbf{Projet} & Opération ponctuelle et ciblée (école, hôpital, route), durée définie (3 à 5 ans), budget propre. & Unité de Gestion de Projet (UGP), souvent parallèle à l'administration. & + Visible, résultats rapides, contrôle fort du donateur. -- Parallélisme (affaiblit les systèmes publics), coûts de transaction élevés, durabilité faible (maintenance ?), fragmentation. \\ \hline
\textbf{Programme sectoriel (approche sectorielle)} & Cadre unique de politique sectorielle (ex. : plan sectoriel de l'éducation), budget commun (État + partenaires), rapport unique, dialogue structuré. & Ministère sectoriel (MINESEC, MINSANTE, MINEPAT) avec appui technique. & + Alignement, harmonisation, renforcement des systèmes, appropriation. -- Lourd à mettre en place, exige une forte capacité de pilotage, lenteur des décaissements. \\ \hline
\textbf{Appui Budgétaire Général (ABG)} & Transfert non affecté au Trésor, conditionné à des réformes macroéconomiques et transversales (gouvernance, gestion des finances publiques). & Trésor public (MINFI) / MINEPAT. & + Souveraineté budgétaire, flexibilité, alignement total, prévisibilité. -- Risque de fongibilité (fuite vers des dépenses non prioritaires), conditionnalité politique forte, suspension en cas de rupture du dialogue. \textbf{Quasi disparu au Cameroun} depuis la fin des années 2000 (tensions sur la gouvernance). \\ \hline
\textbf{Appui Budgétaire Sectoriel (ABS)} & Transfert au Trésor conditionné à des indicateurs sectoriels. Forme hybride entre projet et ABG. & Ministère sectoriel + MINFI (comptabilité). & + Renforce les systèmes publics (budget, comptabilité, marchés), alignement sur la SND30, prévisibilité. -- Exige un système fiable de gestion des finances publiques, auditables, et des indicateurs mesurables. \\ \hline
\textbf{Garanties} & Instrument de réduction du risque (politique, change, rupture de contrat) pour mobiliser des capitaux privés (banques, obligataires) vers les infrastructures et les partenariats public-privé. & Partenaire privé / PPP. & + Effet de levier, partage du risque, pas de don direct. -- Complexité juridique, coût de la prime de garantie, risque moral. \\ \hline
\end{tabularx}
\end{center}

\courssub{Conditionnalités : le prix de l'aide}

L'aide n'est jamais politiquement « gratuite ». Les conditionnalités structurent la relation entre donateur et bénéficiaire.

\begin{itemize}
    \item \textbf{Conditionnalités politiques / de gouvernance :} respect des droits de l'homme, État de droit, démocratie, transparence électorale, lutte contre la corruption, liberté de la presse et d'association. Elles figurent dans les accords-cadres (accord de Cotonou ACP-UE, article 9 sur les « éléments essentiels » ; accord post-Cotonou de Samoa, 2023). Leur non-respect peut entraîner la suspension ou la révision de la coopération.
    \item \textbf{Conditionnalités économiques (programmes FMI) :} ancrage macroéconomique. Le programme appuyé par la \textbf{Facilité élargie de crédit (FEC) 2021-2025} comporte des conditionnalités quantitatives (plafonds du déficit budgétaire, de la dette, des réserves nettes, non-accumulation d'arriérés) et structurelles (mobilisation des recettes non pétrolières, gestion des dépenses, transparence du secteur pétrolier et de la SNH, secteur financier, climat des affaires, gouvernance des entreprises publiques), avec des revues semestrielles. L'accord du FMI sert de « sceau d'approbation » indispensable au déblocage des financements des autres partenaires (Banque mondiale, BAfD, UE, AFD, marchés obligataires).
    \item \textbf{Conditionnalités environnementales et sociales (sauvegardes) :} obligatoires pour les projets de la Banque mondiale (Cadre environnemental et social, 2018), de la BAfD et de l'AFD. Études d'impact environnemental et social, plans de gestion, plans de réinstallation, consultation des populations affectées. Leur non-respect peut entraîner la suspension des décaissements.
\end{itemize}

\courssub{Efficacité de l'aide : des principes de Paris à la question de la maîtrise}

L'histoire de l'aide est jalonnée de déclarations internationales visant à en améliorer l'efficacité. Le Cameroun y a souscrit.

\begin{aretenir}[Jalons de l'efficacité de l'aide (principes directeurs)]
\begin{itemize}
    \item \textbf{Déclaration de Paris (2005) :} cinq piliers : \textbf{appropriation} (le pays pilote son développement), \textbf{alignement} (les donateurs utilisent les systèmes du pays), \textbf{harmonisation} (coordination entre donateurs), \textbf{gestion axée sur les résultats}, \textbf{responsabilité mutuelle}.
    \item \textbf{Agenda d'Accra (2008) :} accent sur la \textbf{prévisibilité} (engagements pluriannuels), l'utilisation des \textbf{systèmes du pays} (finances publiques, passation des marchés, statistiques), des conditionnalités moins nombreuses et plus pertinentes, et le rôle de la \textbf{société civile} et des parlementaires.
    \item \textbf{Partenariat de Busan (2011) :} partenariat mondial pour une coopération efficace au service du développement. Élargissement aux acteurs non traditionnels : Chine, Brésil, fondations, secteur privé, société civile. Principes : appropriation, priorité aux résultats, partenariats inclusifs, transparence et responsabilité mutuelle. Fin de la stricte dichotomie donateur/bénéficiaire.
\end{itemize}
\end{aretenir}

\begin{methode}[Grille d'analyse d'un projet ou dispositif d'aide (critères de l'OCDE/CAD)]
\textbf{Étape 1 -- Pertinence :} l'intervention répond-elle aux besoins réels des populations cibles et aux priorités de la SND30 et des plans sectoriels ?
\textbf{Étape 2 -- Cohérence :} est-elle compatible avec les autres interventions (autres donateurs, politiques nationales, ODD) ? Y a-t-il synergies, doublons ou contradictions ?
\textbf{Étape 3 -- Efficacité :} les objectifs spécifiques sont-ils atteints (indicateurs de résultats) ?
\textbf{Étape 4 -- Efficience :} le rapport coûts/résultats est-il optimal (coût unitaire, délais, gaspillage, coûts de transaction, aide liée) ?
\textbf{Étape 5 -- Impact :} quels changements durables (positifs ou négatifs, directs ou indirects) sur le bien-être, la pauvreté, les inégalités, l'environnement, le genre ?
\textbf{Étape 6 -- Durabilité :} les bénéfices se maintiennent-ils après la fin du financement externe (appropriation locale, capacité de maintenance, financement récurrent par le budget de l'État) ?
\textbf{Étape 7 -- Responsabilité mutuelle :} le donateur rend-il des comptes à ses contribuables, le bénéficiaire à son Parlement et à ses citoyens ? Les données sont-elles transparentes ?
\end{methode}

\begin{savaistu}[Le Cameroun et l'initiative PPTE (Pays Pauvres Très Endettés)]
Le Cameroun a atteint le \textbf{point d'achèvement de l'initiative PPTE en avril 2006}. Cela a entraîné l'annulation d'une part très importante de sa dette extérieure (plusieurs milliers de milliards de FCFA en valeur nominale, soit une réduction de l'endettement extérieur d'environ 80\,\% en valeur actualisée) par les créanciers bilatéraux (Club de Paris) et multilatéraux (FMI, Banque mondiale, BAfD). Cette bouffée d'oxygène a permis de réorienter des ressources vers l'éducation (gratuité du primaire), la santé (paludisme, VIH) et les routes rurales. Cependant, l'endettement a fortement remonté depuis (Eurobonds, dette intérieure, prêts chinois), posant à nouveau la question de la soutenabilité : le FMI classe désormais le risque de surendettement du Cameroun comme \textbf{élevé}.
\end{savaistu}

\courssub{Nouveaux acteurs et modalités innovantes}

Le paysage de la coopération ne se résume plus aux pays du CAD/OCDE (donateurs traditionnels).

\begin{itemize}
    \item \textbf{Pays émergents / coopération Sud-Sud :} \textbf{Chine} (prêts de l'Exim Bank of China pour les infrastructures lourdes : barrages, routes, stades, palais des congrès ; absence de conditionnalité politique de gouvernance, mais clauses de confidentialité et recours à des entreprises chinoises), \textbf{Brésil} (coopération technique agricole via l'EMBRAPA, santé tropicale), \textbf{Inde} (lignes de crédit, médicaments génériques, TIC, formations), \textbf{Turquie} (agence TIKA, construction, santé, éducation), \textbf{pays du Golfe} (Fonds saoudien et koweïtien de développement, BADEA, Banque islamique de développement).
    \item \textbf{Nouvelles banques de développement :} la \textbf{Nouvelle Banque de Développement (NBD, banque des BRICS)} et la \textbf{Banque Asiatique d'Investissement pour les Infrastructures (AIIB)} offrent aux pays africains de nouvelles sources de financement des infrastructures (énergie, transport, eau), avec des conditionnalités différentes de celles du FMI. Le Cameroun suit de près ces mécanismes dans le cadre de la diversification de ses partenaires.
    \item \textbf{Fondations philanthropiques mondiales :} \textbf{Fondation Gates} (vaccins via Gavi, poliomyélite, paludisme, assainissement, agriculture), \textbf{Clinton Health Access Initiative} (négociation des prix des médicaments), \textbf{Fondation Mastercard} (emploi des jeunes, formation, inclusion financière). Modalités : subventions directes aux ONG et au gouvernement, partenariats public-privé. Avantages : agilité, résultats mesurables. Limites : fragmentation, contournement des systèmes du pays, agenda propre.
    \item \textbf{ONG internationales (MSF, Action contre la Faim, Plan International, CARE, CRS, World Vision, IRC, NRC) :} opérationnelles dans l'urgence (Extrême-Nord, Nord-Ouest/Sud-Ouest, région de l'Est accueillant des réfugiés centrafricains) et le développement. Elles complètent, parfois se substituent à l'État dans les zones fragiles. Coordination via le système des « clusters » humanitaires (ONU/OCHA).
    \item \textbf{Financement innovant et climatique :} Fonds vert pour le climat (FVC), Fonds d'adaptation, mécanismes de marché carbone de l'Accord de Paris (article 6), obligations vertes, finance carbone forestière (REDD+).
\end{itemize}

\begin{exoresolu}[Analyse de l'alignement d'un projet hypothétique sur les principes de Busan]
\textbf{Énoncé :} Une banque multilatérale finance un « Projet de Développement des Compétences pour l'Employabilité (PDCE) » au Cameroun (2023-2028, 100 millions USD). Le projet crée une UGP distincte au sein du MINEFOP, recrute son propre personnel (rémunéré sur les fonds du projet), utilise ses propres manuels de procédures (passation des marchés, gestion financière), rend compte trimestriellement au bailleur, et forme 5\,000 jeunes dans des centres privés sélectionnés par appel d'offres international. L'État camerounais ne gère pas les fonds. Évaluer l'alignement de ce projet sur les principes de Busan (appropriation, alignement, harmonisation, gestion axée sur les résultats, responsabilité mutuelle).
\tcblower
\textbf{Solution détaillée :}
\begin{enumerate}
    \item \textbf{Appropriation (faible) :} le pays n'est pas « au volant ». L'UGP parallèle affaiblit le MINEFOP. Les décisions clés (choix des centres, des cursus, des bénéficiaires) sont validées par le bailleur. La SND30 et le plan sectoriel emploi-formation ne sont pas le point de départ unique.
    \item \textbf{Alignement (faible) :} utilisation de systèmes financiers et de passation des marchés \textbf{parallèles} (procédures du bailleur) au lieu des systèmes nationaux (code des marchés publics, comptabilité du Trésor). Non-alignement sur le budget-programme de l'État.
    \item \textbf{Harmonisation (faible) :} projet « autonome », sans fonds communs avec les autres partenaires de la formation professionnelle (AFD, GIZ). Risque de doublons (mêmes centres, mêmes jeunes). Absence de dialogue sectoriel unique.
    \item \textbf{Gestion axée sur les résultats (moyenne) :} un cadre de résultats existe (nombre de formés, taux d'insertion à 6 mois), mais les indicateurs portent surtout sur les « produits » (personnes formées) plutôt que sur l'« impact » (emploi durable, revenu). Les données sont produites par l'UGP et non validées par l'INS ou le ministère.
    \item \textbf{Responsabilité mutuelle (faible) :} la redevabilité s'exerce principalement envers le conseil d'administration du bailleur. La redevabilité envers le Parlement camerounais, la Cour des comptes et les citoyens est quasi inexistante.
\end{enumerate}
\textbf{Conclusion :} il s'agit d'un projet « à l'ancienne » (approche projet classique), peu aligné sur les principes de Paris et de Busan. \textbf{Recommandation :} basculer vers un appui budgétaire sectoriel ou une approche programmatique sectorielle : fonds transitant par le Trésor, indicateurs de décaissement (ex. : taux d'insertion certifié par l'INS, nombre de centres publics accrédités), utilisation des systèmes nationaux, dialogue sectoriel unique piloté par le MINEPAT et le MINEFOP.
\end{exoresolu}

\courssub{Visualisation des flux et de la dépendance}

\begin{popfigure}
\centering
\begin{tikzpicture}[
    source/.style={draw=popblue, thick, fill=popblue!10, rounded corners, minimum width=3cm, minimum height=1cm, font=\small\bfseries, align=center},
    sector/.style={draw=popgreen, thick, fill=popgreen!10, rounded corners, minimum width=3cm, minimum height=1cm, font=\small\bfseries, align=center},
    arrow/.style={-Stealth, thick, draw=popdark},
    lab/.style={font=\tiny, text=popdark, midway, fill=white, inner sep=1pt}
]
% Sources (gauche)
\node[source, align=center] (apd_bil) at (0, 5){APD bilatérale\\ (France, Allemagne,\\ Japon, USA, Chine)};
\node[source, align=center] (apd_mul) at (0, 3){APD multilatérale\\ (Banque mondiale, BAfD,\\ UE, FMI, ONU)};
\node[source, align=center] (ide) at (0, 1){IDE\\ (pétrole, télécoms,\\ agro-industrie, énergie)};
\node[source, align=center] (dias) at (0, -1){Diaspora\\ (environ 300 Mds FCFA/an)};
\node[source, align=center] (own) at (0, -3){Recettes propres\\ (pétrole, impôts,\\ douanes, forêts)};

% Secteurs (droite)
\node[sector, align=center] (infra) at (10, 5){Infrastructures\\ (routes, barrages,\\ ports, énergie, numérique)};
\node[sector, align=center] (social) at (10, 2.5){Social\\ (éducation, santé,\\ protection sociale, eau)};
\node[sector, align=center] (gov) at (10, 0){Gouvernance\\ (finances publiques,\\ décentralisation, justice)};
\node[sector, align=center] (prod) at (10, -2.5){Secteur productif\\ (agro-industrie, PME,\\ mines, industrie)};

% Flux APD bilatérale
\draw[arrow, popblue!70, line width=4pt] (apd_bil.east) -- ++(1,0) |- (infra.west) node[lab, pos=0.7, above] {fort};
\draw[arrow, popblue!70, line width=3pt] (apd_bil.east) -- ++(1,0) |- (social.west) node[lab, pos=0.7, above] {moyen};
\draw[arrow, popblue!70, line width=2pt] (apd_bil.east) -- ++(1,0) |- (gov.west) node[lab, pos=0.7, above] {faible};

% Flux APD multilatérale
\draw[arrow, poppurple!70, line width=4pt] (apd_mul.east) -- ++(1,0) |- (social.west) node[lab, pos=0.3, below] {très fort};
\draw[arrow, poppurple!70, line width=3pt] (apd_mul.east) -- ++(1,0) |- (gov.west) node[lab, pos=0.3, below] {fort};
\draw[arrow, poppurple!70, line width=2pt] (apd_mul.east) -- ++(1,0) |- (infra.west);
\draw[arrow, poppurple!70, line width=1.5pt] (apd_mul.east) -- ++(1,0) |- (prod.west);

% Flux IDE
\draw[arrow, poporange!70, line width=5pt] (ide.east) -- ++(1,0) |- (infra.west) node[lab, pos=0.7, above] {majeur};
\draw[arrow, poporange!70, line width=3pt] (ide.east) -- ++(1,0) |- (prod.west) node[lab, pos=0.7, above] {croissant};

% Flux diaspora
\draw[arrow, poppink!70, line width=3pt] (dias.east) -- ++(1,0) |- (social.west) node[lab, pos=0.5, below] {consommation, santé, éducation};
\draw[arrow, poppink!70, line width=2pt] (dias.east) -- ++(1,0) |- (prod.west) node[lab, pos=0.5, below] {PME familiales};
\draw[arrow, poppink!70, line width=1.5pt] (dias.east) -- ++(1,0) |- (infra.west) node[lab, pos=0.5, below] {logement};

% Recettes propres
\draw[arrow, popgreen!70, line width=6pt] (own.east) -- ++(1,0) |- (infra.west);
\draw[arrow, popgreen!70, line width=5pt] (own.east) -- ++(1,0) |- (social.west);
\draw[arrow, popgreen!70, line width=4pt] (own.east) -- ++(1,0) |- (gov.west);
\draw[arrow, popgreen!70, line width=3pt] (own.east) -- ++(1,0) |- (prod.west);

% Légende
\node[font=\small\bfseries, anchor=north west] at (0, -4.5) {Épaisseur des flèches = volume financier relatif};
\node[font=\tiny, anchor=north west, align=left] at (0, -5) {Bleu : APD bilatérale \quad Violet : APD multilatérale \quad Orange : IDE\\ Rose : diaspora \quad Vert : recettes propres};
\end{tikzpicture}
\legende{Diagramme de flux simplifié (type Sankey) : sources de financement du développement vers les secteurs bénéficiaires au Cameroun. L'épaisseur des flèches représente l'ordre de grandeur des flux financiers annuels.}
\end{popfigure}

\begin{popfigure}
\centering
\begin{tikzpicture}
\begin{axis}[
    width=\linewidth,
    height=9cm,
    grid=major,
    grid style={popline},
    xlabel={Année},
    ylabel={Ratio (\%)},
    legend pos=north east,
    legend cell align=left,
    tick label style={font=\small},
    label style={font=\small},
    xmin=2000, xmax=2023,
    xtick={2000,2005,2010,2015,2020,2023},
    ymin=0, ymax=30,
    axis line style={popdark},
    tick style={popdark},
]
% APD / PIB : baisse tendancielle de ~6 % à ~2-3 %
\addplot[popcurve, mark=*, mark size=2pt, mark options={fill=popblue}] coordinates {
    (2000, 5.8) (2001, 5.5) (2002, 5.2) (2003, 5.0) (2004, 4.8) (2005, 4.5)
    (2006, 4.0) (2007, 3.8) (2008, 3.5) (2009, 3.8) (2010, 3.6)
    (2011, 3.4) (2012, 3.2) (2013, 3.0) (2014, 2.9) (2015, 2.8)
    (2016, 2.7) (2017, 2.6) (2018, 2.5) (2019, 2.4) (2020, 2.8)
    (2021, 2.6) (2022, 2.4) (2023, 2.3)
};
\addlegendentry{APD / PIB (\%)}

% APD / Budget de l'État : baisse plus marquée
\addplot[poporange, mark=square*, mark size=2pt, mark options={fill=poporange}, thick] coordinates {
    (2000, 28) (2001, 26) (2002, 24) (2003, 23) (2004, 22) (2005, 20)
    (2006, 18) (2007, 17) (2008, 16) (2009, 18) (2010, 17)
    (2011, 16) (2012, 15) (2013, 14) (2014, 13) (2015, 12)
    (2016, 12) (2017, 11) (2018, 11) (2019, 10) (2020, 13)
    (2021, 12) (2022, 11) (2023, 10)
};
\addlegendentry{APD / Budget de l'État (\%)}

% Annotations
\node[font=\tiny, popblue, anchor=south] at (axis cs:2006,4.0) {Point d'achèvement PPTE};
\node[font=\tiny, poporange, anchor=north] at (axis cs:2020,13) {Choc COVID-19};
\node[font=\tiny, popgreen, anchor=south] at (axis cs:2023,10) {Objectif SND30 : moins de 10\,\%};
\end{axis}
\end{tikzpicture}
\legende{Évolution indicative des ratios APD/PIB et APD/Budget de l'État (2000-2023). Tendance structurelle à la baisse de la dépendance à l'aide, malgré le choc conjoncturel de la COVID-19 (2020). Source : élaboration pédagogique à partir des données Banque mondiale, FMI (rapports article IV) et MINFI (lois de finances).}
\end{popfigure}

\courssub{Problèmes et défis : le Cameroun est-il maître de son développement ? (Transition vers le Dossier 5)}

L'analyse de l'aide révèle des tensions structurelles non résolues, qui constituent le cœur du \textbf{Dossier 5 : les problèmes du Cameroun dans les relations internationales} :

\begin{itemize}
    \item \textbf{Fragmentation et coûts de transaction :} plus de 30 partenaires techniques et financiers, plus d'une centaine de projets actifs, de multiples UGP, des procédures distinctes et des missions d'évaluation simultanées. La charge administrative est écrasante pour le MINEPAT, le MINFI et les ministères sectoriels.
    \item \textbf{Volatilité et imprévisibilité :} les décaissements sont souvent inférieurs aux engagements (taux d'exécution fréquemment de 60 à 70\,\%). Suspensions politiques (gouvernance, droits humains) et conditionnalités du FMI en « stop-and-go ».
    \item \textbf{Aide liée et fuite de valeur :} surcoûts de 15 à 30\,\%, faibles retombées locales (emplois, PME, transfert de compétences).
    \item \textbf{Endettement croissant et risque de soutenabilité :} passage du statut de pays ayant bénéficié de l'initiative PPTE (2006) à un risque de surendettement jugé élevé par le FMI. Le service de la dette absorbe une part croissante des recettes de l'État. Accumulation d'Eurobonds, de dette chinoise (clauses parfois opaques) et de dette intérieure (arriérés envers les PME).
    \item \textbf{Appropriation limitée :} politiques publiques parfois « écrites pour les donateurs » (documents de stratégie, matrices de réformes du FMI) plutôt que portées par le leadership politique national et le Parlement. Évaluation \emph{ex post} faible.
    \item \textbf{Conditionnalité asymétrique :} exigence de transparence et de gouvernance du côté du bénéficiaire, mais opacité des motivations politiques des donateurs (géopolitique, intérêts commerciaux, migrations).
    \item \textbf{Négligence du financement interne :} pression fiscale faible (environ 12 à 13\,\% du PIB), niches d'exonérations massives, évasion et flux financiers illicites, gestion perfectible des rentes pétrolières et minières. L'aide peut désinciter à la réforme fiscale interne (effet de substitution).
\end{itemize}

\begin{exercice}[Analyse critique d'un document de stratégie de coopération]
\textbf{Document :} extrait d'un document-cadre de coopération avec le Cameroun (par exemple le cadre de partenariat pays de la Banque mondiale 2022-2026, ou le cadre de coopération des Nations unies 2022-2026).
\textbf{Consigne :} à partir de la grille OCDE/CAD (méthode du 5.5) et des principes de Busan, identifier 3 forces et 3 faiblesses de l'alignement de ce document sur la SND30 et la souveraineté nationale. Proposer 2 mesures concrètes pour renforcer l'appropriation camerounaise.
\end{exercice}

\begin{corrige}[Exercice n°1 -- Pistes de correction]
\textbf{Forces :} (1) alignement explicite sur les piliers de la SND30 (éducation, santé, gouvernance, climat) ; (2) recours à l'approche programme et aux appuis budgétaires sectoriels dans certains secteurs (éducation, décentralisation) ; (3) association de la société civile et du Parlement au dialogue (comités de pilotage).
\textbf{Faiblesses :} (1) maintien d'UGP parallèles pour les projets « phares » (non-utilisation des systèmes nationaux) ; (2) conditionnalités politiques non négociables (clauses de suspension) asymétriques ; (3) part d'aide liée et recours aux experts des donateurs encore élevée ; (4) indicateurs portant sur les « produits » (activités réalisées) plutôt que sur l'« impact » (emploi, pauvreté).
\textbf{Mesures :} (1) fusionner les UGP sectorielles dans les secrétariats généraux et directions centrales des ministères (renforcement des capacités internes) ; (2) négocier un « pacte » unique pluriannuel (cadre de partenariat pays unique) remplaçant les accords bilatéraux multiples, piloté par le MINEPAT, avec des fonds communs (fonds fiduciaire multi-bailleurs) et une matrice de résultats unique alignée sur la SND30.
\end{corrige}

\coursec{Gestion des conflits et sécurité collective : Contribution camerounaise (TD2)}

Après avoir étudié, dans la section précédente, l'aide au développement comme instrument de la coopération multilatérale, nous abordons maintenant une dimension plus sensible de l'action internationale du Cameroun : la \cle{paix} et la \cle{sécurité}. En effet, le développement est impossible sans stabilité. Un pays qui reçoit des financements mais dont les frontières sont menacées ou dont le territoire est déstabilisé ne peut pas transformer cette aide en progrès durable. Le Cameroun l'a bien compris : il s'est imposé, depuis son indépendance, comme un acteur majeur de la \cle{gestion des conflits} en Afrique centrale, à la fois comme \emph{médiateur}, comme \emph{contributeur de troupes} et comme \emph{hôte diplomatique}. Ce travail dirigé (TD2) nous invite à évaluer cette contribution, ses succès, mais aussi ses limites et ses coûts.

\courssub{La doctrine de politique étrangère du Cameroun en matière de sécurité}

\begin{definition}[Doctrine de politique étrangère]
Une \cle{doctrine de politique étrangère} est l'ensemble des principes constants qui guident l'action d'un État sur la scène internationale. Elle traduit ses valeurs, ses intérêts et sa vision du monde, et permet de prévoir ses réactions face aux crises.
\end{definition}

L'action sécuritaire du Cameroun repose sur quatre piliers, hérités de la Charte de l'ONU (1945), de l'Acte constitutif de l'Union Africaine (2000) et du Protocole relatif à la création du Conseil de paix et de sécurité de l'UA (2002) :

\begin{itemize}
\item \textbf{La \cle{non-ingérence}} dans les affaires intérieures des autres États : le Cameroun refuse d'intervenir unilatéralement dans les crises internes de ses voisins et n'agit que sur demande, sur mandat régional (UA, CEEAC) ou dans le cadre d'accords multilatéraux (ONU).
\item \textbf{Le \cle{règlement pacifique des différends}} : priorité à la négociation, à la médiation, à l'arbitrage et au droit international plutôt qu'à la force. L'exemple emblématique est le conflit du Bakassi avec le Nigeria, réglé devant la Cour internationale de Justice (arrêt du 10 octobre 2002) puis par les Accords de Greentree (12 juin 2006).
\item \textbf{La \cle{solidarité africaine}} : un conflit chez un voisin menace toute la sous-région (réfugiés, trafics d'armes, insécurité transfrontalière, propagation du terrorisme). Le Cameroun s'engage donc dans les mécanismes collectifs (COPAX, FOMUC/MICOPAX, FMM).
\item \textbf{La défense de l'\cle{intégrité territoriale}} : la protection des frontières et de la souveraineté nationale est un devoir constitutionnel (préambule de la Constitution du 18 janvier 1996). C'est le fondement de l'engagement militaire contre Boko Haram dans l'Extrême-Nord et de la réponse à la crise du NOSO.
\end{itemize}

\begin{aretenir}[Les quatre piliers de la doctrine camerounaise]
Non-ingérence ; règlement pacifique des différends ; solidarité africaine ; défense de l'intégrité territoriale. Ces principes expliquent à la fois les médiations (volet diplomatique) et les interventions militaires (volet sécuritaire, toujours encadré par un mandat régional ou international).
\end{aretenir}

\courssub{Les médiations directes : le Cameroun, artisan de la paix}

Le Cameroun a acquis une réputation de \emph{médiateur crédible} en Afrique centrale, grâce à sa stabilité relative, sa position géographique centrale et son histoire de règlement pacifique de ses propres différends frontaliers.

\begin{itemize}
\item \textbf{Tchad (1979 et 2008)} : lors des guerres civiles tchadiennes, le Cameroun a accueilli des pourparlers et facilité la réconciliation entre factions, notamment en 1979 (conférences de Lagos et Kano auxquelles il participe activement) et en 2008 lors de l'offensive rebelle sur N'Djamena (soutien logistique, diplomatique et accueil de réfugiés).
\item \textbf{République centrafricaine (2008, 2013, 2019)} : le Cameroun a participé aux \emph{Accords de Libreville} (2008, dialogue politique inclusif) puis, en janvier 2013, à l'accord de cessez-le-feu de Libreville entre le régime Bozizé et la coalition Séléka. Yaoundé a ensuite accueilli des pourparlers inter-centrafricains et soutenu l'Accord de Khartoum (février 2019) sous l'égide de l'UA et de la CEEAC.
\item \textbf{Nigeria (Bakassi)} : après l'arrêt de la CIJ, la \emph{Commission mixte Cameroun-Nigeria} (CMCN), présidée par le représentant spécial du Secrétaire général de l'ONU et basée à Yaoundé, a supervisé le retrait de l'administration nigériane, la démarcation de la frontière terrestre et maritime, et le transfert de souveraineté (2006-2013). C'est un modèle mondial de règlement pacifique d'un conflit frontalier.
\item \textbf{Guinée équatoriale} : le Cameroun a facilité les négociations sur la délimitation des frontières maritimes dans le golfe de Guinée, zone riche en hydrocarbures, aboutissant à un accord de principe sur le tracé de la frontière maritime.
\end{itemize}

\begin{exemplebox}[La Commission mixte Cameroun-Nigeria : une médiation réussie]
Créée en novembre 2002 sous l'égide de l'ONU, la Commission mixte (siège à Yaoundé) a organisé le retrait de l'administration nigériane du Bakassi, la démarcation de la frontière terrestre (piliers) et maritime, et la protection des populations (Accords de Greentree, 12 juin 2006). Le 14 août 2013, le transfert définitif de souveraineté a eu lieu sans effusion de sang, alors que la presse annonçait une guerre certaine en 2002. Leçon : le droit international et la médiation patiente peuvent désamorcer les conflits les plus explosifs.
\end{exemplebox}

\courssub{L'engagement militaire régional : FOMAC et Force Mixte Multinationale}

\begin{definition}[Force Mixte Multinationale (FMM)]
La \cle{Force Mixte Multinationale} (FMM) est une force militaire régionale regroupant le Nigeria, le Niger, le Tchad, le Cameroun et le Bénin, créée en 1994 (réactivée et renforcée en 2015) sous mandat de l'Union Africaine et appuyée par la résolution 2349 (2017) du Conseil de sécurité de l'ONU, pour combattre le groupe terroriste Boko Haram et sa branche ISWAP dans le bassin du Lac Tchad. Son quartier général est à N'Djamena (Tchad).
\end{definition}

\textbf{La FOMAC en Centrafrique.} Dans le cadre de la CEEAC, le Cameroun a fourni des contingents à la Force multinationale de l'Afrique centrale : FOMUC (2002-2008) puis MICOPAX (2008-2013) et enfin des contributions à la MISCA (mission africaine) puis à la MINUSCA (mission onusienne) après la crise de 2013. Ces opérations relèvent du COPAX (Conseil de paix et de sécurité de l'Afrique centrale), étudié au dossier 4.

\textbf{La FMM contre Boko Haram.} Depuis 2009, la secte Boko Haram (puis sa branche État islamique en Afrique de l'Ouest, ISWAP) mène des attaques meurtrières dans l'Extrême-Nord du Cameroun (enlèvements, attentats-suicides, incursions, pillages). En 2014-2015, les pays du bassin du Lac Tchad décident de mutualiser leurs forces. La FMM est divisée en secteurs : le \emph{secteur 1} est commandé depuis Mora (Cameroun) et couvre l'Extrême-Nord camerounais. Les opérations conjointes (notamment 2015-2016) ont permis de reprendre la majeure partie des territoires occupés et de réduire fortement la capacité de nuisance du groupe, sans toutefois l'éliminer totalement.

\begin{savaistu}[Un Camerounais à la tête de la FMM]
Le Cameroun a fourni le premier commandant de la Force Mixte Multinationale réactivée : le général de brigade Bouba Dobekreo (2015-2016). Ce choix témoigne de la confiance des partenaires régionaux dans l'armée camerounaise et de la place centrale du Cameroun dans la sécurité du bassin du Lac Tchad.
\end{savaistu}

\begin{popfigure}
\begin{center}
\begin{tikzpicture}[scale=1.0, every node/.style={font=\small}]
% Fond : zone Cameroun schématisée
\fill[popcream] (0,0) rectangle (9,6);
\draw[popdark, thick] (0,0) rectangle (9,6);
% Pays voisins (étiquettes)
\node[popmuted] at (1.2,6.35) {\textbf{TCHAD}};
\node[popmuted] at (8.6,6.35) {\textbf{RCA}};
\node[popmuted] at (-0.6,4.2) {\textbf{NIGERIA}};
\node[popmuted] at (4.5,-0.45) {\textbf{Gabon / Congo / G. équatoriale}};
% Zone Secteur 1 FMM (Extrême-Nord)
\fill[poporange!45] (2.2,4.4) -- (6.8,4.4) -- (6.8,6) -- (2.2,6) -- cycle;
\draw[poporange, very thick] (2.2,4.4) rectangle (6.8,6);
\node[popdark, align=center] at (4.5,5.55) {\textbf{Secteur 1 FMM}\\ (Extrême-Nord)};
% Base Mora
\fill[popink] (3.4,4.9) circle (0.09);
\node[popink, anchor=west] at (3.55,4.9) {Mora (QG secteur 1)};
% Lac Tchad
\fill[popblue!50] (6.0,5.6) ellipse (0.7 and 0.35);
\node[popblue, font=\footnotesize] at (6.0,5.6) {Lac Tchad};
% Front NOSO
\fill[poppink!55] (0,2.2) -- (2.2,2.2) -- (2.2,4.0) -- (0,4.0) -- cycle;
\draw[poppink, very thick] (0,2.2) rectangle (2.2,4.0);
\node[popdark, align=center] at (1.1,3.1) {\textbf{Front NOSO}\\ (Nord-Ouest /\\ Sud-Ouest)};
% Zone Bakassi
\fill[popgreen!45] (0,0.9) -- (1.6,0.9) -- (1.6,2.2) -- (0,2.2) -- cycle;
\draw[popgreen, very thick] (0,0.9) rectangle (1.6,2.2);
\node[popdark, align=center, font=\footnotesize] at (0.8,1.55) {\textbf{Zone}\\ \textbf{Bakassi}};
% Front Est (RCA / 3R)
\fill[poppurple!35] (6.8,1.2) -- (9,1.2) -- (9,4.4) -- (6.8,4.4) -- cycle;
\draw[poppurple, very thick] (6.8,1.2) rectangle (9,4.4);
\node[popdark, align=center] at (7.9,2.8) {\textbf{Front Est}\\ (incursions\\ 3R / RCA)};
% Yaoundé
\fill[popink] (4.2,1.4) circle (0.1);
\node[popink, anchor=west] at (4.35,1.4) {Yaoundé (hub diplomatique)};
% Flèches de mouvements terroristes
\draw[-{Stealth[length=3mm]}, poppink, very thick] (7.0,5.9) .. controls (5.6,5.4) .. (4.6,5.0);
\draw[-{Stealth[length=3mm]}, poppink, very thick] (0.2,5.2) .. controls (1.4,5.0) .. (2.6,4.8);
\draw[-{Stealth[length=3mm]}, poppurple, very thick] (9.2,3.4) -- (8.3,3.0);
% Nord
\draw[-{Stealth[length=3mm]}, popdark, thick] (8.6,0.3) -- (8.6,0.9);
\node[popdark, anchor=west] at (8.7,0.6) {N};
\end{tikzpicture}
\end{center}
\legende{Carte schématique des théâtres d'opérations et zones d'influence sécuritaire du Cameroun : secteur 1 de la FMM (orange), front NOSO (rose), front Est face aux incursions des groupes armés centrafricains (violet), zone du Bakassi (vert). Les flèches indiquent les axes d'incursion des groupes armés. Croquis simplifié, sans prétention à l'exactitude cartographique.}
\end{popfigure}

\courssub{La lutte anti-terroriste intérieure et la protection des populations}

\textbf{L'opération « Épervier ».} Dans l'Extrême-Nord, l'armée camerounaise conduit l'opération \emph{Épervier} (lancée en 2014) : sécurisation des localités frontalières, patrouilles conjointes avec les comités de vigilance, protection des marchés, des écoles et des axes routiers, frappes ciblées contre les bases de Boko Haram/ISWAP. Cette opération nationale s'inscrit dans le dispositif régional de la FMM (secteur 1) et bénéficie d'appuis en renseignement et en formation de partenaires bilatéraux (France, USA, Israël).

\textbf{Réfugiés et personnes déplacées internes (PDI).} Les conflits voisins (RCA, Nigeria) et le terrorisme ont provoqué d'importants mouvements de population : selon le HCR et l'OIM (données 2023-2024), le Cameroun accueille plus de 480 000 réfugiés (majoritairement Centrafricains à l'Est et Nigérians à l'Extrême-Nord) et compte environ 1 million de déplacés internes (Extrême-Nord, Nord-Ouest, Sud-Ouest). Le Cameroun applique le \cle{droit international humanitaire} (Conventions de Genève de 1949 et Protocoles additionnels de 1977, Convention de l'OUA de 1969 régissant les aspects propres aux problèmes des réfugiés en Afrique) et sa \emph{loi n° 2005/006 du 27 juillet 2005} portant statut des réfugiés, en coopération avec le HCR, l'OIM et les ONG humanitaires.

\begin{attention}[Droit d'asile et rapatriement « volontaire » : un dilemme]
Les accords tripartites Cameroun-RCA-HCR organisent le rapatriement « volontaire » des réfugiés centrafricains. Or, si la sécurité n'est pas réellement rétablie dans les zones de retour, ce rapatriement peut devenir un retour forcé déguisé, contraire au principe de \emph{non-refoulement} (article 33 de la Convention de Genève de 1951, article II-3 de la Convention de l'OUA de 1969). Le dilemme est à la fois \textbf{éthique} (protéger la dignité des réfugiés), \textbf{humanitaire} (ne pas renvoyer des personnes vers le danger) et \textbf{sécuritaire} (risque d'infiltration de combattants parmi les réfugiés, pression sur les ressources locales). En examen, montrez toujours les deux faces du problème et citez les instruments juridiques.
\end{attention}

\courssub{La diplomatie préventive : Yaoundé, hub diplomatique}

La \cle{diplomatie préventive} consiste à agir \emph{avant} qu'un conflit n'éclate ou ne reparte : bons offices, pourparlers, facilitation, déploiement d'observateurs. Yaoundé s'est imposée comme une capitale de la médiation en Afrique centrale : accords et pourparlers sur la RCA (2008, 2013, 2019), siège de la Commission mixte Cameroun-Nigeria (depuis 2002), réunions régulières du COPAX et de la CEEAC. En septembre-octobre 2019, le \emph{Grand Dialogue National} (« Major National Dialogue ») a réuni à Yaoundé des délégués pour chercher une issue à la crise anglophone ; les pourparlers ultérieurs avec les groupes séparatistes armés ont toutefois avorté, montrant les limites de la médiation interne lorsque la confiance fait défaut et que les positions sont irréconciliables.

\begin{popfigure}
\begin{center}
\begin{tikzpicture}[scale=1.0, every node/.style={font=\footnotesize}]
% Axe chronologique
\draw[-{Stealth[length=3mm]}, popdark, very thick] (0,0) -- (14.2,0);
\foreach \x/\y in {0.4/1975, 2.2/1979, 4.0/2002, 5.6/2006, 7.2/2008, 8.8/2013, 10.4/2015, 12.0/2019, 13.6/2024}{
  \draw[popdark] (\x,0.12) -- (\x,-0.12);
  \node[popdark, below] at (\x,-0.15) {\y};
}
% Événements alternés haut/bas
\draw[popblue, thick] (0.4,0.12) -- (0.4,0.9);
\node[popblue, above, align=center, text width=2.6cm] at (0.4,0.9) {1975 : premières médiations régionales};
\draw[poporange, thick] (2.2,0.12) -- (2.2,0.9);
\node[poporange, above, align=center, text width=2.6cm] at (2.2,0.9) {1979 : rôle dans la crise tchadienne};
\draw[popgreen, thick] (4.0,-0.12) -- (4.0,-1.0);
\node[popgreen, below, align=center, text width=2.8cm] at (4.0,-1.0) {2002 : arrêt CIJ Bakassi ; FOMUC en RCA};
\draw[poppurple, thick] (5.6,0.12) -- (5.6,0.9);
\node[poppurple, above, align=center, text width=2.6cm] at (5.6,0.9) {2006 : Accords de Greentree};
\draw[popblue, thick] (7.2,-0.12) -- (7.2,-1.0);
\node[popblue, below, align=center, text width=2.8cm] at (7.2,-1.0) {2008 : Accords de Libreville (RCA) ; MICOPAX};
\draw[poporange, thick] (8.8,0.12) -- (8.8,0.9);
\node[poporange, above, align=center, text width=2.6cm] at (8.8,0.9) {2013 : cessez-le-feu de Libreville ; transfert du Bakassi};
\draw[poppink, thick] (10.4,-0.12) -- (10.4,-1.0);
\node[poppink, below, align=center, text width=2.8cm] at (10.4,-1.0) {2015 : FMM réactivée ; opération Épervier};
\draw[popgold, thick] (12.0,0.12) -- (12.0,0.9);
\node[popgold, above, align=center, text width=2.6cm] at (12.0,0.9) {2019 : Grand Dialogue National ; Accord Khartoum};
\draw[popdark, thick] (13.6,-0.12) -- (13.6,-1.0);
\node[popdark, below, align=center, text width=2.6cm] at (13.6,-1.0) {2024 : lutte continue Boko Haram/ISWAP ; NOSO};
\end{tikzpicture}
\end{center}
\legende{Chronologie simplifiée des grandes médiations et interventions militaires camerounaises (1975-2024).}
\end{popfigure}

\courssub{Analyser un accord de cessez-le-feu : méthode et application}

\begin{methode}[Analyser un accord de cessez-le-feu ou de paix]
Face à un document diplomatique (extrait d'accord, article de presse, rapport d'organisation internationale), procédez en quatre étapes :
\begin{enumerate}
\item \textbf{Identifier les parties prenantes} : qui signe ? (gouvernement, groupes armés, opposition politique, représentants de la société civile).
\item \textbf{Identifier les garants et facilitateurs} : qui supervise et parraine ? (État médiateur, UA, ONU, CEEAC, communauté internationale).
\item \textbf{Repérer le mécanisme de vérification et de suivi} : commission de suivi mixte, observateurs militaires/civils, sanctions prévues en cas de violation, calendrier de mise en œuvre.
\item \textbf{Relever les dispositions DDR/RSS} : \cle{Désarmement} (remise des armes), \cle{Démobilisation} (dissolution des unités combattantes), \cle{Réinsertion} (retour des ex-combattants à la vie civile : formation, emploi, appui psychosocial), \cle{Réforme du secteur de sécurité} (RSS : intégration partielle dans l'armée/police, vetting).
\end{enumerate}
Concluez en évaluant la \emph{crédibilité} et la \emph{durabilité} de l'accord : un accord sans mécanisme de vérification indépendant, sans DDR/RSS financé et sans volonté politique réelle des parties reste souvent lettre morte.
\end{methode}

\begin{exoresolu}[Analyse de l'Accord de Khartoum (RCA, 2019)]
\textbf{Énoncé.} En février 2019, l'Accord politique pour la paix et la réconciliation en République centrafricaine est signé à Khartoum (Soudan) entre le gouvernement centrafricain et 14 groupes armés, sous l'égide de l'Union Africaine, avec la facilitation de la CEEAC. Le Cameroun, voisin direct, soutient le processus. À l'aide de la méthode en quatre étapes, analysez cet accord et expliquez pourquoi sa mise en œuvre reste fragile.
\tcblower
\textbf{Solution détaillée.}
\begin{enumerate}
\item \emph{Parties prenantes} : le gouvernement de la RCA (président Faustin-Archange Touadéra) d'une part, 14 groupes armés (ex-Séléka, anti-Balaka, 3R, UPC, MPC, FPRC, etc.) d'autre part.
\item \emph{Garants et facilitateurs} : l'Union Africaine (garant principal), la CEEAC (facilitatrice), la MINUSCA (mission onusienne d'appui) ; les États voisins (Cameroun, Tchad, RDC, Congo, Soudan) apportent un appui diplomatique et logistique.
\item \emph{Mécanisme de vérification} : un comité de suivi national (CNS) et un mécanisme de vérification régional appuyés par la MINUSCA, chargés de constater les violations et de proposer des sanctions.
\item \emph{Dispositions DDR/RSS} : l'accord prévoit le désarmement, la démobilisation, la réinsertion (DDR) des combattants, et l'intégration d'une partie d'entre eux dans des unités mixtes de sécurité (USMS) dans le cadre de la RSS.
\end{enumerate}
\emph{Évaluation critique} : l'accord reste fragile car plusieurs groupes armés continuent de violer le cessez-le-feu, le DDR/RSS avance lentement (manque de financements pérennes, absence de cantonnement effectif, méfiance mutuelle) et l'État centrafricain ne contrôle pas tout son territoire (présence de Wagner/AFRICA CORPS depuis 2018 complexifie le paysage sécuritaire). Pour le Cameroun, l'enjeu est direct : l'instabilité centrafricaine entretient les incursions des groupes armés (notamment les « 3R » - Retour, Réclamation, Réhabilitation) dans les régions de l'Est et de l'Adamaoua et maintient le flux de réfugiés. Cela illustre la \emph{solidarité africaine} : la paix chez le voisin est une condition nécessaire de la sécurité nationale.
\end{exoresolu}

\courssub{Les coûts de l'engagement sécuritaire du Cameroun}

La contribution à la paix a un prix, qu'il faut savoir évaluer honnêtement pour mesurer l'effort national et ses arbitrages budgétaires :

\begin{itemize}
\item \textbf{Coûts budgétaires directs} : le budget du Ministère de la Défense (MINDEF) représente environ 300 à 350 milliards de FCFA par an (lois de finances 2020-2024), soit un effort de guerre estimé entre 5 et 10 \% du budget général de l'État. Ces ressources sont autant de moyens en moins pour l'éducation, la santé, les infrastructures ou l'agriculture (coût d'opportunité).
\item \textbf{Coûts humains} : pertes militaires (plusieurs milliers de soldats tués depuis 2014) et civiles (populations, humanitaires, enseignants), environ 1 million de personnes déplacées internes (PDI), plus de 480 000 réfugiés accueillis (charge logistique, sanitaire, scolaire, foncière).
\item \textbf{Coûts socio-économiques} : tensions sur la cohésion nationale (stigmatisation de certaines communautés), écoles fermées pendant des années (générations sacrifiées), risque de radicalisation des jeunes désœuvrés dans les zones de conflit, traumatismes psychologiques durables, effondrement de l'économie locale (agriculture, commerce transfrontalier, tourisme).
\item \textbf{Coûts diplomatiques} : exposition aux critiques internationales (droits de l'homme, droit humanitaire), conditionnalité de certains appuis budgétaires ou militaires, nécessité de rendre des comptes auprès des partenaires (UE, USA, UA, ONU).
\end{itemize}

\begin{exemplebox}[Le coût chiffré de la crise anglophone (NOSO)]
Selon des estimations de la Banque mondiale, du FMI et d'instituts de recherche (2020-2023), le coût économique cumulé de la crise du Nord-Ouest et du Sud-Ouest dépasse 1 000 milliards de FCFA : infrastructures détruites (écoles, centres de santé, ponts, marchés, installations CDC, PAMOL, Pamol Plantations), pertes de production agricole (cacao, café, banane, huile de palme) et industrielle, baisse du PIB régional estimée à plus de 30\%, dépenses militaires supplémentaires, prise en charge des PDI et des réfugiés au Nigeria. Ce montant représente plusieurs années d'investissement public dans ces deux régions : la paix est donc aussi un \emph{investissement économique} rentable.
\end{exemplebox}

\begin{attention}[La crise anglophone : un conflit interne qui internationalise la diplomatie]
La crise du NOSO est juridiquement un conflit \emph{interne} (non-international au sens du DIH). Pourtant, elle a provoqué des pressions internationales : résolutions du Parlement européen (2019, 2021), auditions au Congrès américain, préoccupations exprimées par l'UA et le Commonwealth (dont le Cameroun est membre depuis 1995), suspension partielle de l'éligibilité à l'AGOA (African Growth and Opportunity Act) par les USA en 2019. Résultat : la doctrine de \emph{non-ingérence}, pilier de la diplomatie camerounaise, se trouve mise à l'épreuve — le Cameroun invoque la non-ingérence pour refuser les ingérences extérieures, tout en étant sommé de rendre des comptes sur le respect des droits humains. Au Bac, évitez de présenter le NOSO comme un conflit international : c'est un conflit interne \emph{à dimension internationale} (répercussions transfrontalières, diaspora, pression diplomatique).
\end{attention}

\begin{exercice}[Pour s'entraîner]
1. Définissez la Force Mixte Multinationale (FMM), citez ses États membres, son quartier général et le secteur commandé par le Cameroun. 2. Expliquez en quoi le règlement du conflit du Bakassi illustre la doctrine camerounaise du règlement pacifique des différends (citez les étapes juridiques et diplomatiques). 3. En vous appuyant sur la méthode d'analyse d'un accord de paix (4 étapes), rédigez un paragraphe argumenté (12 à 15 lignes) répondant à la question : « La contribution du Cameroun à la gestion des conflits en Afrique renforce-t-elle sa sécurité nationale ? »
\end{exercice}

\begin{corrige}[Exercice n° 1]
\textbf{1.} La FMM est une force militaire régionale créée en 1994 (réactivée/renforcée en 2015) sous mandat de l'UA (appuyée par la résolution 2349 de l'ONU en 2017) pour combattre Boko Haram/ISWAP dans le bassin du Lac Tchad. Membres : Nigeria, Niger, Tchad, Cameroun, Bénin. Quartier général : N'Djamena (Tchad) ; le Cameroun commande le secteur 1 (basé à Mora, Extrême-Nord).
\textbf{2.} Face au différend frontalier du Bakassi, le Cameroun a choisi le droit (saisine de la CIJ en 1994, arrêt du 10 octobre 2002) puis la négociation diplomatique (Accords de Greentree, 12 juin 2006 ; mise en place de la Commission mixte Cameroun-Nigeria à Yaoundé sous égide ONU) plutôt que la guerre, aboutissant à un transfert de souveraineté pacifique et ordonné le 14 août 2013 : c'est l'application concrète et aboutie du principe de règlement pacifique des différends (Article 33 de la Charte de l'ONU, Article 4 de l'Acte constitutif de l'UA).
\textbf{3.} Éléments attendus pour la dissertation : \textbf{Thèse} : oui, la contribution régionale renforce la sécurité nationale car la sécurité nationale dépend de la stabilité régionale (effet domino : conflits voisins $\rightarrow$ réfugiés, trafics d'armes, incursions groupes armés 3R à l'Est) ; la FMM a affaibli Boko Haram/ISWAP dans le bassin du Lac Tchad ; les médiations (Bakassi, RCA, Tchad) renforcent le prestige diplomatique, le leadership régional et la légitimité internationale du Cameroun. \textbf{Antithèse/Nuance} : coûts budgétaires et humains très élevés (effort de guerre 5-10\% budget), conflit interne du NOSO non résolu malgré le Grand Dialogue, dépendance aux partenaires pour l'équipement/renseignement/formation, risque d'enlisement (terrorisme résiduel, banditisme), pression diplomatique sur les droits humains. \textbf{Synthèse} : la contribution régionale est une condition \emph{nécessaire} mais non \emph{suffisante} de la sécurité nationale ; elle doit s'accompagner d'une gouvernance interne inclusive, d'une justice transitionnelle efficace au NOSO et d'une diversification des partenariats pour réduire la vulnérabilité stratégique.
\end{corrige}

\begin{aretenir}[L'essentiel de la section]
Le Cameroun contribue à la gestion des conflits par la \emph{médiation} (Tchad, RCA, Bakassi, Guinée équatoriale), par l'\emph{engagement militaire régional} (FOMUC/MICOPAX/MISCA/MINUSCA en RCA, FMM contre Boko Haram/ISWAP), par la \emph{diplomatie préventive} (Yaoundé, hub diplomatique, siège CMCN, COPAX) et par la \emph{protection des populations} (réfugiés, PDI, droit international humanitaire, loi 2005/006). Cette action, guidée par les quatre piliers de sa doctrine (non-ingérence, règlement pacifique, solidarité africaine, intégrité territoriale), a des succès reconnus (Bakassi, affaiblissement Boko Haram, leadership régional) mais aussi des coûts lourds (budgétaires, humains, socio-économiques, diplomatiques) et des limites persistantes (crise anglophone, terrorisme résiduel, instabilité centrafricaine). La paix régionale et la sécurité nationale sont indissociables : l'une conditionne l'autre.
\end{aretenir}

\coursec{Problèmes et défis des relations internationales du Cameroun (Dossier 5)}

Après avoir analysé la contribution du Cameroun à la gestion des conflits et à la sécurité collective au sein de l'Union Africaine et du COPAX (section précédente), il convient d'interroger la contrepartie : quel est le prix à payer pour cet engagement diplomatique et militaire ? Si le Cameroun est un acteur reconnu de la stabilité régionale, sa marge de manœuvre internationale est structurellement contrainte par des fragilités économiques, des héritages historiques et des rapports de force géopolitiques nouveaux. Cette section dresse le bilan lucide des vulnérabilités qui limitent la souveraineté effective du pays sur la scène mondiale.

\courssub{L'endettement extérieur : le frein budgétaire à la souveraineté}

La dette publique est devenue la contrainte macroéconomique majeure qui conditionne la politique étrangère du Cameroun. Le service de la dette absorbe une part croissante des recettes budgétaires, réduisant l'espace fiscal nécessaire aux investissements sociaux (santé, éducation, infrastructures) et contraignant le pays à négocier sous contrainte avec ses créanciers.

\begin{definition}[Eurobonds]
Les \cle{Eurobonds} sont des obligations souveraines émises par l'État camerounais sur les marchés financiers internationaux (places de Londres, Luxembourg, New York), libellées en devises fortes (principalement l'euro et le dollar US). Contrairement aux prêts concessionnels des institutions de Bretton Woods (taux faibles, longs délais de remboursement, période de grâce), les Eurobonds portent des taux d'intérêt élevés (généralement entre 5\,\% et 9\,\% selon les émissions) et des maturités plus courtes (7 à 12 ans). Le Cameroun y a eu recours à partir de 2015 (première émission d'environ 750 millions de dollars US), puis en 2021 et 2024, pour refinancer la dette arrivant à échéance et financer le budget de l'État.
\end{definition}

\begin{propriete}[Le « paradoxe de la rente »]
L'exploitation pétrolière, amorcée à la fin des années 1970, a généré d'importantes recettes d'exportation, mais a structurellement \textbf{retardé la diversification économique} et \textbf{fragilisé la souveraineté budgétaire}. La facilité des « dollars du pétrole » a découragé l'industrialisation et la modernisation de l'agriculture d'exportation. Lorsque le contre-choc pétrolier de 1986 a frappé, l'État, privé d'une partie de sa rente, s'est endetté massivement pour maintenir les dépenses publiques, entrant dans le cycle des Plans d'Ajustement Structurel (PAS) imposés par le FMI et la Banque mondiale à la fin des années 1980. La rente a ainsi agi comme un « anesthésiant » qui a masqué l'absence de transformation structurelle de l'économie.
\end{propriete}

\begin{exemplebox}[Structure de la dette publique camerounaise (ordres de grandeur 2023-2024)]
La dette publique totale dépasse 12 000 milliards de FCFA, soit environ 45\,\% du PIB. Sa structure révèle une triple dépendance :
\begin{itemize}
    \item \textbf{Créanciers multilatéraux (FMI, Banque mondiale, BAD, BDEAC) :} environ 40\,\% à 45\,\%. Prêts concessionnels, mais conditionnés à des réformes structurelles.
    \item \textbf{Créanciers bilatéraux (Chine via Exim Bank, France/AFD, Japon/JICA, Turquie, Inde) :} environ 30\,\% à 35\,\%. La Chine est le premier créancier bilatéral (financement d'infrastructures : barrages, routes, ports, stades). Taux variables, prêts souvent liés à des marchés attribués à des entreprises chinoises.
    \item \textbf{Marchés financiers (Eurobonds) :} environ 20\,\% à 25\,\%. La part la plus coûteuse et la plus risquée (risque de change, risque de refinancement à échéance).
\end{itemize}
\end{exemplebox}

\begin{popfigure}
\begin{tikzpicture}
\begin{axis}[
    width=\linewidth,
    height=0.6\linewidth,
    grid=major,
    grid style={popline},
    xlabel={Année},
    ylabel={Pourcentage (\%)},
    legend pos=north west,
    legend cell align=left,
    xtick={2000,2005,2010,2015,2020,2024},
    xticklabel style={rotate=45, anchor=east},
    ymin=0, ymax=100,
    axis line style={popdark},
    tick style={popdark},
]
\addplot[popblue, thick, mark=*, mark size=2] coordinates {
    (2000, 15) (2005, 12) (2010, 8) (2015, 14) (2020, 22) (2024, 24)
};
\addlegendentry{Service dette / Recettes budgétaires (\%)}
\addplot[poppink, dashed, thick, no marks] coordinates {
    (2000, 22) (2024, 22)
};
\addlegendentry{Seuil d'alerte (22\%)}
\addplot[poporange, thick, mark=square*, mark size=2] coordinates {
    (2000, 88) (2005, 85) (2010, 82) (2015, 78) (2020, 80) (2024, 76)
};
\addlegendentry{Part exportations primaires / Exportations totales (\%)}
\end{axis}
\end{tikzpicture}
\legende{Évolution indicative du poids du service de la dette sur les recettes budgétaires (bleu) et de la part des produits primaires dans les exportations (orange) de 2000 à 2024. Sources : rapports FMI (Article IV), MINFI, INS. Le franchissement du seuil d'alerte coïncide avec la période des émissions d'Eurobonds. La dépendance aux matières premières reste structurellement élevée (plus de 75\,\%).}
\end{popfigure}

\begin{exoresolu}[Analyse de la soutenabilité de la dette]
\textbf{Énoncé :} En 2023, les recettes budgétaires de l'État du Cameroun s'élèvent à environ 4 500 milliards de FCFA. Le service de la dette (capital + intérêts) est estimé à 950 milliards de FCFA. On considère qu'un ratio service de la dette / recettes supérieur à 22\,\% signale un risque élevé de détresse de la dette. Calculer le ratio et interpréter la situation pour la marge de manœuvre diplomatique du pays.
\tcblower
\textbf{Solution détaillée :}
\[
\text{Ratio} = \frac{\text{Service de la dette}}{\text{Recettes budgétaires}} \times 100 = \frac{950}{4500} \times 100 \approx 21,1\,\%
\]
Ce ratio (21,1\,\%) se situe \textbf{juste en dessous du seuil critique de 22\,\%}, mais il est en hausse constante (il était d'environ 14\,\% en 2015).
\textbf{Conséquences diplomatiques :}
\begin{enumerate}
    \item \textbf{Conditionnalité du FMI :} pour maintenir l'accès aux financements concessionnels, le Cameroun doit accepter des « actions préalables » (réforme des subventions aux carburants, élargissement de l'assiette fiscale, maîtrise de la masse salariale). Cela réduit l'autonomie de décision budgétaire.
    \item \textbf{Risque de refinancement :} d'importants remboursements d'Eurobonds arrivent à échéance dans les années 2025-2027. Le pays doit émettre de nouvelle dette (souvent plus chère) ou puiser dans les réserves de change, ce qui affaiblit sa position de négociation vis-à-vis de ses créanciers (Chine, UE).
    \item \textbf{Arbitrage social :} chaque franc consacré au service de la dette est un franc en moins pour la santé et l'éducation, ce qui alimente le mécontentement social que les partenaires internationaux surveillent de près (droits humains, gouvernance).
\end{enumerate}
\end{exoresolu}

\begin{attention}[Souveraineté formelle et souveraineté effective]
Ne pas confondre la \textbf{souveraineté formelle} (siège à l'ONU, armée nationale, drapeau, hymne, droit de vote dans les instances internationales) et la \textbf{souveraineté effective} (capacité réelle à décider seul de son budget, de sa politique économique, de sa sécurité alimentaire et énergétique, de son modèle de développement sans diktat extérieur). Un pays très endetté perd une partie de sa souveraineté effective : ses choix budgétaires sont surveillés par le FMI, ses infrastructures stratégiques sont construites par des entreprises étrangères, sa sécurité sous-régionale dépend en partie de l'appui logistique de partenaires extérieurs. C'est le cœur du Dossier 5.
\end{attention}

\courssub{Dépendance structurelle et déséquilibre commercial chronique}

L'économie camerounaise reste une économie extravertie : elle exporte des produits primaires à faible valeur ajoutée et importe des biens manufacturés à forte valeur ajoutée. Ce piège structurel creuse le déficit commercial et maintient le pays en position de preneur de prix (\emph{price taker}) sur les marchés mondiaux.

\begin{center}
\begin{tabularx}{\linewidth}{|l|X|X|X|}
\hline
\rowcolor{popblueL}
\textbf{Flux} & \textbf{Produits principaux} & \textbf{Poids} & \textbf{Vulnérabilité} \\
\hline
Exportations & Pétrole brut, bois, cacao, café, coton, aluminium & Plus de 75\,\% des exportations & Volatilité des cours mondiaux, termes de l'échange défavorables, faible transformation locale \\
\hline
Importations & Produits pétroliers raffinés, machines et équipements, médicaments, céréales (blé, riz), véhicules & Part dominante des importations & Dépendance énergétique, dépendance alimentaire, sorties de devises \\
\hline
\end{tabularx}
\end{center}

\begin{methode}[Analyse SWOT de la position géopolitique du Cameroun]
Pour structurer une argumentation en dissertation ou à l'oral (épreuve d'ECM), utilisez la matrice SWOT (FFOM en français) :
\begin{enumerate}
    \item \textbf{Forces (internes) :} stabilité relative, position charnière entre l'Afrique centrale et l'Afrique de l'Ouest, forêt du bassin du Congo (atout pour la diplomatie climatique), armée expérimentée (engagement contre Boko Haram et dans le COPAX), double appartenance Francophonie / Commonwealth, ressources diversifiées (agriculture, mines, hydroélectricité).
    \item \textbf{Faiblesses (internes) :} endettement élevé, crise anglophone (Nord-Ouest / Sud-Ouest) non résolue, gouvernance critiquée (corruption perçue), infrastructures déficitaires (énergie, transport), jeunesse sous-employée, dépendance aux importations alimentaires et énergétiques.
    \item \textbf{Opportunités (externes) :} ZLECAf (marché africain de plus d'un milliard de consommateurs), transition énergétique mondiale (demande de minerais stratégiques, financement carbone), diversification des partenaires (Chine, Turquie, pays du Golfe, Inde, Brésil), revendication africaine d'une réforme du Conseil de sécurité de l'ONU.
    \item \textbf{Menaces (externes) :} rivalités entre grandes puissances en Afrique centrale et au Sahel, protectionnisme (règlement européen sur la déforestation, taxe carbone aux frontières), changement climatique (sécheresses dans le Nord, montée des eaux à Douala), terrorisme transfrontalier (Boko Haram / ISWAP, groupes armés centrafricains), durcissement des conditions de financement international (hausse des taux d'intérêt).
\end{enumerate}
\textbf{Synthèse :} le Cameroun est un « État pivot » fragile : ses atouts géographiques et démographiques sont neutralisés par ses faiblesses structurelles. Sa diplomatie doit transformer les opportunités (ZLECAf, climat) en leviers pour corriger les faiblesses (industrialisation, endettement) avant que les menaces ne deviennent existentielles.
\end{methode}

\courssub{Conditionnalités, ingérences et pressions extérieures}

L'aide au développement et les financements internationaux ne sont jamais neutres. Ils s'accompagnent de conditionnalités politiques et économiques qui pèsent sur la souveraineté.

\begin{itemize}
    \item \textbf{Conditionnalités économiques (FMI, Banque mondiale) :} les programmes d'appui imposent la réduction progressive des subventions aux carburants (hausse des prix à la pompe en 2023-2024), la maîtrise de la masse salariale et la réforme de la fiscalité. Ces mesures, techniquement justifiées, ont un coût social et politique élevé (baisse du pouvoir d'achat, tensions sociales) que l'État assume seul face à sa population.
    \item \textbf{Conditionnalités politiques et de gouvernance (UE, USA) :} l'Union européenne (Accord de partenariat post-Cotonou, dit Accord de Samoa, signé en 2023) et les États-Unis (AGOA) lient l'aide et l'accès préférentiel aux marchés au respect de l'État de droit, des droits humains et de la lutte contre la corruption.
    \item \textbf{La crise anglophone comme facteur de pression internationale :} depuis 2016, la gestion de la crise du Nord-Ouest et du Sud-Ouest a valu au Cameroun :
    \begin{itemize}
        \item des mesures restrictives américaines (restrictions de visas, suspension partielle de certaines coopérations militaires) ;
        \item des résolutions critiques du Parlement européen appelant au dialogue et surveillant l'usage des fonds européens ;
        \item une attention accrue des organes des Nations Unies (Haut-Commissariat aux droits de l'homme) sur la situation humanitaire et les droits humains.
    \end{itemize}
    \item \textbf{Contentieux frontaliers et perception de cession :} l'arrêt de la Cour internationale de Justice (2002) attribuant la presqu'île de Bakassi au Cameroun a été appliqué (retrait nigérian en 2006-2008, accord de Greentree). Si le Cameroun a juridiquement gagné ce contentieux, la gestion de la frontière avec le Nigeria (pêche dans le lac Tchad, sécurité face à Boko Haram) demeure un sujet sensible de coopération obligée.
\end{itemize}

\courssub{Concurrence des puissances : le risque d'instrumentalisation}

Le Cameroun est devenu un terrain de compétition entre puissances établies et puissances émergentes. Cette multipolarité offre des alternatives de financement mais crée des risques d'alignement forcé ou d'instrumentalisation.

\begin{itemize}
    \item \textbf{France et partenaires traditionnels (UE, USA) :} partenaires historiques (langue, droit, accords de coopération et de défense, formation des élites, premiers investisseurs). Pressions sur la gouvernance, les droits humains et le climat. Perception critique d'une partie de l'opinion (« Françafrique »).
    \item \textbf{Chine :} premier partenaire commercial et premier créancier bilatéral, constructeur d'infrastructures « clé en main » (barrages, routes, stades, bâtiments publics). Modèle fondé sur la non-ingérence politique. Risques : opacité de certains contrats, endettement peu transparent, faible transfert de technologie et d'emplois locaux.
    \item \textbf{Russie :} présence marquée en Centrafrique voisine (coopération militaire et sécuritaire). Le Cameroun ménage un équilibre : abstention lors de votes à l'ONU concernant la guerre en Ukraine, signature d'un accord de coopération militaire avec Moscou en 2022, tout en maintenant ses partenariats occidentaux.
    \item \textbf{États-Unis :} partenaire sécuritaire clé (formation, renseignement dans la lutte contre Boko Haram), marché ouvert par l'AGOA, mais pressions sur les droits humains et la crise anglophone.
    \item \textbf{Pays du Golfe, Turquie, Brésil, Inde :} nouveaux partenaires (investissements, agro-industrie, santé, défense). Diversification bienvenue, mais risque de fragmentation des partenariats en l'absence d'un cadre stratégique national cohérent.
\end{itemize}

\begin{savaistu}[Le Cameroun et la réforme du Conseil de sécurité de l'ONU]
Le Cameroun soutient la \textbf{position commune africaine} dite « Consensus d'Ezulwini » (adoptée en 2005 en Eswatini) portée par le Comité des dix chefs d'État de l'Union Africaine. La revendication africaine est claire : \textbf{deux sièges permanents avec droit de veto} et des sièges non permanents supplémentaires pour l'Afrique. Le Cameroun plaide régulièrement à la tribune de l'Assemblée générale des Nations Unies pour cette réforme, arguant que l'Afrique est la seule région sans siège permanent alors qu'elle fournit la majorité des dossiers traités par le Conseil de sécurité. Obtenir cette réforme serait le symbole de la reconnaissance du rôle de l'Afrique — et du Cameroun — dans la gouvernance mondiale.
\end{savaistu}

\courssub{Défis émergents : climat, numérique, migrations, réforme de la gouvernance mondiale}

Au-delà des problèmes classiques, de nouvelles frontières diplomatiques s'ouvrent, nécessitant des compétences techniques que l'administration peine parfois à mobiliser.

\begin{itemize}
    \item \textbf{Diplomatie climatique (bassin du Congo) :} le Cameroun abrite une partie de la deuxième forêt tropicale du monde. Enjeux : monétiser le carbone forestier (mécanisme REDD+, crédits carbone), accéder aux financements climatiques internationaux, tout en protégeant les droits des populations locales et des peuples autochtones (Baka, Bagyeli). Risque : accaparement des terres forestières au nom de la compensation carbone.
    \item \textbf{Diplomatie numérique :} cybersécurité (loi de 2010 sur la cybersécurité et la cybercriminalité, Convention de l'Union Africaine dite « Convention de Malabo » de 2014), souveraineté des données, régulation des réseaux sociaux (lutte contre les discours de haine et les fausses informations), fiscalité des grandes entreprises numériques. Le Cameroun doit choisir ses normes entre les modèles européen, chinois et africain.
    \item \textbf{Gestion des flux migratoires :} le Cameroun est à la fois pays de transit (migrants d'Afrique centrale et de l'Ouest vers le Maghreb et l'Europe), pays d'accueil (centaines de milliers de réfugiés centrafricains et nigérians) et pays d'émigration (diaspora qualifiée, transferts d'argent importants). La coopération avec l'UE (retours volontaires assistés, gestion des frontières) doit être conciliée avec la souveraineté nationale en matière d'asile.
    \item \textbf{Réforme du système financier international :} plaidoyer, aux côtés du Groupe des 77 et de l'Union Africaine, pour une meilleure représentation de l'Afrique dans les institutions de Bretton Woods, la réallocation des Droits de Tirage Spéciaux (DTS) du FMI et des mécanismes de conversion de la dette en investissements climatiques (\emph{debt-for-nature swaps}).
\end{itemize}

\begin{popfigure}
\begin{tikzpicture}[
    mindmap,
    grow cyclic,
    every node/.style=concept,
    concept color=popblue,
    level 1/.append style={level distance=4.5cm, sibling angle=90},
    level 2/.append style={level distance=3.2cm, sibling angle=45},
    text=white,
    font=\small
]
\node[concept] {Problèmes des RI du Cameroun}
    child[concept color=poporange] { node {Racines structurelles}
        child { node {Économie de rente (pétrole, bois)} }
        child { node {Héritage colonial (frontières, double tutelle)} }
        child { node {Modèle extraverti (export primaire / import manufacturé)} }
    }
    child[concept color=poppurple] { node {Vulnérabilités majeures}
        child { node {Endettement extérieur (Eurobonds, Chine, FMI)} }
        child { node {Déséquilibre commercial chronique} }
        child { node {Crises internes (NOSO, terrorisme, gouvernance)} }
    }
    child[concept color=popgreen] { node {Pertes de souveraineté}
        child { node {Conditionnalités FMI / UE} }
        child { node {Rivalités des puissances} }
        child { node {Pressions droits humains / sanctions} }
        child { node {Opacité de certains contrats} }
    }
    child[concept color=poppink] { node {Conséquences}
        child { node {Pauvreté et inégalités persistantes} }
        child { node {Fuite des cerveaux, jeunesse désœuvrée} }
        child { node {Poids diplomatique limité} }
        child { node {Risque de rupture sociale} }
    };
\end{tikzpicture}
\legende{Carte mentale des problèmes des relations internationales du Cameroun. Les racines (structure économique, histoire) nourrissent les vulnérabilités actuelles, qui entraînent des pertes de souveraineté effective et des conséquences sociales et diplomatiques. La diplomatie nationale doit agir à tous les niveaux.}
\end{popfigure}

\begin{aretenir}[Synthèse du Dossier 5 : les 5 défis majeurs]
\begin{enumerate}
    \item \textbf{Le piège de la dette :} un service de la dette supérieur à 20\,\% des recettes budgétaires entraîne une perte d'autonomie budgétaire, la conditionnalité du FMI et un risque de refinancement des Eurobonds.
    \item \textbf{Le piège de la rente :} des exportations primaires représentant plus de 75\,\% des ventes à l'étranger exposent le pays à la volatilité des cours, à des termes de l'échange défavorables et à la désindustrialisation.
    \item \textbf{Le piège de la conditionnalité :} l'aide et les financements sont liés à des réformes économiques (subventions, masse salariale) et à des exigences de gouvernance (droits humains, corruption, crise anglophone), ce qui crée une tension entre souveraineté formelle et souveraineté effective.
    \item \textbf{Le piège de la multipolarité :} la concurrence France / Chine / USA / Russie / Turquie offre des opportunités de diversification, mais comporte des risques d'instrumentalisation, de fragmentation des partenariats et d'opacité des contrats.
    \item \textbf{Les nouveaux fronts :} climat (forêt du bassin du Congo), numérique (données, cybersécurité), migrations (transit, diaspora, réfugiés), réforme de l'ONU (siège permanent pour l'Afrique). Ils exigent une expertise technique et une stratégie nationale cohérente.
\end{enumerate}
\end{aretenir}

\begin{exercice}[Sujet type Probatoire / Baccalauréat technique — ECM]
\textbf{Sujet :} « Le Cameroun, malgré son engagement actif dans le multilatéralisme (ONU, UA, CEMAC, Francophonie, Commonwealth), voit sa marge de manœuvre diplomatique réduite par des vulnérabilités structurelles. »
À partir de cette affirmation et de vos connaissances, analysez les principaux problèmes qui entravent l'efficacité de l'action extérieure du Cameroun et proposez des pistes pour renforcer sa souveraineté effective.
\textbf{Attendus :} introduction (définition de la souveraineté effective, problématique, annonce du plan) ; Partie 1 : vulnérabilités économiques (dette, dépendance commerciale) ; Partie 2 : vulnérabilités politiques et sécuritaires (crise anglophone, conditionnalités, concurrence des puissances) ; Partie 3 : pistes de résorption (ZLECAf, industrialisation, gestion de la dette, diplomatie climatique et numérique, réforme de l'ONU) ; conclusion.
\end{exercice}

\begin{corrige}[Exercice n°1 — Éléments de corrigé]
\textbf{Introduction :} définir la souveraineté effective (capacité réelle à décider et à agir par soi-même) par opposition à la souveraineté formelle. Problématique : dans quelle mesure l'engagement multilatéral du Cameroun est-il limité par ses vulnérabilités internes ? Thèse : l'engagement multilatéral est réel (casques bleus, médiations, COPAX), mais l'autonomie stratégique est entravée par l'endettement, l'extraversion économique et les crises internes.
\textbf{I. L'étau économique : dette et dépendance.}
\begin{itemize}
    \item Dette publique supérieure à 45\,\% du PIB ; service de la dette voisin de 21\,\% des recettes (seuil d'alerte : 22\,\%). Eurobonds coûteux et risqués. Conditionnalité du FMI (subventions, masse salariale) assimilable à une tutelle budgétaire.
    \item Structure des exportations : pétrole, bois, cacao, café, coton, aluminium (plus de 75\,\%). Importations : produits raffinés, machines, médicaments, céréales. Faible industrialisation, ZLECAf encore sous-exploitée.
\end{itemize}
\textbf{II. L'étau politico-sécuritaire et géopolitique.}
\begin{itemize}
    \item Crise anglophone : pressions américaines et européennes, résolutions du Parlement européen, image internationale ternie, coût militaire et budgétaire.
    \item Concurrence des puissances : Chine (infrastructures, dette) contre Occident (gouvernance, droits humains) ; risque d'instrumentalisation ; équilibre précaire (abstention à l'ONU sur l'Ukraine).
    \item Contentieux frontaliers : Bakassi (arrêt CIJ 2002 appliqué), tensions récurrentes avec le Nigeria (lac Tchad, sécurité).
\end{itemize}
\textbf{III. Pistes pour reconquérir la souveraineté effective.}
\begin{itemize}
    \item \textbf{Économie :} transformation locale des matières premières (cacao, bois, coton, minerais), exploitation de la ZLECAf, gestion rigoureuse de la dette (restructuration, conversions dette-climat), mobilisation des recettes intérieures (numérisation fiscale, lutte contre la fraude).
    \item \textbf{Gouvernance :} résolution inclusive de la crise anglophone (dialogue, décentralisation effective), lutte contre la corruption, indépendance de la justice, afin de lever les conditionnalités liées aux droits humains.
    \item \textbf{Diplomatie :} diversification des partenariats « gagnant-gagnant » (Turquie, Brésil, Inde, pays du Golfe) sans alignement exclusif ; leadership climatique (forêt du bassin du Congo comme atout carbone) ; diplomatie numérique (souveraineté des données) ; plaidoyer pour la réforme du Conseil de sécurité (Consensus d'Ezulwini).
    \item \textbf{Sécurité :} développement d'une industrie de défense locale, mutualisation des efforts sous-régionaux (COPAX, Force multinationale mixte contre Boko Haram).
\end{itemize}
\textbf{Conclusion :} la souveraineté effective se conquiert par la transformation structurelle de l'économie et la consolidation de la gouvernance interne. Le multilatéralisme n'est pas une fin en soi mais un levier au service de l'intérêt national, à condition de disposer d'une stratégie claire et des moyens de sa mise en œuvre.
\end{corrige}

\newpage

\coursec{Activité d'intégration}

\courssub{Objectifs de l'activité}
\emph{Mobiliser l'ensemble des savoirs et savoir-faire de l'unité « Le Cameroun dans la coopération multilatérale » à travers une simulation de négociation internationale réaliste.}
Plus précisément, l'élève doit être capable de :
\begin{itemize}
    \item Identifier et caractériser les \cle{acteurs} (États, organisations internationales, institutions financières, société civile) et leurs \cle{intérêts} respectifs vis-à-vis du Cameroun.
    \item Exploiter des \cle{données chiffrées} (PIB, dette, APD, flux d'IDE, effectifs Casques bleus) pour étayer une argumentation.
    \item Appliquer les principes de l'\cle{efficacité de l'aide} (appropriation, alignement, harmonisation, gestion axée sur les résultats, reddition de comptes mutuelle) issus de la Déclaration de Paris (2005) et du Partenariat de Busan (2011).
    \item Négocier des \cle{compromis} entre conditionnalités de gouvernance (partenaires traditionnels) et non-ingérence / infrastructures (partenaires émergents).
    \item Rédiger des documents diplomatiques formels : \cle{Mémorandum de position} et \cle{Déclaration finale}.
\end{itemize}

\courssub{Situation}
\textbf{Contexte :} Nous sommes en janvier 2030. Le Cameroun arrive à l'échéance de sa \cle{Stratégie Nationale de Développement 2020-2030 (SND30)}. Le bilan est contrasté : le taux de pauvreté stagne autour de \num{37,5}\% (ENEC 2022), le ratio dette/PIB atteint \num{45,8}\% (FMI, 2023), et les besoins de financement pour l'industrialisation (objectif « Pays émergent » à l'horizon 2035) sont estimés à \SI{12000}{milliards} FCFA sur la période 2030-2035. Le Premier Ministre, Chef du Gouvernement, convoque une \textbf{Conférence des Partenaires pour le Développement du Cameroun 2030 (CPDC 2030)} à Yaoundé pour négocier un \emph{Nouveau Cadre de Partenariat Stratégique (NCPS)}.

\textbf{Données macroéconomiques de référence (2023-2024) :}
\begin{center}
\begin{tabularx}{\linewidth}{|l|X|X|}\hline
\textbf{Indicateur} & \textbf{Valeur} & \textbf{Source} \\ \hline
PIB nominal & \SI{27500}{milliards} FCFA & INS / FMI \\ \hline
Croissance réelle & 4,0\% (prévision 2024) & FMI \\ \hline
Dette publique brute & 45,8\% du PIB & Minfi / FMI \\ \hline
Service de la dette / Recettes & 42\% & FMI (DSA 2023) \\ \hline
APD nette reçue (2022) & \SI{650}{milliards} FCFA & OCDE/CAD \\ \hline
Part APD / PIB & 2,4\% & Calcul \\ \hline
IDE entrants (stock) & \SI{3200}{milliards} FCFA & CNUCED \\ \hline
Rang contributeur ONU (Casques bleus) & 14\textsuperscript{e} mondial (\num{1250} pers.) & DPKO ONU \\ \hline
Transferts diaspora (estimation 2023) & \SI{400}{milliards} FCFA & Banque Mondiale / BEAC \\ \hline
\end{tabularx}
\end{center}

\textbf{Les 6 Délégations et leurs « Lignes Rouges » (extraits des dossiers fournis) :}

\begin{enumerate}
    \item \textbf{Gouvernement Camerounais (Chef de file : PM)} : \textit{Position :} « Appropriation nationale » (SND30). \textit{Intérêts :} Budgets d'appui sectoriels (éducation, santé), traitement de la dette (Cadre Commun G20, FRD/FMI), infrastructures (barrage Nachtigal, autoroute Yaoundé-Douala, port Kribi II). \textit{Lignes rouges :} Refus de la conditionnalité politique \emph{ex ante} (droits de l'homme, décentralisation) comme prérequis au décaissement ; refus de la révision du Code des Investissements 2013 imposée.
    \item \textbf{Partenaires Traditionnels (UE, France, Japon, Allemagne)} : \textit{Position :} « Partenariat politique renforcé » (Accord de Samoa, post-Cotonou). \textit{Intérêts :} Gouvernance démocratique, État de droit, lutte contre la corruption (CONAC, COSUMAF), gestion durable forêts (APV-FLEGT), climat (NDC). \textit{Lignes rouges :} Pas de budget d'appui sans réformes PFM (Gestion Finances Publiques) vérifiées par le FMI ; conditionnalité droits humains (crise NOSO, libertés publiques) ; transparence contrats chinois.
    \item \textbf{Partenaires Émergents (Chine, Brésil, Turquie, Russie, Inde)} : \textit{Position :} « Coopération Sud-Sud, gagnant-gagnant, non-ingérence ». \textit{Intérêts :} Accès ressources (bois, minerais, pétrole), marchés travaux (BTP), débouchés agricoles (cacao, coton). \textit{Lignes rouges :} Refus de toute conditionnalité politique/gouvernance ; préférence pour prêts concessionnels/EXIM Bank vs dons ; clauses de main-d'œuvre et matériaux d'origine.
    \item \textbf{Institutions Financières Internationales (FMI, BM, BAD, BEI, BOAD)} : \textit{Position :} « Viabilité macroéconomique et transformation structurelle ». \textit{Intérêts :} Réformes structurelles (subventions carburant, recettes fiscales, secteur financier), climat des affaires (cadre B-READY), projets bankables. \textit{Lignes rouges :} Programme FMI (FEC/MEDC) comme \emph{anchor} pour tout appui budgétaire ; plafonnement emprunts non concessionnels ; audit dette contractée hors budget.
    \item \textbf{Société Civile / Diaspora (ONG locales, Syndicats, CACO, Diaspora USA/France/Allemagne)} : \textit{Position :} « Redevabilité, inclusion, droits humains ». \textit{Intérêts :} Suivi budget participatif, protection défenseurs DDH, engagement diaspora (transferts \SI{400}{milliards} FCFA/an), emploi jeunes. \textit{Lignes rouges :} Refus austérité sociale (subventions carburant) sans filets de protection ; exigence siège observateur dans Comité de pilotage NCPS ; transparence contrats miniers/pétroliers (ITIE).
    \item \textbf{Observateurs (UA, CEEAC/COPAX, ONU, OIF, Commonwealth)} : \textit{Position :} « Paix, sécurité, intégration régionale, valeurs partagées ». \textit{Intérêts :} Contribution Cameroun FMM (Bassin Lac Tchad) et FOMAC/COPAX (RCA), mise en œuvre ZLECAf, réforme secteur sécurité (RSS), bilinguisme effectif (OIF), valeurs Harare (Commonwealth). \textit{Lignes rouges :} Liaison aide sécurité/développement (nexus HDP) ; respect engagements climat (Accord Paris) ; résolution pacifique crise NOSO.
\end{enumerate}

\courssub{Consignes}
\textbf{Séance 1 — Préparation (2h) : Analyse stratégique et rédaction du Mémorandum}
\begin{enumerate}
    \item \textbf{Analyse documentaire :} Chaque délégation étudie son « Dossier Acteur » (3 pages : discours officiels, données OCDE/FMI, traités signés). Identifier les \cle{3 priorités absolues} et les \cle{2 concessions possibles}.
    \item \textbf{Cartographie des alliances :} Dresser un tableau « Alliés naturels / Adversaires potentiels / Neutres à convaincre » pour chacun des 5 points de l'ordre du jour.
    \item \textbf{Rédaction du Mémorandum de position (1 page max, police 11, interligne 1,15) :} Structure imposée :
    \begin{enumerate}
        \item Rappel des engagements antérieurs (ex: CPD 2016, Accord de Samoa, TICAD, FOCAC).
        \item Diagnostic partagé (accord/désaccord sur chiffres).
        \item Proposition concrète par point de l'OdJ (montants, conditionnalités, calendrier, instance de suivi).
        \item « Lignes rouges » non négociables.
    \end{enumerate}
\end{enumerate}

\textbf{Séance 2 — Négociation (3h) : Table ronde « CPDC 2030 »}
L'enseignant préside (Facilitateur ONU). Ordre du jour strict (30 min par point) :
\begin{enumerate}
    \item \textbf{Cadre macroéconomique \& Dette :} FMI présente scénario baseline/réformes. Gouvernement demande rééchelonnement (G20 Cadre Commun). Partenaires traditionnels lient à gouvernance. Chine propose swap dette/infrastructure.
    \item \textbf{Infrastructures prioritaires :} Port Kribi II, Autoroute Yaoundé-Douala, Énergie (solaire Nord, barrage Kikot). Japon/UE : appel d'offres international, normes ESG. Chine : clé en main, financement EXIM, rapidité. BAD : cofinancement.
    \item \textbf{Gouvernance, Droits Humains, Décentralisation :} UE/USA/Commonwealth : conditionnalité décaissements (indicateurs CONAC, élections locales, code pénal). Gouvernement : souveraineté, lois existantes suffisantes. Société civile : mise en œuvre effective, pas de textes seulement.
    \item \textbf{Sécurité régionale \& Nexus HDP :} UA/COPAX : contribution FMM/FOMAC (logistique, renseignement), RSS. ONU : protection civils, DDR. Gouvernement : équipement armée, formation, refus ingérence.
    \item \textbf{Climat, Forêts \& Biodiversité :} CAFI/UE : paiement performance (REDD+), APV-FLEGT. Gouvernement : paiement services écosystémiques, souveraineté foncière. Brésil/Indonésie : alliance pays forestiers tropicaux.
\end{enumerate}
\textit{Consigne :} Chaque intervention \textless 2 min. Noteur désigné par délégation. Recherche de consensus \emph{« à la japonaise »} (nemawashi) avant vote final à main levée.

\textbf{Séance 3 — Synthèse \& Évaluation (2h) : Déclaration de Yaoundé 2030}
\begin{enumerate}
    \item \textbf{Rédaction collective (45 min) :} Groupe de rédaction (1 rep. par délégation) produit la \textbf{Déclaration de Yaoundé 2030} (2 pages) : Préambule (reconnaissance acquis), 5 engagements chiffrés (ex: « Les partenaires s'engagent à mobiliser \SI{3500}{milliards} FCFA d'APD nouvelle sur 2030-2035 »), Mécanisme de suivi (Comité de pilotage paritaire semestriel), Clause de révision (2027).
    \item \textbf{Signature symbolique \& Débriefing (30 min) :} Lecture plénière, signatures. Enseignant anime le retour critique : « Qu'est-ce qui a bloqué ? Quelle donnée a fait basculer l'argument ? ».
    \item \textbf{Auto-évaluation par grille de compétences (45 min) :} Chaque élève remplit la grille individuelle (voir \emph{Ressources}).
\end{enumerate}

\courssub{Ressources à mobiliser}
\emph{Concepts-clés (rappel de cours) :}
\begin{itemize}
    \item \textbf{Typologie de l'aide :} Dons (grants) vs Prêts concessionnels (grant element \textgreater 25\%) vs Prêts non concessionnels. \textbf{Appui budgétaire} (général/secteur) vs \textbf{Appui projet}.
    \item \textbf{Principes d'efficacité (Paris 2005 / Busan 2011) :} \cle{Appropriation} (leadership pays), \cle{Alignement} (systèmes pays), \cle{Harmonisation} (procédures donateurs), \cle{Gestion axée résultats}, \cle{Reddition comptes mutuelle}.
    \item \textbf{Conditionnalités :} \emph{Ex ante} (prérequis) vs \emph{Ex post} (tranches). \textbf{Conditionnalité politique} (droits humains, démocratie) vs \textbf{Conditionnalité technique} (PFM, passation marchés).
    \item \textbf{Architecture onusienne :} CSNU (maintien paix), ECOSOC (développement), PNUD (coordination pays). \textbf{Contribution Cameroun :} \num{1250} casques bleus (MINUSCA, MONUSCO, ex-MINUSMA), 4\textsuperscript{e} contributeur africain.
    \item \textbf{Intégration régionale :} CEMAC (union monétaire, marché commun), CEEAC (sécurité via COPAX/FOMAC), ZLECAf (zone libre-échange continentale). \textbf{COPAX :} Mécanisme prévention/gestion/règlement conflits (État-major de planification, FOMAC).
    \item \textbf{Partenariats historiques :} France (Accords coopération 1960/1974, zone franc), Royaume-Uni (Commonwealth depuis 1995, anglophonie), Allemagne (coopération technique GIZ/KfW).
    \item \textbf{Nouveaux partenariats :} Chine (FOCAC, BRI / Nouvelles Routes Soie), Japon (TICAD, qualité infrastructure), Brésil (coopération technique agriculture/Embrapa), Turquie (investissements BTP/santé).
\end{itemize}

\courssub{Démarche et corrigé commenté}
\begin{exoresolu}[Guide pédagogique \& Éléments de corrigé attendus]
\textbf{1. Séance 1 — Attendus du Mémorandum (Exemple : Délégation Gouvernement)}
\begin{itemize}
    \item \textbf{Diagnostic :} « La SND30 a permis la résilience face au Covid-19 et à la guerre Ukraine (croissance 3,6\% moy. 2020-23). Cependant, le gap de financement infrastructurel est de \SI{4500}{milliards} FCFA/an. La dette, bien que viable (FMI), absorbe 42\% des recettes, limitant l'investissement social. »
    \item \textbf{Propositions OdJ 1 (Macro) :} Demande appui budgétaire \SI{500}{milliards} FCFA/an (2024-2026) sans conditionnalité politique \emph{ex ante}. Acceptation revue FMI technique (PFM) semestrielle. Demande traitement dette bilatérale Chine (Club de Paris + créanciers non-membres) via Cadre Commun G20 et adhésion FRD (FMI).
    \item \textbf{Propositions OdJ 2 (Infra) :} Port Kribi II : Partenariat Public-Privé (PPP) cofinancé BAD/BEI/EXIM Chine. Autoroute Ydé-Dla : Appel d'offres international (normes UE/Japon) mais financement mixte. Énergie : Solaire Nord (BAD/UE), Kikot (Chine).
    \item \textbf{Lignes rouges :} Aucune ingérence dans le calendrier électoral ou le statut des régions (décentralisation actée loi 2019). Code Investissements 2013 intangible (sécurité juridique).
\end{itemize}

\textbf{2. Séance 2 — Dynamique de négociation (Points de blocage \& Solutions types)}
\begin{center}
\begin{tabularx}{\linewidth}{|l|X|X|}\hline
\textbf{Point OdJ} & \textbf{Conflit principal} & \textbf{Compromis « Déclaration Yaoundé 2030 » (Exemple)} \\ \hline
1. Macro/Dette & FMI/UE : Réforme subventions carburant (coût \SI{1200}{mds} FCFA/an) vs Gov : Risque social explosion. & \textbf{Accord :} Suppression progressive (3 ans) + Filets protection (Transferts monétaires ciblés \SI{150}{mds}/an, Banque Mondiale). Rééchelonnement dette bilatérale Chine accepté (maturité +5 ans, taux -1\%). \\ \hline
2. Infrastructures & Japon/UE : Normes ESG, transparence, délais longs vs Chine : Rapidité, clauses préférence nationale, opacité coûts. & \textbf{Accord :} Cadre unique « Infrastructures Durables Cameroun » (normes ESG obligatoires \emph{ex post} audit BAD). Port Kribi II : Concession 30 ans consortium international (lead Chine + opérateur portuaire singapourien/néerlandais). \\ \hline
3. Gouvernance/DH & UE/USA : Conditionnalité décaissement (Libération détenus politiques, Code pénal, Élections locales) vs Gov : Souveraineté, lois votées. & \textbf{Accord :} « Dialogue politique structuré » trimestriel (Art. 8 Accord Samoa). Indicateurs gouvernance (CONAC, ITIE, Décentralisation) suivis par Comité paritaire. Pas de suspension automatique, mais « clause de révision » si régression avérée (ONU/OIF). \\ \hline
4. Sécurité & UA/ONU : Nexus HDP (Humanitaire-Développement-Paix), RSS vs Gov : Équipement offensif, refus ingérence RSS. & \textbf{Accord :} Appui formation/équipement non-létal (UE/USA/Japon) + Contribution FMM/FOMAC \num{500} hommes/an (financement UA/Partenaires). Création « Fonds Stabilisation Extrême-Nord/Nord-Ouest/Sud-Ouest » géré par MINAT/UNDP. \\ \hline
5. Climat/Forêts & CAFI/UE : Paiement performance REDD+ (résultats vérifiés) vs Gov : Paiement efforts (politiques), droits communautés. & \textbf{Accord :} « Partenariat Forêts-Climat 2030 ». Enveloppe \SI{150}{millions} USD (2024-2030). 60\% paiement résultats (déforestation évitée), 40\% mise en œuvre politiques (aménagement, tenure foncière, APV-FLEGT). Brésil/Indonésie : Appui technique surveillance satellite. \\ \hline
\end{tabularx}
\end{center}

\textbf{3. Séance 3 — Structure type « Déclaration de Yaoundé 2030 »}
\begin{verbatim}
DÉCLARATION DE YAOUNDÉ 2030
NOUVEAU CADRE DE PARTENARIAT STRATÉGIQUE CAMEROUN - PARTENAIRES

PRÉAMBULE
Les Parties prenantes... réaffirment leur attachement aux principes de la Charte des Nations Unies,
de l'Acte constitutif de l'UA, de l'Accord de Samoa (UE-ACP), du Partenariat de Busan...
Reconnaissent les progrès réalisés (SND30, émergence port Kribi, contribution paix ONU)...

ENGAGEMENTS QUANTIFIÉS (ÉCHÉANCE 2027)
1. FINANCEMENT : Mobilisation \SI{3500}{milliards} FCFA ressources nouvelles (APD + IDE + Diaspora Bonds)
   dont \SI{1200}{milliards} Appuis Budgétaires (Gouvernance), \SI{1800}{milliards} Infrastructures (PPP),
   \SI{500}{milliards} Climat/Forêts (CAFI/UE/BM).
2. DETTE : Maintien viabilité (ratio dette/PIB < 50\%). Rééchelonnement bilatéral Chine (\SI{1200}{milliards} FCFA).
   Adhésion facilité de résilience et durabilité (FRD) FMI (\SI{400}{millions} USD).
3. GOUVERNANCE : Opérationnalisation CONAC (poursuites effectives), ITIE (rapports annuels), Décentralisation
   (transfert compétences/ressources 2024-2025). Revue conjointe semestrielle.
4. SÉCURITÉ : Contribution FMM/FOMAC/COPAX maintenue. Appui RSS (formation, justice militaire, DDR).
   Fonds Stabilisation \SI{50}{milliards} FCFA/an (Gov + Partenaires).
5. CLIMAT : NDC révisée (-35\% émissions 2030). Zéro déforestation nette 2030. Mécanisme paiement
   services écosystémiques (PSE) national opérationnel 2025.

GOUVERNANCE DU PARTENARIAT
Comité de Pilotage Paritaire (CPP) : Co-présidence PM / Chef de file Partenaires (Présidence UE / Troïka).
Réunion semestrielle. Secrétariat technique permanent (Minepat + Ambassade Chef de file).
Évaluation à mi-parcours 2027 (Indicateurs résultats + Dépenses).

CLAUSE DE RÉVISION ET DE SORTIE
Révision possible sur demande écrite (6 mois préavis). Non-respect engagements majeurs (gouvernance,
dette, sécurité) déclenche clause de consultation (Art. 96 Cotonou/Samoa) avant suspension.

Fait à Yaoundé, le [Date], en 6 exemplaires (Français, Anglais, Chinois, Arabe).
Signatures : PM Cameroun, Ambassadeurs UE/France/Japon/Chine/Brésil, Représentants FMI/BM/BAD/ONU/UA/OIF/Commonwealth,
Présidents Plateformes Société Civile / Diaspora.
\end{verbatim}

\textbf{4. Grille d'évaluation des compétences (Sur 20 pts)}
\begin{center}
\begin{tabularx}{\linewidth}{|l|X|c|}\hline
\textbf{Compétence (Programme ECM)} & \textbf{Indicateurs observables} & \textbf{Barème} \\ \hline
\textbf{Mobilisation des savoirs} (Acteurs, Instruments, Données) & Citation précise chiffres (PIB, APD, Casques bleus), distinction types aide, connaissance accords (Samoa, FOCAC, TICAD, ZLECAf). & /5 \\ \hline
\textbf{Argumentation \& Rhétorique} & Cohérence position, usage concepts (appropriation, conditionnalité, nexus HDP), réponses aux objections, langue diplomatique. & /5 \\ \hline
\textbf{Esprit de compromis \& Négociation} & Concessions réalistes proposées, identification zones d'accord possibles (ZOPA), respect lignes rouges adverses, recherche consensus. & /5 \\ \hline
\textbf{Rédaction professionnelle} & Mémorandum : structure, concision, ton. Déclaration : forme juridique, précision engagements, mécanismes suivi. Orthographe/Syntaxe. & /5 \\ \hline
\textbf{TOTAL} & & /20 \\ \hline
\end{tabularx}
\end{center}
\end{exoresolu}

\courssub{Bilan de l'activité}
Cette simulation a permis de mettre en évidence la \textbf{complexité systémique} de la diplomatie camerounaise contemporaine. Les élèves ont constaté que :
\begin{enumerate}
    \item \textbf{La multipolarité est une réalité opérationnelle :} Le Cameroun ne négocie plus uniquement avec Paris/Bruxelles/Washington (Partenaires Traditionnels) mais doit arbitrer entre le modèle \emph{conditionnel/qualité} (UE/Japon/FMI) et le modèle \emph{rapide/infrastructure/non-ingérence} (Chine/Turquie/Brésil), tout en respectant les cadres contraignants de l'UA/CEMAC/ONU.
    \item \textbf{Les données chiffrées sont des armes de négociation :} La maîtrise du ratio dette/recettes (42\%), du stock d'IDE, des transferts diaspora (\SI{400}{milliards} FCFA) ou des effectifs Casques bleus (\num{1250}) change la donne : le Cameroun n'est pas un « assisté » mais un \cle{partenaire stratégique} (sécurité, forêts, marché, voix ONU).
    \item \textbf{L'appropriation nationale (SND30) est le levier central :} Sans document de stratégie nationale crédible, aligné sur les ODD et l'Agenda 2063, le pays subit les agendas donateurs. Le Mémorandum a forcé la délégation Gouvernement à hiérarchiser \emph{ses} priorités.
    \item \textbf{Le nexus Sécurité-Développement (HDP) s'impose :} L'insécurité (Boko Haram, NOSO, RCA) n'est plus un dossier « défense » mais conditionne l'accès aux financements climat (CAFI), aux appuis budgétaires (FMI/UE) et à la notation souveraine.
    \item \textbf{La société civile et la diaspora sont devenues des acteurs incontournables :} Leur exclusion de la table des négociations (pratique passée) affaiblit la légitimité et le suivi (reddition de comptes mutuelle). Leur intégration dans le Comité de Pilotage Paritaire (CPP) est une innovation majeure de la Déclaration 2030.
\end{enumerate}
\textbf{Compétence transversale validée :} \emph{« Rédiger et justifier une démarche, une réponse ou une conclusion »} appliquée à un contexte professionnel simulé (diplomatie, gestion de projet, relations internationales), préfigurant les épreuves du Probatoire/Baccalauréat Technique (Étude de cas, Dissertation, Commentaire de documents).

\newpage

\coursec{Méthodes et exercices résolus}

\courssub{Méthode 1 : Analyser la participation du Cameroun aux opérations de maintien de la paix de l'ONU}
\begin{methode}[Analyse d'une contribution onusienne]
    \textbf{Objectif :} Expliquer l'engagement du Cameroun dans le maintien de la paix international à partir d'une mission précise.
    \begin{enumerate}
        \item \textbf{Identifier la mission :} Nom (ex. MINUSCA, MONUSCO, MINUSMA), date de déploiement, mandat (Chapitre VI ou VII de la Charte de l'ONU).
        \item \textbf{Préciser la contribution camerounaise :} Effectifs (policiers, militaires, observateurs), matériel, durée, pertes humaines éventuelles.
        \item \textbf{Relier aux principes de l'ONU :} Consentement des parties, impartialité, non-recours à la force sauf légitime défense ou mandat robuste (Chapitre VII).
        \item \textbf{Évaluer l'impact :} Pour le pays hôte (stabilisation, protection des civils) et pour le Cameroun (image, expérience, retombées financières et de formation).
        \item \textbf{Conclure :} Situer cet engagement dans la politique étrangère du Cameroun (respect du droit international, panafricanisme).
    \end{enumerate}
    \textbf{Réflexe Bac :} Citer obligatoirement le \cle{Chapitre VII} si le mandat est robuste et le \cle{numéro de la résolution} du Conseil de Sécurité si connu.
\end{methode}

\begin{exoresolu}[La MINUSCA et l'apport camerounais]
    \textbf{Énoncé :} En vous appuyant sur la Mission multidimensionnelle intégrée des Nations Unies pour la stabilisation en Centrafrique (MINUSCA), montrez comment le Cameroun met en œuvre son engagement pour la paix et la sécurité internationales.
    \tcblower
    \textbf{Solution détaillée :}
    \begin{enumerate}
        \item \textbf{Contexte et mandat :} La MINUSCA est créée par la \textbf{Résolution 2149 (2014)} du Conseil de Sécurité sous \textbf{Chapitre VII}. Son mandat : protection des civils, appui à la transition politique, désarmement (DDRR), réforme du secteur de la sécurité (RSS).
        \item \textbf{Contribution camerounaise :} Le Cameroun est un contributeur important de troupes et de police.
            \begin{itemize}
                \item \textbf{Police :} Unités de police constituées (FPU) pour la sécurisation de Bangui et les patrouilles.
                \item \textbf{Militaires :} Contingents d'infanterie pour la protection des civils et la lutte contre les groupes armés.
                \item \textbf{Personnel :} Officiers d'état-major, observateurs militaires, experts de mission.
            \end{itemize}
        \item \textbf{Respect des principes :} Les troupes camerounaises opèrent sous commandement onusien, dans le respect de l'impartialité (protection de toutes les communautés) et du droit international humanitaire (DIH).
        \item \textbf{Impact et retombées :}
            \begin{itemize}
                \item \emph{Pour la RCA :} Sécurisation des axes routiers, protection des sites de déplacés, appui à l'organisation des scrutins (2015-2016, 2020-2021).
                \item \emph{Pour le Cameroun :} Renforcement du professionnalisme de l'armée (expérience du terrain, interopérabilité), rayonnement diplomatique, indemnités versées par l'ONU aux contingents.
            \end{itemize}
        \item \textbf{Conclusion :} La participation à la MINUSCA illustre la \cle{diplomatie de présence} du Cameroun et son attachement au \cle{multilatéralisme} et à la \cle{sécurité collective} (Articles 1 et 24 de la Charte de l'ONU), tout en défendant ses intérêts de sécurité nationale (stabilité de la frontière Est).
    \end{enumerate}
    \textbf{Phrase de conclusion :} Le Cameroun, par son engagement constant dans la MINUSCA, affirme son rôle d'acteur de la paix en Afrique centrale et son adhésion aux principes de la Charte des Nations Unies.
\end{exoresolu}

\courssub{Méthode 2 : Expliquer le fonctionnement d'un mécanisme de sécurité régionale (COPAX / CEEAC)}
\begin{methode}[Analyse d'une architecture de paix régionale]
    \textbf{Objectif :} Décrire et évaluer un organe de prévention et de gestion des conflits en Afrique centrale.
    \begin{enumerate}
        \item \textbf{Situer l'institution :} Nom complet, traité fondateur (Traité de la CEEAC, 1983 ; Protocole relatif au COPAX, 2000), siège de la CEEAC (Libreville), lien avec l'UA (le Conseil de Paix et de Sécurité de l'UA est le niveau continental, le COPAX le niveau régional).
        \item \textbf{Identifier les organes clés :}
            \begin{itemize}
                \item \textbf{Conférence des Chefs d'État et de Gouvernement :} Décisions politiques suprêmes.
                \item \textbf{Conseil de Paix et de Sécurité de l'Afrique centrale (COPAX) :} Organe de décision et de suivi.
                \item \textbf{Comité d'État-Major :} Planification militaire.
                \item \textbf{Force Multinationale de l'Afrique centrale (FOMAC) :} Outil militaire de réaction.
                \item \textbf{Mécanisme d'Alerte Rapide de l'Afrique centrale (MARAC) :} Analyse des risques (siège : Libreville).
            \end{itemize}
        \item \textbf{Décrire le processus d'intervention :} Alerte (MARAC) $\rightarrow$ Analyse (Comité d'État-Major) $\rightarrow$ Décision politique (COPAX / Conférence) $\rightarrow$ Déploiement (FOMAC) $\rightarrow$ Relais de l'UA ou de l'ONU si nécessaire.
        \item \textbf{Citer une application concrète :} MICOPAX (2008-2013) $\rightarrow$ MISCA (UA, 2013-2014) $\rightarrow$ MINUSCA (ONU, depuis 2014) en RCA.
        \item \textbf{Évaluer les limites :} Financement (dépendance envers les partenaires extérieurs), lenteur décisionnelle, tension entre non-ingérence et ingérence humanitaire.
    \end{enumerate}
    \textbf{Formule à retenir :} \cle{Subsidiarité} : l'UA encourage les Communautés Économiques Régionales (CER) comme la CEEAC à agir en premier, au plus près des crises.
\end{methode}

\begin{exoresolu}[Le COPAX face à la crise centrafricaine (2013)]
    \textbf{Énoncé :} « Le COPAX est l'outil privilégié de la CEEAC pour la gestion des crises. » À partir de l'exemple de la crise centrafricaine de 2013, vérifiez cette affirmation en analysant la chaîne de décision et d'action.
    \tcblower
    \textbf{Solution détaillée :}
    \begin{enumerate}
        \item \textbf{Déclenchement de l'alerte (MARAC) :} Fin 2012, la progression de la coalition Séléka vers Bangui est signalée. \textbf{Action :} Rapports d'alerte précoce transmis aux instances de la CEEAC.
        \item \textbf{Mobilisation diplomatique :} \textbf{Janvier 2013 :} Sommet extraordinaire de la CEEAC à N'Djamena. \textbf{Décision :} Renforcement de la \textbf{FOMAC} (Force Multinationale de l'Afrique centrale, alors déployée sous le nom de MICOPAX) pour sécuriser Bangui et appuyer les forces centrafricaines. Mandat : protection des institutions, cessez-le-feu.
        \item \textbf{Engagement militaire :} Contingents tchadiens, congolais (Brazzaville), camerounais et gabonais. Le Cameroun déploie des éléments à la frontière Est et contribue aux effectifs.
        \item \textbf{Limites observées :} La MICOPAX fait face à l'effondrement de l'État centrafricain, aux exactions (Séléka puis Anti-Balaka), au manque de logistique (transport aérien, renseignement) et à des effectifs insuffisants (environ 2000 hommes).
        \item \textbf{Transfert de charge (Subsidiarité) :} \textbf{Juillet 2013 :} Décision de l'UA $\rightarrow$ la \textbf{MISCA} (Mission internationale de soutien à la Centrafrique sous conduite africaine) absorbe la MICOPAX. \textbf{Avril 2014 :} \textbf{Résolution 2149} $\rightarrow$ la \textbf{MINUSCA} (ONU) prend le relais en septembre 2014.
        \item \textbf{Rôle spécifique du Cameroun :} Membre du COPAX, contributeur de troupes à la FOMAC puis à la MISCA et à la MINUSCA, pays d'accueil de réfugiés centrafricains (région de l'Est), appui aux processus de paix.
    \end{enumerate}
    \textbf{Conclusion :} L'affirmation est \textbf{partiellement vraie}. Le COPAX a réagi rapidement (outil régional de premier rang), mais sa \cle{capacité de projection} et son \cle{autonomie financière} restent limitées, ce qui a nécessité le relais de l'UA puis de l'ONU. C'est la mise en œuvre concrète du principe de \cle{subsidiarité}.
\end{exoresolu}

\courssub{Méthode 3 : Comparer les cadres de coopération institutionnelle (Francophonie, Commonwealth, OCI)}
\begin{methode}[Tableau comparatif de cadres multilatéraux]
    \textbf{Objectif :} Structurer une réponse comparative rigoureuse pour le Bac.
    \begin{enumerate}
        \item \textbf{Construire un tableau à 4 colonnes :} \textbf{Critères} | \textbf{Francophonie (OIF)} | \textbf{Commonwealth} | \textbf{OCI}.
        \item \textbf{Remplir les lignes critiques :}
            \begin{itemize}
                \item \textbf{Année d'adhésion du Cameroun :} 1970 (Agence de coopération culturelle et technique, ancêtre de l'OIF), 1995 (Commonwealth), 1974 (OCI).
                \item \textbf{Fondement identitaire :} Langue française et valeurs partagées | Langue anglaise et histoire commune (ancien Empire britannique) | Solidarité islamique.
                \item \textbf{Organes principaux :} Sommet de la Francophonie, Conférence ministérielle, Secrétaire générale | Sommet des Chefs de gouvernement (CHOGM), Secrétaire général | Sommet de la Conférence islamique, Secrétaire général, comités spécialisés.
                \item \textbf{Domaines d'action prioritaires pour le Cameroun :} Éducation et enseignement supérieur (AUF), formation professionnelle, médias, appui électoral | Commerce intra-Commonwealth, jeunesse, climat, bourses d'études | Solidarité financière (Banque Islamique de Développement), médiation, culture et éducation (ICESCO).
                \item \textbf{Spécificité camerounaise :} \cle{Pays charnière} (bilingue officiel), pont entre les espaces francophone et anglophone. Yaoundé accueille le siège de la \cle{CONFEJES} (Conférence des ministres de la Jeunesse et des Sports de la Francophonie).
            \end{itemize}
        \item \textbf{Rédiger la synthèse :} 3 paragraphes (un par organisation) + 1 paragraphe transversal (complémentarité, diplomatie multidimensionnelle).
    \end{enumerate}
    \textbf{Astuce rédaction :} Utiliser les connecteurs logiques : \emph{« Tandis que... », « À l'inverse... », « En complément... »}.
\end{methode}

\begin{exoresolu}[Le Cameroun, carrefour de la Francophonie et du Commonwealth]
    \textbf{Énoncé :} En quoi l'appartenance simultanée à la Francophonie et au Commonwealth constitue-t-elle un atout diplomatique et économique majeur pour le Cameroun ?
    \tcblower
    \textbf{Solution détaillée :}
    \begin{center}
    \begin{tabularx}{\linewidth}{|l|X|X|}\hline
    \rowcolor{popblueL} \textbf{Critères} & \textbf{Francophonie (OIF)} & \textbf{Commonwealth} \\ \hline
    \textbf{Adhésion} & 1970 (membre de l'Agence de coopération culturelle et technique) & 1995 (adhésion ouverte à un pays majoritairement non anglophone d'origine) \\ \hline
    \textbf{Levier éducatif} & AUF (antenne de Yaoundé), CONFEJES, bourses de la coopération francophone & Bourses Commonwealth, partenariats universitaires avec le Royaume-Uni \\ \hline
    \textbf{Levier économique} & Espace linguistique partagé facilitant les échanges ; investissements (France, Canada/Québec, Afrique francophone) & Accès aux marchés anglophones (Nigeria, Afrique du Sud, Royaume-Uni, Canada, Inde) ; études du Commonwealth évaluant à environ 19\% la réduction des coûts de transaction entre membres \\ \hline
    \textbf{Poids diplomatique} & Tribune pour la paix et la diversité culturelle & Voix dans le concert des 56 États membres, défense des petits États, climat \\ \hline
    \textbf{Spécificité} & Siège d'institutions francophones (CONFEJES à Yaoundé) & Modèle de \cle{bilinguisme officiel} cité en exemple \\ \hline
    \end{tabularx}
    \end{center}
    \textbf{Analyse synthétique :}
    \begin{enumerate}
        \item \textbf{Diplomatie multidimensionnelle :} Le Cameroun négocie sur deux tableaux linguistiques et culturels. Exemple : dans la gestion de la crise anglophone (Nord-Ouest/Sud-Ouest), il peut solliciter l'expertise de l'OIF (décentralisation, éducation bilingue) comme celle du Commonwealth (observation électorale, gouvernance locale).
        \item \textbf{Diversification des partenariats économiques :} Réduction de la dépendance vis-à-vis de la France et de la zone franc. Ouverture vers le \cle{Nigeria} (premier partenaire commercial africain du Cameroun, membre du Commonwealth), l'\cle{Inde} et le \cle{Royaume-Uni}.
        \item \textbf{Rayonnement culturel :} Promotion du \cle{bilinguisme} comme modèle d'intégration nationale et régionale (CEMAC/CEEAC).
    \end{enumerate}
    \textbf{Conclusion :} Cette double appartenance confère au Cameroun une \cle{centralité géopolitique} unique en Afrique, lui permettant de mobiliser des ressources variées (financières, techniques, humaines) et de jouer un rôle de \cle{pont} entre l'Afrique francophone et anglophone.
\end{exoresolu}

\courssub{Méthode 4 : Évaluer une coopération bilatérale majeure (Modèle Acteurs / Instruments / Résultats)}
\begin{methode}[Grille d'analyse d'un partenariat bilatéral]
    \textbf{Objectif :} Produire une analyse structurée (ex. Cameroun-Chine, Cameroun-France, Cameroun-USA).
    \begin{enumerate}
        \item \textbf{Contexte historique :} Date d'établissement des relations, accords-cadres (traités d'amitié, commissions mixtes).
        \item \textbf{Identifier les ACTEURS :} État (Présidence, MINREX, ministères sectoriels), agences (AFD, JICA, USAID, Exim Bank of China), entreprises (BTP, énergie, télécoms), société civile et diaspora.
        \item \textbf{Lister les INSTRUMENTS :}
            \begin{itemize}
                \item \textbf{Aide publique au développement (APD) :} Dons, prêts concessionnels (taux faibles, longue durée), prêts non concessionnels.
                \item \textbf{Investissements directs étrangers (IDE) :} PPP (Partenariats Public-Privé), concessions (Port de Kribi, barrages).
                \item \textbf{Coopération technique :} Formations (bourses), expertise, volontaires.
                \item \textbf{Diplomatie économique :} Forums d'affaires, commissions mixtes.
            \end{itemize}
        \item \textbf{Quantifier les RÉSULTATS (chiffres clés Bac) :} Montants des engagements, projets phares (barrage de Nachtigal, autoroute Yaoundé-Nsimalen, hôpitaux), emplois créés, transfert de technologie.
        \item \textbf{Analyser les CRITIQUES / DÉFIS :} Endettement (ratio dette/PIB), contenu local (emploi de la main-d'œuvre locale), gouvernance (transparence des marchés), conditionnalités (gouvernance, droits de l'homme, environnement).
    \end{enumerate}
    \textbf{Mots-clés :} \cle{APD}, \cle{Prêt concessionnel}, \cle{PPP}, \cle{Contenu local}, \cle{Viabilité de la dette}.
\end{methode}

\begin{exoresolu}[La coopération sino-camerounaise : le modèle du « gagnant-gagnant » ?]
    \textbf{Énoncé :} À l'aide d'exemples précis, analysez les modalités et les enjeux de la coopération entre le Cameroun et la République populaire de Chine depuis 2000.
    \tcblower
    \textbf{Solution détaillée :}
    \begin{enumerate}
        \item \textbf{Cadre institutionnel :} Forum sur la Coopération Sino-Africaine (FOCAC, depuis 2000), Commission Mixte Cameroun-Chine, plans d'action triennaux.
        \item \textbf{Instruments financiers dominants :} \textbf{Prêts concessionnels} (Exim Bank of China) et \textbf{prêts commerciaux} (aux conditions du marché). \textbf{Peu de dons} (sauf infrastructures symboliques : Palais des Congrès de Yaoundé, stades, siège de l'Assemblée Nationale).
        \item \textbf{Projets phares (réalisations) :}
            \begin{itemize}
                \item \textbf{Énergie :} Barrage de \textbf{Memve'ele} (211 MW), participation au projet \textbf{Nachtigal} (420 MW), centrales thermiques (Kribi, Dibamba).
                \item \textbf{Transport :} Autoroute \textbf{Yaoundé-Nsimalen}, route d'accès au \textbf{port en eau profonde de Kribi} (construit par China Harbour Engineering Company), extension de l'aéroport de Yaoundé-Nsimalen.
                \item \textbf{Eau et assainissement :} Projet d'alimentation en eau potable de Yaoundé à partir de la Sanaga.
                \item \textbf{Numérique :} Réseau national à fibre optique (Huawei), centres de données.
            \end{itemize}
        \item \textbf{Bilan « Gagnant-Gagnant » :}
            \begin{itemize}
                \item \textbf{Cameroun :} Infrastructures lourdes réalisées rapidement, financement sans conditionnalités politiques explicites (principe de non-ingérence), coûts souvent compétitifs, formation d'ingénieurs.
                \item \textbf{Chine :} Accès aux ressources (bois, coton, minerais), marchés pour ses entreprises publiques, influence diplomatique (soutien au principe d'une seule Chine), excédent commercial structurel.
            \end{itemize}
        \item \textbf{Points de vigilance (critiques) :}
            \begin{itemize}
                \item \textbf{Endettement :} La Chine est le premier créancier bilatéral du Cameroun (risque de surendettement signalé par le FMI).
                \item \textbf{Contenu local :} Recours massif à la main-d'œuvre chinoise, faible sous-traitance locale, transfert de technologie limité (projets « clé en main »).
                \item \textbf{Gouvernance :} Marchés de gré à gré fréquents, opacité de certains contrats.
                \item \textbf{Environnement et social :} Déforestation (exploitation du bois), déplacements de populations (barrages), normes sociales discutées.
            \end{itemize}
        \item \textbf{Évolution récente :} FOCAC 2021 (Dakar) et 2024 (Pékin) : recentrage sur les « petits et beaux projets », l'agriculture, la santé, l'énergie verte, et annonces d'annulation des dettes à taux zéro arrivant à échéance.
    \end{enumerate}
    \textbf{Conclusion :} La coopération sino-camerounaise est un \cle{moteur de transformation structurelle} (infrastructures) mais pose un \cle{risque de dépendance financière et technologique}. Le défi camerounais : négocier des \cle{clauses de contenu local} contraignantes et diversifier les sources de financement.
\end{exoresolu}

\courssub{Méthode 5 : Analyser l'efficacité de l'aide au développement (Cadre de Paris / Accra / Busan)}
\begin{methode}[Grille d'évaluation de l'efficacité de l'aide (Dossier 3)]
    \textbf{Objectif :} Juger de la qualité de l'aide reçue au regard des principes internationaux.
    \begin{enumerate}
        \item \textbf{Rappeler les 5 principes de la \cle{Déclaration de Paris (2005)}, prolongée par Accra (2008) et Busan (2011) :}
            \begin{enumerate}
                \item \textbf{Appropriation (Ownership) :} Le pays définit sa propre stratégie (SND30).
                \item \textbf{Alignement :} Les donateurs s'alignent sur les systèmes nationaux (budget, passation des marchés, évaluation).
                \item \textbf{Harmonisation :} Les donateurs se coordonnent entre eux (réduction de la fragmentation, missions conjointes).
                \item \textbf{Gestion axée sur les résultats (GAR) :} Indicateurs de performance mesurables.
                \item \textbf{Responsabilisation mutuelle :} Redevabilité des donateurs et des bénéficiaires (revues sectorielles annuelles).
            \end{enumerate}
        \item \textbf{Appliquer au Cameroun :}
            \begin{itemize}
                \item \textbf{Cadre :} \cle{Document de Stratégie de Croissance et d'Emploi (DSCE)} puis \cle{SND30} (Stratégie Nationale de Développement 2020-2030).
                \item \textbf{Coordination :} Concertation avec les \cle{Partenaires Techniques et Financiers (PTF)} (UE, Banque Mondiale, AFD, BAD...).
                \item \textbf{Alignement budgétaire :} Part de l'aide \textbf{budgétaire} (versée au Trésor) par rapport à l'aide \textbf{projet} (gérée par le donateur). Progrès : développement de l'aide budgétaire liée aux réformes des finances publiques.
                \item \textbf{Fragmentation :} Nombreux donateurs (UE, France, Banque Mondiale, BAD, FMI, agences de l'ONU, Chine, Japon, Allemagne, USA, Suisse...). Problème : \cle{coûts de transaction} élevés pour l'administration camerounaise.
                \item \textbf{Conditionnalités :} \cle{Conditionnalités politiques} (gouvernance, droits de l'homme, lutte contre la corruption — FMI, UE, USA) contre \cle{non-conditionnalité politique} (Chine, certains pays arabes).
            \end{itemize}
        \item \textbf{Indicateurs de résultat :} Taux d'alignement sur les systèmes du pays, prévisibilité de l'aide (décaissements par rapport aux engagements), utilisation des systèmes financiers nationaux.
    \end{enumerate}
    \textbf{Formule Bac :} \cle{Efficacité de l'aide} = résultats de développement obtenus / (ressources mobilisées + coûts de transaction).
\end{methode}

\begin{exoresolu}[L'aide budgétaire contre l'aide projet au Cameroun]
    \textbf{Énoncé :} « L'aide budgétaire est plus efficace que l'aide projet pour renforcer les capacités de l'État camerounais. » Discutez cette affirmation en vous appuyant sur les principes de la Déclaration de Paris et la pratique camerounaise (2010-2024).
    \tcblower
    \textbf{Solution détaillée :}
    \begin{enumerate}
        \item \textbf{Définitions :}
            \begin{itemize}
                \item \textbf{Aide projet :} Financement d'une activité précise (ex. construction d'une école, campagne de vaccination), gérée par une unité de projet du donateur, en parallèle de l'administration. \emph{Faible alignement, forte fragmentation.}
                \item \textbf{Aide budgétaire (appui programme) :} Transfert direct au Trésor public, inscrit en Loi de Finances, géré par les procédures nationales. \emph{Fort alignement, forte appropriation.}
            \end{itemize}
        \item \textbf{Arguments POUR l'aide budgétaire (alignement / appropriation) :}
            \begin{itemize}
                \item Renforce les \textbf{systèmes nationaux} (budget, comptabilité, audit, passation des marchés). Exemple : appui à la réforme de la \cle{chaîne de la dépense publique}.
                \item Réduit les \textbf{coûts de transaction} (procédures et rapports unifiés).
                \item Soutient les \textbf{réformes structurelles} (matrices de réformes liées aux décaissements du FMI, de l'UE, de la Banque Mondiale). Exemples : adhésion à l'initiative de transparence des industries extractives (ITIE), renforcement de la Cour des Comptes.
                \item Prévisibilité : engagements pluriannuels inscrits dans le Cadre de Dépenses à Moyen Terme (CDMT).
            \end{itemize}
        \item \textbf{Arguments CONTRE / Limites (risques) :}
            \begin{itemize}
                \item \textbf{Fongibilité :} L'argent versé libère des ressources nationales qui peuvent être affectées à des dépenses non prioritaires si la \textbf{gouvernance est faible}.
                \item \textbf{Conditionnalités lourdes :} Les « \cle{actions préalables} » au décaissement sont parfois perçues comme une ingérence (ex. réforme des prix des carburants, réforme des retraites).
                \item \textbf{Capacité d'absorption :} Trésorerie tendue $\rightarrow$ retards de paiement des fournisseurs (arriérés intérieurs) malgré l'aide budgétaire.
                \item \textbf{Suspensions possibles :} Gel des décaissements en cas d'écarts par rapport au programme (instabilité du financement).
            \end{itemize}
        \item \textbf{Réalité camerounaise :} \textbf{Hybridation nécessaire.}
            \begin{itemize}
                \item Aide budgétaire : secteurs « mûrs » (santé, éducation, finances publiques) via les appuis sectoriels.
                \item Aide projet : infrastructures lourdes (barrages, routes, ports) où l'État manque d'expertise de maîtrise d'ouvrage (Chine, BAD, BID, AFD).
            \end{itemize}
    \end{enumerate}
    \textbf{Conclusion :} L'affirmation est \textbf{théoriquement fondée} (principes d'alignement et d'appropriation) mais \textbf{pratiquement conditionnée} à la qualité de la \cle{gouvernance économique et financière}. Le Cameroun vise la \cle{mutualisation} (fonds communs sectoriels : éducation, santé) pour combiner les avantages des deux modalités.
\end{exoresolu}

\courssub{Méthode 6 : Rédiger une analyse de contribution à la gestion d'un conflit (TD2)}
\begin{methode}[Structure type « Étude de cas : Médiation / Maintien de la paix »]
    \textbf{Objectif :} Répondre à une question de type « Montrez la contribution du Cameroun à la résolution du conflit X ».
    \begin{enumerate}
        \item \textbf{Introduction :} Contexte du conflit (acteurs, causes, dates), position officielle du Cameroun (communiqués du MINREX, discours du Chef de l'État), cadre juridique (Acte constitutif de l'UA, Charte de l'ONU chapitre VIII, Traité de la CEEAC).
        \item \textbf{Volet diplomatique et politique :}
            \begin{itemize}
                \item Médiation directe (Chef de l'État, envoyés spéciaux).
                \item Accueil de pourparlers (Yaoundé, Garoua).
                \item Résolutions soutenues ou parrainées par le Cameroun (UA, CEEAC, ONU).
            \end{itemize}
        \item \textbf{Volet militaire et sécuritaire :}
            \begin{itemize}
                \item Contribution de troupes et de policiers (FOMAC, MISCA, MINUSCA, Force Multinationale Mixte du bassin du Lac Tchad).
                \item Sécurisation des frontières (opérations nationales de lutte contre le terrorisme et la criminalité transfrontalière).
                \item Désarmement, Démobilisation, Réinsertion et Rapatriement (DDRR).
            \end{itemize}
        \item \textbf{Volet humanitaire et social :}
            \begin{itemize}
                \item Accueil des réfugiés et des personnes déplacées internes (politique de « porte ouverte », sites de Minawao, Gado, Borgop).
                \item Appui logistique (bases de Garoua et de Douala pour les corridors humanitaires).
                \item Réconciliation communautaire (chefs traditionnels, leaders religieux, ONG locales).
            \end{itemize}
        \item \textbf{Bilan critique :} Succès (accords signés, zones sécurisées), échecs et limites (reprise des violences, impunité, coût budgétaire, radicalisation résiduelle).
        \item \textbf{Conclusion :} Apport spécifique du Cameroun (\cle{puissance médiane}, \cle{ancrage régional}, \cle{diplomatie de proximité}).
    \end{enumerate}
    \textbf{Connecteurs imposés :} \emph{« En premier lieu... », « Parallèlement... », « Toutefois... », « Au final... »}.
\end{methode}

\begin{exoresolu}[La contribution du Cameroun à la lutte contre Boko Haram (Force Multinationale Mixte)]
    \textbf{Énoncé :} Évaluez la contribution du Cameroun au sein de la Force Multinationale Mixte (FMM) face à l'insurrection de Boko Haram dans le bassin du Lac Tchad (2015-2024).
    \tcblower
    \textbf{Solution détaillée :}
    \begin{enumerate}
        \item \textbf{Cadre juridique et institutionnel :}
            \begin{itemize}
                \item Réactivation de la Force Multinationale Mixte par le Nigeria, le Niger, le Tchad, le Cameroun et le Bénin (\textbf{2014-2015}), sous l'égide de la Commission du Bassin du Lac Tchad (CBLT).
                \item Endossement par le Conseil de Paix et de Sécurité de l'\textbf{UA} (2015) et appui de l'\textbf{ONU} (Résolution 2349 de 2017 sur le bassin du Lac Tchad). Quartier général : \textbf{N'Djamena}. Secteur 1 (Cameroun) : QG à \textbf{Mora}.
            \end{itemize}
        \item \textbf{Contribution militaire (Secteur 1 — Mora) :}
            \begin{itemize}
                \item \textbf{Effectifs :} Environ 2000 à 3000 hommes (BIR — Bataillon d'Intervention Rapide, bataillons d'infanterie, génie, artillerie).
                \item \textbf{Opérations :} Opérations nationales de sécurisation de la frontière et opérations conjointes dans les îles du Lac Tchad.
                \item \textbf{Résultats :} Reprise de localités (Fotokol, Amchidé), destruction de bases, affaiblissement des chefs (mort d'Abubakar Shekau en 2021 dans un contexte de rivalité avec l'ISWAP), libération d'otages.
            \end{itemize}
        \item \textbf{Contribution au renseignement et à la logistique :}
            \begin{itemize}
                \item Partage de renseignement entre les pays membres de la FMM.
                \item Bases de \textbf{Garoua} et \textbf{Maroua} : plateformes logistiques, évacuations sanitaires, surveillance aérienne.
                \item Corridor Douala-N'Djamena (port, rail, route) pour le ravitaillement des forces.
            \end{itemize}
        \item \textbf{Contribution humanitaire et à la stabilisation :}
            \begin{itemize}
                \item \textbf{Accueil :} Plus de 100 000 réfugiés nigérians (camp de Minawao) et plusieurs centaines de milliers de déplacés internes camerounais (Extrême-Nord).
                \item \textbf{Reconstruction :} \cle{Plan d'Urgence Triennal (PLANUT)} et programmes de relèvement du Nord : reconstruction d'écoles, de centres de santé, de marchés.
                \item \textbf{DDR :} Dispositifs de désengagement et de réinsertion des ex-combattants repentis.
            \end{itemize}
        \item \textbf{Bilan critique (coûts et limites) :}
            \begin{itemize}
                \item \textbf{Humain :} Des milliers de morts (civils et militaires) côté camerounais depuis 2014.
                \item \textbf{Financier :} Un effort de guerre et de reconstruction pesant lourdement sur le budget de l'État.
                \item \textbf{Sécuritaire :} Résilience des factions (ISWAP, JAS), engins explosifs improvisés (IED), radicalisation des jeunes, porosité des frontières.
                \item \textbf{Gouvernance :} Accusations de violations des droits de l'homme (rapports d'ONG) $\rightarrow$ risque de suspension de certaines aides militaires partenaires (ex. loi Leahy aux USA).
            \end{itemize}
    \end{enumerate}
    \textbf{Conclusion :} Le Cameroun est le \cle{pilier occidental} de la FMM (géographie, logistique, troupes aguerries). Sa contribution est \cle{déterminante mais insuffisante} sans règlement politique des causes profondes (pauvreté, exclusion, gouvernance locale) et sans justice transitionnelle crédible.
\end{exoresolu}

\courssub{Méthode 7 : Synthétiser les problèmes et défis des relations internationales du Cameroun (Dossier 5)}
\begin{methode}[Classification thématique des défis diplomatiques]
    \textbf{Objectif :} Produire une synthèse structurée pour une dissertation ou une question de synthèse.
    \begin{enumerate}
        \item \textbf{Classez les défis en 4 catégories majeures :}
            \begin{enumerate}
                \item \textbf{Défis sécuritaires (exogènes et endogènes) :} Terrorisme (Boko Haram, ISWAP), insécurité maritime (Golfe de Guinée), crise anglophone (Nord-Ouest/Sud-Ouest), réfugiés et déplacés internes, criminalité transfrontalière (trafics d'armes, de drogue, de bois, d'or).
                \item \textbf{Défis économiques et financiers :} \cle{Endettement extérieur} (risque de surendettement signalé par le FMI, poids de la Chine), \cle{dépendance extravertie} (exportations brutes : pétrole, bois, cacao, coton, aluminium), \cle{balance commerciale} structurellement déficitaire hors pétrole, \cle{fuites illicites de capitaux}, \cle{climat des affaires} dégradé, \cle{concentration des partenaires} (UE, Chine).
                \item \textbf{Défis politiques et institutionnels (image et soft power) :} \cle{Gouvernance} (corruption — indice CPI de Transparency International, justice, droits de l'homme), \cle{crise anglophone} (perception internationale, conditionnalités de l'aide), \cle{respect des engagements internationaux} (droits de l'homme, climat, ODD).
                \item \textbf{Défis diplomatiques et stratégiques :} \cle{Marginalisation dans les nouvelles architectures} (G20, BRICS+), \cle{rivalités des grandes puissances} (Occident / Chine / Russie — neutralité difficile), \cle{intégration régionale inachevée} (CEMAC/CEEAC : libre circulation freinée, infrastructures manquantes), \cle{contentieux frontaliers} (Bakassi réglé par l'arrêt CIJ de 2002 et l'accord de Greentree de 2006, mais tensions ponctuelles ; frontière Est instable).
            \end{enumerate}
        \item \textbf{Hiérarchisez :} Court terme (sécurité, humanitaire) contre moyen et long terme (structurel : dette, industrialisation, gouvernance).
        \item \textbf{Proposez des pistes de réponse (prospective) :} Diplomatie économique offensive, réformes de gouvernance, accélération de la ZLECAf, renforcement des capacités des diplomates (ENAM/IRIC), mobilisation de la diaspora.
    \end{enumerate}
    \textbf{Mots-clés transverses :} \cle{Résilience}, \cle{autonomisation stratégique}, \cle{diplomatie économique}, \cle{paix positive}.
\end{methode}

\begin{exoresolu}[Les défis majeurs de la politique étrangère camerounaise à l'horizon 2030]
    \textbf{Énoncé :} « La politique étrangère du Cameroun fait face à une crise de l'efficacité. » En vous appuyant sur le Dossier 5, identifiez trois défis structurels et proposez pour chacun une réponse stratégique adaptée.
    \tcblower
    \textbf{Solution détaillée :}
    \textbf{Introduction :} La SND30 vise l'émergence du Cameroun à l'horizon 2035. La diplomatie doit être le levier extérieur de cette transformation. Le diagnostic révèle un écart entre le potentiel (ressources, position géographique) et les résultats (investissements, influence, sécurité).
    \begin{enumerate}
        \item \textbf{Défi 1 : Le piège de l'endettement et de la dépendance extravertie (économique).}
            \begin{itemize}
                \item \textbf{Constat :} Dette publique élevée (risque de surendettement selon le FMI), service de la dette absorbant une part importante des recettes budgétaires. Exportations dominées par les produits bruts (pétrole, bois, cacao, aluminium). Partenaires concentrés (UE, Chine, Inde).
                \item \textbf{Réponse stratégique :} \textbf{Diplomatie économique de rupture.}
                    \begin{itemize}
                        \item Négociation de clauses de \cle{contenu local} et de \cle{transfert de technologie} dans tous les grands contrats (mines, énergie, BTP).
                        \item Promotion active de la \cle{ZLECAf} (Zone de Libre-Échange Continentale Africaine) : cibler les marchés du Nigeria, de l'Égypte et de l'Afrique du Sud pour diversifier les exportations transformées (agro-industrie, bois de 2\textsuperscript{e} et 3\textsuperscript{e} transformation).
                        \item Gestion rigoureuse de la dette : \cle{rééchelonnements}, \cle{conversions dette-investissement} (échanges dette-nature, dette-climat), plafonnement des emprunts non concessionnels.
                    \end{itemize}
            \end{itemize}
        \item \textbf{Défi 2 : L'insécurité multidimensionnelle et l'érosion du « soft power » (sécuritaire et politique).}
            \begin{itemize}
                \item \textbf{Constat :} Trois fronts actifs (Extrême-Nord / Boko Haram, Nord-Ouest et Sud-Ouest / crise anglophone, Est / frontière centrafricaine). Coût humain et financier massif. Image internationale dégradée (rapports d'ONG sur les droits de l'homme, indices de corruption). Conditionnalités de l'aide (UE, USA, FMI).
                \item \textbf{Réponse stratégique :} \textbf{Diplomatie de la responsabilité et de la prévention.}
                    \begin{itemize}
                        \item \textbf{Justice transitionnelle crédible} (vérité, réconciliation, réparation des victimes) pour lever les conditionnalités.
                        \item \textbf{Diplomatie préventive régionale :} Leadership effectif dans le COPAX, la FMM et la Commission du Bassin du Lac Tchad (CBLT) : mutualisation du renseignement, sécurisation des corridors humanitaires, développement transfrontalier.
                        \item \textbf{Communication stratégique :} Narratif national unifié (ambassadeurs, diaspora, think-tanks comme l'IRIC) pour défendre l'image du pays.
                    \end{itemize}
            \end{itemize}
        \item \textbf{Défi 3 : L'inadaptation de l'appareil diplomatique aux nouveaux enjeux (institutionnel et stratégique).}
            \begin{itemize}
                \item \textbf{Constat :} Diplomatie « de cérémonie » plutôt que « diplomatie d'affaires et technique ». Sous-représentation dans les organisations techniques (OMC, OMS, UIT). Faible mobilisation de la diaspora qualifiée (transferts importants mais peu d'investissement productif). Absence des grands clubs (G20, BRICS+).
                \item \textbf{Réponse stratégique :} \textbf{Refonte de l'outil diplomatique (modernisation).}
                    \begin{itemize}
                        \item \textbf{Recrutement et formation :} Profils d'économistes, juristes internationaux et ingénieurs ; formation continue (numérique, climat, négociation commerciale) à l'ENAM et à l'IRIC.
                        \item \textbf{Diplomatie économique décentralisée :} Attachés commerciaux efficaces, missions économiques ciblées (FOCAC, sommets UK-Africa, salons technologiques).
                        \item \textbf{Engagement de la diaspora :} Cadre institutionnel dédié, vote de la diaspora, émissions obligataires diaspora (« Diaspora Bonds ») pour financer les infrastructures.
                        \item \textbf{Candidatures stratégiques :} Postes clés à l'ONU, l'UA, l'OIF, l'OCI et l'OMC pour peser sur l'élaboration des normes internationales.
                    \end{itemize}
            \end{itemize}
    \end{enumerate}
    \textbf{Conclusion :} Le Cameroun doit passer d'une \cle{diplomatie de position} (héritage historique) à une \cle{diplomatie de projection} (intérêts vitaux, performance, résultats mesurables). La réussite de la SND30 en dépend.
\end{exoresolu}

\begin{attention}[Les 5 erreurs qui coûtent des points au Bac]
    \begin{enumerate}
        \item \textbf{Confondre les cadres institutionnels :} Mélanger les organes de l'\textbf{UA} (Conseil de Paix et de Sécurité, Assemblée de l'Union, Cour de justice) et ceux de la \textbf{CEEAC/COPAX} (Conférence des Chefs d'État, COPAX, MARAC, FOMAC). \emph{Parade :} Faire un tableau comparatif des deux architectures (niveau continental contre niveau régional) et mémoriser les sièges (Addis-Abeba pour l'UA, Libreville pour la CEEAC).
        \item \textbf{Citer des chiffres approximatifs ou datés :} Dire « la Chine est le premier partenaire » sans nuancer (premier pour les \emph{importations} et les \emph{infrastructures}, mais l'\textbf{UE} reste un partenaire majeur pour \emph{l'APD} et les \emph{exportations} hors pétrole). \emph{Parade :} Mémoriser 3 ou 4 repères chiffrés récents (part de la Chine dans la dette bilatérale, montants du FOCAC, IDE français, APD de l'UE).
        \item \textbf{Oublier la distinction « Aide / Investissement / Prêt » :} Traiter l'argent chinois comme de l'« aide » (don). \emph{Parade :} Utiliser le vocabulaire technique exact : \cle{prêt concessionnel}, \cle{prêt commercial}, \cle{don}, \cle{IDE}, \cle{PPP}. La très grande majorité du financement chinois est de la \textbf{dette souveraine}.
        \item \textbf{Réponse descriptive au lieu d'analytique :} Lister les organisations auxquelles le Cameroun appartient (ONU, UA, CEMAC, CEEAC, OIF, Commonwealth, OCI...) sans expliquer \emph{ce que le Cameroun y fait} (négociations, leadership, retombées concrètes). \emph{Parade :} Appliquer la méthode \textbf{Acteur / Instrument / Résultat / Limite} pour chaque partenariat.
        \item \textbf{Négliger le volet « Problèmes / Défis » (Dossier 5) :} Présenter une vision idyllique (« tout va bien, le Cameroun est respecté »). \emph{Parade :} Intégrer systématiquement une partie « Limites / Défis / Conditionnalités » dans chaque développement (gouvernance, dette, droits de l'homme, sécurité, fragmentation de l'aide). Le correcteur attend un \textbf{esprit critique}.
    \end{enumerate}
\end{attention}

\newpage

\coursec{Exercices}

\courssub{Je vérifie mes connaissances}

\begin{exercice}[QCM : Acteurs et cadres de la coopération]
\emph{Choisissez la bonne réponse parmi les propositions.}
\begin{enumerate}
    \item L'organe principal de l'ONU chargé du maintien de la paix et de la sécurité internationale est :
    \begin{itemize}
        \item[A.] L'Assemblée générale
        \item[B.] Le Conseil de sécurité
        \item[C.] Le Secrétariat
        \item[D.] La Cour internationale de justice
    \end{itemize}
    \item Le \cle{COPAX} (Conseil de paix et de sécurité de l'Afrique centrale) est une institution de :
    \begin{itemize}
        \item[A.] La CEMAC
        \item[B.] La CEEAC
        \item[C.] L'Union Africaine
        \item[D.] La CIRGL
    \end{itemize}
    \item L'accord de partenariat qui lie les pays ACP (Afrique, Caraïbes, Pacifique) à l'Union européenne s'appelle actuellement :
    \begin{itemize}
        \item[A.] Convention de Lomé
        \item[B.] Accord de Cotonou
        \item[C.] Accord de partenariat UE-ACP (Samoa)
        \item[D.] Accord d'association euro-méditerranéen
    \end{itemize}
    \item Le Cameroun est membre de la \cle{Francophonie} et du \cle{Commonwealth}. Cette particularité s'explique par :
    \begin{itemize}
        \item[A.] Son adhésion simultanée aux deux organisations en 1960
        \item[B.] Son histoire coloniale partagée entre la France et le Royaume-Uni
        \item[C.] Une décision de l'ONU en 1961
        \item[D.] Sa position géographique de charnière
    \end{itemize}
    \item L'\cle{APD} (Aide publique au développement) se caractérise principalement par :
    \begin{itemize}
        \item[A.] Un taux d'intérêt de marché et une maturité courte
        \item[B.] Un don pur et simple sans contrepartie
        \item[C.] Un caractère concessionnel (don ou prêt à taux très faible, longue maturité) et l'objectif de développement
        \item[D.] Un investissement direct étranger dans le secteur privé
    \end{itemize}
\end{enumerate}
\end{exercice}

\begin{exercice}[Vrai ou Faux justifié : Réalités de la coopération]
\emph{Indiquez si les affirmations sont vraies ou fausses et justifiez brièvement votre réponse en vous appuyant sur le programme.}
\begin{enumerate}
    \item Le Cameroun n'a jamais fourni de contingents militaires aux opérations de maintien de la paix (OMP) de l'ONU.
    \item La Chine est devenue le premier partenaire commercial et le premier bailleur de fonds bilatéral du Cameroun en volume d'engagements récents (projets d'infrastructures).
    \item L'aide du FMI au Cameroun (via la FEC - Facilité élargie de crédit) est dénuée de toute conditionnalité macroéconomique.
    \item Le \cle{COPAX} a pour mission exclusive la gestion des catastrophes naturelles en Afrique centrale.
    \item La « coopération triangulaire » (ex: Cameroun - Brésil - Pays tiers) vise à mobiliser l'expertise d'un pays émergent au profit d'un troisième pays avec l'appui d'un bailleur traditionnel.
\end{enumerate}
\end{exercice}

\begin{exercice}[Définitions et mise en relation : Concepts clés]
\emph{Définissez précisément les termes suivants (2 à 3 lignes chacun) puis reliez chaque concept à l'institution ou au partenaire principal concerné.}
\begin{enumerate}
    \item \cle{Conditionnalité} (dans le cadre de l'APD / FMI)
    \item \cle{Diplomatie économique}
    \item \cle{Sécurité collective} (au sens de l'UA / ONU)
    \item \cle{Partenariat gagnant-gagnant} (rhétorique sino-africaine)
    \item \cle{Intégration sous-régionale} (CEMAC / CEEAC)
\end{enumerate}
\emph{Reliez :}
\begin{center}
\begin{tabularx}{\linewidth}{|l|X|}\hline
Concept & Institution / Partenaire principal \\ \hline
COPAX & \\ \hline
FEC (Facilité élargie de crédit) & \\ \hline
OMP (Casques bleus) & \\ \hline
Forum sur la coopération sino-africaine (FOCAC) & \\ \hline
Accord de Cotonou / Samoa & \\ \hline
\end{tabularx}
\end{center}
\end{exercice}

\begin{exercice}[Repères chiffrés et chronologiques]
\emph{Complétez le tableau ci-dessous avec les données exactes (années, sigles, montants approximatifs ou pourcentages) issues du cours.}
\begin{center}
\begin{tabularx}{\linewidth}{|l|X|}\hline
\textbf{Indicateur / Événement} & \textbf{Valeur / Date / Sigle} \\ \hline
Année d'adhésion du Cameroun à l'ONU & \\ \hline
Année d'adhésion à l'UA (ex-OUA) & \\ \hline
Année d'adhésion au Commonwealth & \\ \hline
Part approximative de la dette extérieure camerounaise détenue par la Chine (environ) & \\ \hline
Nombre de piliers de la stratégie FEC 2021-2024 (FMI) & \\ \hline
Nom de la force multinationale mixte (Lac Tchad) & \\ \hline
Sigle de l'instrument financier de l'UE pour l'APD (2021-2027) & \\ \hline
\end{tabularx}
\end{center}
\end{exercice}

\courssub{Je m'entraîne}

\begin{exercice}[Analyse de document : Le discours de la Tribune onusienne]
\emph{Document : Extrait du discours du Chef de l'État camerounais à la 78e Assemblée générale de l'ONU (septembre 2023).}
\begin{quote}
\emph{« Le Cameroun réaffirme son attachement au multilatéralisme [...] Nous plaidons pour une réforme du Conseil de sécurité afin qu'il reflète les réalités géopolitiques actuelles, avec une représentation africaine permanente. [...] Notre pays continue de payer un lourd tribut à la lutte contre le terrorisme et l'extrémisme violent, notamment dans le bassin du Lac Tchad et dans les régions du Nord-Ouest et du Sud-Ouest. [...] Nous appelons à un partenariat rénové, fondé sur le respect de la souveraineté et l'équité. »}
\end{quote}
\begin{enumerate}
    \item Identifiez les trois registres de l'action diplomatique camerounaise évoqués dans cet extrait (sécurité, institutionnel, économique).
    \item Expliquez la demande de « réforme du Conseil de sécurité » : quel est l'enjeu pour l'Afrique (consensus d'Ezulwini) ?
    \item Montrez en quoi la mention du « lourd tribut » dans le bassin du Lac Tchad et les régions anglophones illustre le lien entre \cle{sécurité intérieure} et \cle{diplomatie multilatérale}.
    \item Quel est le sens de l'expression « partenariat rénové fondé sur le respect de la souveraineté » vis-à-vis des partenaires traditionnels (UE, FMI) et nouveaux (Chine) ?
\end{enumerate}
\end{exercice}

\begin{exercice}[Étude de cas : La coopération Cameroun - Chine : infrastructures et dette]
\emph{Contexte : Le Cameroun a signé de nombreux accords de prêts concessionnels et commerciaux avec la Chine (Exim Bank, China Development Bank) pour financer des projets structurants : barrage de Memve'ele, port en eau profonde de Kribi (phase 1), autoroute Yaoundé-Nsimalen, stade d'Olembé, centrale hydroélectrique de Nachtigal (partenaire chinois minoritaire), hôpital de référence de Yaoundé.}
\begin{enumerate}
    \item Citez trois avantages majeurs de ce partenariat pour le \cle{développement} camerounais (délais, coûts, technologie, conditionnalité politique).
    \item Identifiez deux risques majeurs soulignés par les rapports de la Banque mondiale et du FMI (soutenabilité de la dette, opacité des clauses, main-d'œuvre importée, impact environnemental).
    \item Le concept de « \cle{piège de la dette} » (debt-trap diplomacy) est-il pertinent pour analyser la situation camerounaise ? Nuancez votre réponse en distinguant dette concessionnelle et dette commerciale.
    \item Proposez deux mesures de \cle{gouvernance} que le Cameroun pourrait prendre pour optimiser le retour sur investissement de ces projets (ex: clause de contenu local, audit indépendant, rééchelonnement).
\end{enumerate}
\end{exercice}

\begin{exercice}[Schéma synthèse : Architecture institutionnelle de la coopération]
\emph{Réalisez un organigramme complet (à l'aide d'un schéma arborescent ou d'un tableau structuré) classant l'ensemble des partenaires et cadres de coopération du Cameroun selon les trois niveaux suivants. Utilisez les sigles officiels.}
\begin{enumerate}
    \item \textbf{Niveau Universel (Multilatéral global)} : ONU (AG, CS, ECOSOC, OMS, UNESCO, PNUD, OMP), FMI, BM, OMC.
    \item \textbf{Niveau Régional et Sous-régional (Afrique)} : UA (CPS/COPAX), CEEAC, CEMAC (BEAC, COBAC), CLC (Commission Lac Tchad), FOMAC.
    \item \textbf{Niveau Bilatéral et Interrégional} : Partenaires historiques (France, UK), Partenaires européens (UE), Partenaires émergents (Chine, Japon - TICAD, Brésil, Turquie, Russie, Inde, Corée du Sud), Blocs (Francophonie, Commonwealth, OCI).
\end{enumerate}
\emph{Indiquez pour chaque niveau un objectif stratégique principal pour le Cameroun (ex: Paix/Sécurité, Intégration économique, Diversification/Désendettement).}
\end{exercice}

\begin{exercice}[Argumentation dirigée : Diversification vs Dépendance]
\emph{Sujet : « La diversification des partenariats internationaux (Chine, Brésil, Turquie, Russie, Inde) permet-elle au Cameroun de réduire sa dépendance vis-à-vis de ses partenaires traditionnels (France, UE, FMI) ou crée-t-elle de nouvelles vulnérabilités ? »}
\emph{Rédigez un paragraphe argumenté d'une vingtaine de lignes structuré ainsi :}
\begin{enumerate}
    \item Thèse : La diversification offre des marges de manœuvre (alternatives de financement, non-alignement).
    \item Antithèse : Elle génère de nouveaux risques (fragmentation de l'aide, dette opaque, alignement géopolitique forcé, affaiblissement des normes sociales/environnementales).
    \item Synthèse : La \cle{souveraineté} réside dans la capacité à \emph{négocier} et \emph{coordonner} l'ensemble des flux (Cadre national de partenariat) plutôt qu'à les substituer.
\end{enumerate}
\end{exercice}

\courssub{Je me prépare au Bac}

\begin{exercice}[Sujet type Bac : Analyse de documents - Diversification des partenariats et nouveaux risques]
\emph{Durée conseillée : 1h30. Vous disposez de trois documents.}
\vspace{0.5cm}
\noindent\textbf{Document 1 : Carte schématique des flux d'APD nette reçus par le Cameroun (Moyenne 2018-2022, en millions USD).}
\begin{center}
\begin{tabularx}{\linewidth}{|l|X|}\hline
\textbf{Bailleur} & \textbf{Montant (M USD) / Part \%} \\ \hline
Union Européenne (Institutions + États membres) & 350 / 35\% \\ \hline
France (AFD, Budget) & 180 / 18\% \\ \hline
Banque Mondiale (IDA) & 200 / 20\% \\ \hline
FMI (Facilités) & 50 / 5\% \\ \hline
Chine (Prêts concessionnels + Dons) & 120 / 12\% \\ \hline
Autres (BAfD, Pays arabes, Japon, USA, ONU) & 100 / 10\% \\ \hline
\end{tabularx}
\end{center}
\emph{Source : OCDE/DAC, données retraitées.}
\vspace{0.5cm}
\noindent\textbf{Document 2 : Graphique - Évolution de l'encours de la dette publique extérieure du Cameroun vis-à-vis de la Chine (2010-2023).}
\begin{center}
\begin{tikzpicture}
\begin{axis}[
    width=\linewidth, height=6cm,
    xlabel={Année}, ylabel={Milliards USD},
    xtick={2010,2012,2014,2016,2018,2020,2022},
    ytick={0,1,2,3,4,5},
    ymin=0, ymax=5.5,
    grid=major,
    popcurve,
    thick
]
\addplot coordinates {
(2010,0.3) (2012,0.8) (2014,1.5) (2016,2.2) (2018,3.0) (2020,3.8) (2022,4.2) (2023,4.0)
};
\end{axis}
\end{tikzpicture}
\legende{Évolution de la dette envers la Chine (Données simplifiées FMI/Banque Mondiale).}
\end{center}
\vspace{0.5cm}
\noindent\textbf{Document 3 : Extrait du discours du Président de la République à l'Assemblée Générale de l'ONU (2022).}
\begin{quote}
\emph{« L'Afrique ne veut pas être le théâtre de la rivalité des grandes puissances. Elle aspire à des partenariats gagnant-gagnant, respectueux de sa souveraineté et de ses priorités de développement définies dans l'Agenda 2063. Le Cameroun, fidèle à sa vocation de paix, continuera d'œuvrer pour un multilatéralisme efficace, seul garant de la sécurité collective. »}
\end{quote}
\vspace{0.5cm}
\noindent\textbf{Travail demandé :}
\begin{enumerate}
    \item \textbf{Analyse des documents (6 pts) :}
    \begin{enumerate}
        \item À partir du Document 1, calculez la part combinée des partenaires « traditionnels » (UE + France + FMI + BM) et comparez-la à celle de la Chine. Que révèle ce chiffre sur la réalité de la \cle{diversification} ?
        \item À partir du Document 2, calculez le taux de croissance annuel moyen (TCAM) de la dette envers la Chine sur la période 2010-2023. Interprétez l'infléchissement observé en 2023.
        \item Identifiez dans le Document 3 les trois principes directeurs de la diplomatie camerounaise énoncés par le Chef de l'État.
    \end{enumerate}
    \item \textbf{Mise en relation et synthèse (8 pts) :}
    \emph{Montrez, en vous appuyant sur l'ensemble des documents et vos connaissances, que la diversification des partenariats est une stratégie de réduction de la dépendance pour le Cameroun, tout en soulignant les nouveaux risques qu'elle génère (soutenabilité de la dette, fragmentation de l'aide, alignements géopolitiques, conditionnalités différenciées).}
    \item \textbf{Écriture d'une recommandation (6 pts) :}
    \emph{En tant qu'expert en relations internationales, rédigez une note de synthèse à l'intention du Ministre des Relations Extérieures proposant trois axes prioritaires pour une \cle{gouvernance} optimisée de la coopération multilatérale et bilatérale du Cameroun à l'horizon 2030.}
\end{enumerate}
\end{exercice}

\begin{exercice}[Sujet type Bac : Dissertation courte - OMP et Intérêt national]
\emph{Sujet : « La participation du Cameroun aux opérations de maintien de la paix (OMP) de l'ONU : entre devoir moral, intérêt national et contrainte capacitaire. Discutez. »}
\emph{Durée conseillée : 1h30.}
\emph{Consignes :}
\begin{enumerate}
    \item Analysez le sujet : définissez les termes (OMP, devoir moral/R2P, intérêt national : formation, équipement, rayonnement, revenus, contrainte capacitaire : logistique, formation, pertes humaines, fatigue opérationnelle).
    \item Proposez un plan détaillé (Introduction, 2 ou 3 parties équilibrées, Conclusion).
    \item Rédigez l'introduction (accroche, définition, problématique, annonce du plan).
    \item Rédigez le développement en mobilisant des exemples précis : MINUSCA (RCA), MINUSMA (Mali - retrait), MONUSCO (RDC), MINURSO (Sahara), Force multinationale mixte (Lac Tchad - cadre UA/ONU), contributions policières (FPU), hôpital de campagne.
    \item Rédigez la conclusion (bilan, ouverture sur la réforme des OMP / Partenariat stratégique pour les opérations de paix).
\end{enumerate}
\end{exercice}

\begin{exercice}[Sujet type Bac : Croquis commenté - Le Cameroun, carrefour de coopérations]
\emph{Consigne : « Réalisez un croquis commenté intitulé : \textbf{Le Cameroun, carrefour de coopérations multilatérales et bilatérales}. La légende organisée en 3 items doit faire apparaître :}
\begin{enumerate}
    \item \textbf{Ancrages institutionnels} (Sièges, organisations, statuts).
    \item \textbf{Flux financiers et projets structurants} (APD, Investissements, Dette, Infrastructures).
    \item \textbf{Enjeux sécuritaires et diplomatiques} (OMP, Frontières, Conflits, Médiation).
\end{enumerate}
\emph{Le croquis sera réalisé sur une carte de base du Cameroun (contour, capitales régionales, villes principales, axes routiers/portuaires, pays limitrophes). Vous utiliserez des figurés de surface, linéaires et ponctuels adaptés. La légende sera hiérarchisée et quantitative si possible. »}
\vspace{0.5cm}
\emph{Indications pour la réalisation (non fournies à l'élève le jour J, mais utiles pour la révision) :}
\begin{itemize}
    \item \textbf{Fond de carte :} Contour pays, Yaoundé (capitale politique), Douala (capitale économique), Garoua, Maroua, Ngaoundéré, Bertoua, Bafoussam, Bamenda, Buea, Ebolowa, Kribi. Pays limitrophes : Nigeria, Tchad, RCA, Congo, Gabon, Guinée équatoriale.
    \item \textbf{Item 1 (Ancrages) :} Point rouge à Yaoundé (Siège UA - Représentation, Siège CEMAC, Siège COPAX - État-major FOMAC). Drapeaux/Logos : ONU, UA, CEEAC, CEMAC, Francophonie, Commonwealth, OCI, UE (Délégation). Flèches entrantes : Ambassades/Représentations.
    \item \textbf{Item 2 (Flux/Projets) :} Flèches proportionnelles (largeur) vers Douala/Kribi (Port, Zone industrielle, Pipeline Tchad-Cameroun, Barrage Memve'ele, Autoroute, Stade Olembé, Centrale Nachtigal). Pictogrammes : Usine, Barrage, Port, Route, Hôpital. Codes couleur : Bleu (UE/France), Rouge (Chine), Vert (Banque Mondiale/BAfD), Orange (Pays émergents). Bulles chiffrées : Montants APD / Investissements.
    \item \textbf{Item 3 (Sécurité) :} Zones hachurées : Extrême-Nord (Boko Haram), Nord-Ouest/Sud-Ouest (Crise anglophone), Est (Réfugiés RCA). Flèches sortantes : Contingents OMP (RCA, Mali, RDC, Sahara). Icônes : Casque bleu, Main serrée (Médiation), Base FOMAC (Libreville/Yaoundé).
\end{itemize}
\end{exercice}

\newpage

\coursec{Corrigés des exercices}

\coursec{Le Cameroun dans la coopération multilatérale}
\courssub{Corrigés des exercices d'application et sujets type Baccalauréat}

\begin{corrige}[Exercice 1 — QCM : Acteurs et cadres de la coopération]
\begin{enumerate}
    \item \textbf{Réponse B — Le Conseil de sécurité.} Selon la Charte des Nations unies (chapitre VII), le Conseil de sécurité a la \cle{responsabilité principale} du maintien de la paix et de la sécurité internationales ; il peut autoriser des sanctions, des embargos et des opérations de maintien de la paix (OMP). L'Assemblée générale délibère mais ne décide pas des OMP ; le Secrétariat exécute ; la CIJ tranche les différends juridiques entre États.
    \item \textbf{Réponse B — La CEEAC.} Le COPAX est l'organe de paix et de sécurité de la Communauté économique des États de l'Afrique centrale (CEEAC), créé par le protocole de 1999 (pacte d'assistance mutuelle). Il s'appuie sur la FOMAC (Force multinationale de l'Afrique centrale). Piège : le confondre avec la CEMAC (intégration économique et monétaire) ou avec le CPS de l'Union africaine.
    \item \textbf{Réponse C — Accord de partenariat UE-ACP (Samoa).} Signé en novembre 2023 à Apia (Samoa), il succède à l'accord de Cotonou (2000-2020, prolongé). La convention de Lomé (1975) est l'ancien cadre.
    \item \textbf{Réponse B — Son histoire coloniale partagée entre la France et le Royaume-Uni.} Après 1916, le Cameroun est partagé en deux mandats puis deux tutelles : française (Cameroun oriental) et britannique (Cameroun septentrional et méridional). Cette double filiation explique le bilinguisme officiel et l'appartenance aux deux organisations (Francophonie dès 1960, Commonwealth en 1995).
    \item \textbf{Réponse C — Caractère concessionnel et objectif de développement.} L'APD est définie par l'OCDE/CAD : dons ou prêts à conditions très favorables (taux faible, longue maturité, différé d'amortissement) accordés par des organismes publics avec pour objectif principal le développement économique et le bien-être. Elle n'est ni un don systématique (B est trop restrictif), ni un investissement privé (D).
\end{enumerate}
\end{corrige}

\begin{corrige}[Exercice 2 — Vrai ou Faux justifié : Réalités de la coopération]
\begin{enumerate}
    \item \textbf{FAUX.} Le Cameroun fournit des contingents militaires et policiers aux OMP de l'ONU : il est notamment l'un des principaux contributeurs de troupes à la MINUSCA (République centrafricaine) avec un bataillon d'infanterie, et il a participé à la MINUSMA (Mali) et à d'autres missions (observateurs militaires, unités de police constituées FPU). Cette participation relève du dossier « La participation du Cameroun dans les missions des Nations unies ».
    \item \textbf{VRAI.} La Chine est devenue le premier partenaire commercial du Cameroun et le premier bailleur bilatéral en volume d'engagements récents, principalement via les prêts d'Exim Bank of China pour les infrastructures (Kribi, Memve'ele, autoroutes, stades). Elle détient une part majoritaire (environ 60 à 70 \%) de la dette extérieure bilatérale camerounaise.
    \item \textbf{FAUX.} Les programmes du FMI (FEC — Facilité élargie de crédit, comme celle de 2017-2020 puis 2021-2024) sont assortis de \cle{conditionnalités} macroéconomiques : réduction du déficit budgétaire, maîtrise de l'endettement, réformes structurelles (gouvernance, transparence des marchés publics, subventions). Les décaissements sont par tranches, conditionnés à des revues périodiques.
    \item \textbf{FAUX.} Le COPAX a une mission beaucoup plus large : prévention, gestion et règlement des conflits en Afrique centrale, maintien et consolidation de la paix, diplomatie préventive, lutte contre la criminalité transfrontalière et le terrorisme. La gestion des catastrophes n'est qu'une composante marginale (mécanisme d'assistance humanitaire).
    \item \textbf{VRAI.} La coopération triangulaire associe un pays émergent (ex : Brésil, expertise agricole tropicale), un pays en développement bénéficiaire (ex : le Cameroun ou un pays tiers africain) et un bailleur traditionnel ou multilatéral qui cofinance (ex : Japon, UE). Elle valorise le transfert de savoir-faire Sud-Sud adapté aux contextes tropicaux.
\end{enumerate}
\end{corrige}

\begin{corrige}[Exercice 3 — Définitions et mise en relation : Concepts clés]
\textbf{Définitions :}
\begin{enumerate}
    \item \cle{Conditionnalité} : ensemble d'exigences (réformes budgétaires, structurelles, de gouvernance) qu'un bailleur (FMI, Banque mondiale, UE) impose à un État bénéficiaire pour décaisser son aide ou ses prêts. Elle vise à garantir l'efficacité de l'aide et la soutenabilité macroéconomique, mais est critiquée comme une entorse à la souveraineté.
    \item \cle{Diplomatie économique} : utilisation des instruments diplomatiques (ambassades, accords, sommets) au service des intérêts économiques du pays : promotion des exportations, attraction des investissements directs étrangers, recherche de financements et de débouchés.
    \item \cle{Sécurité collective} : principe selon lequel la paix et la sécurité de chaque État membre concernent l'ensemble de la communauté (ONU, chapitre VII ; UA, article 4h — droit d'intervention). Une agression ou une menace contre un membre appelle une réponse collective, y compris coercitive.
    \item \cle{Partenariat gagnant-gagnant} : rhétorique de la coopération sino-africaine (FOCAC) affirmant que les échanges profitent aux deux parties sans ingérence politique : infrastructures et financements pour l'Afrique, matières premières et débouchés pour la Chine. À analyser de manière critique (asymétrie des gains, endettement).
    \item \cle{Intégration sous-régionale} : processus de rapprochement économique, monétaire et politique entre États voisins (CEMAC : monnaie commune FCFA, BEAC ; CEEAC : libre circulation, paix et sécurité) visant à créer un marché commun et à mutualiser les politiques publiques.
\end{enumerate}
\textbf{Mise en relation :}
\begin{center}
\begin{tabularx}{\linewidth}{|l|X|}\hline
\rowcolor{popblueL}
\textbf{Concept} & \textbf{Institution / Partenaire principal} \\ \hline
COPAX & CEEAC (Communauté économique des États de l'Afrique centrale) \\ \hline
FEC (Facilité élargie de crédit) & FMI (Fonds monétaire international) \\ \hline
OMP (Casques bleus) & ONU (Conseil de sécurité / Département des opérations de paix) \\ \hline
Forum sur la coopération sino-africaine (FOCAC) & Chine (et États africains membres) \\ \hline
Accord de Cotonou / Samoa & Union européenne et pays ACP \\ \hline
\end{tabularx}
\end{center}
\end{corrige}

\begin{corrige}[Exercice 4 — Repères chiffrés et chronologiques]
\begin{center}
\begin{tabularx}{\linewidth}{|l|X|}\hline
\rowcolor{popblueL}
\textbf{Indicateur / Événement} & \textbf{Valeur / Date / Sigle} \\ \hline
Année d'adhésion du Cameroun à l'ONU & 1960 (20 septembre 1960) \\ \hline
Année d'adhésion à l'UA (ex-OUA) & 1963 (membre fondateur de l'OUA, Addis-Abeba) \\ \hline
Année d'adhésion au Commonwealth & 1995 \\ \hline
Part approximative de la dette extérieure camerounaise détenue par la Chine (environ) & Environ 60 à 70 \% de la dette bilatérale (environ 25 à 30 \% de la dette extérieure totale) \\ \hline
Nombre de piliers de la stratégie FEC 2021-2024 (FMI) & 3 piliers \\ \hline
Nom de la force multinationale mixte (Lac Tchad) & FMM (Force multinationale mixte) — MNJTF en anglais \\ \hline
Sigle de l'instrument financier de l'UE pour l'APD (2021-2027) & NDICI — Global Europe \\ \hline
\end{tabularx}
\end{center}
\textbf{Points de vigilance :} ne pas confondre adhésion à l'OUA (1963) et création de l'UA (2002) ; le Commonwealth date de 1995, pas de 1960 ; la FMM associe Cameroun, Nigeria, Niger, Tchad et Bénin contre Boko Haram.
\end{corrige}

\begin{corrige}[Exercice 5 — Analyse de document : Le discours de la Tribune onusienne]
\begin{enumerate}
    \item \textbf{Les trois registres :}
    \begin{itemize}
        \item \emph{Registre institutionnel} : attachement au multilatéralisme et plaidoyer pour la réforme du Conseil de sécurité.
        \item \emph{Registre sécuritaire} : lutte contre le terrorisme (bassin du Lac Tchad) et gestion de la crise anglophone (Nord-Ouest/Sud-Ouest).
        \item \emph{Registre économique et normatif} : appel à un « partenariat rénové » fondé sur l'équité (financement du développement, respect de la souveraineté).
    \end{itemize}
    \item \textbf{La réforme du Conseil de sécurité :} l'Afrique, qui compte 54 États membres de l'ONU et fournit l'essentiel des dossiers traités par le Conseil, n'y dispose d'aucun siège permanent avec droit de veto. Le \cle{consensus d'Ezulwini} (position commune africaine, 2005) revendique au moins deux sièges permanents (avec droit de veto) et cinq sièges non permanents pour l'Afrique. L'enjeu : corriger une injustice historique de représentation et rendre la gouvernance mondiale plus légitime et plus efficace.
    \item \textbf{Lien sécurité intérieure / diplomatie multilatérale :} le « lourd tribut » (pertes humaines, déplacés, coûts budgétaires) montre que les crises internes ont une dimension transnationale (Boko Haram vient du bassin du Lac Tchad) : elles justifient la mobilisation de cadres multilatéraux (FMM, appui de l'ONU, financements internationaux) et donnent au Cameroun une légitimité pour réclamer solidarité internationale et ressources. La sécurité intérieure devient ainsi un argument diplomatique.
    \item \textbf{« Partenariat rénové fondé sur le respect de la souveraineté » :} vis-à-vis des partenaires traditionnels (UE, FMI), c'est une demande d'allègement des conditionnalités perçues comme intrusives ; vis-à-vis des nouveaux partenaires (Chine), c'est la valorisation du principe de non-ingérence, tout en rappelant que le Cameroun entend définir lui-même ses priorités (Agenda 2063, SND30) et ne pas être l'otage des rivalités entre puissances.
\end{enumerate}
\end{corrige}

\begin{corrige}[Exercice 6 — Étude de cas : La coopération Cameroun - Chine]
\begin{enumerate}
    \item \textbf{Trois avantages majeurs :}
    \begin{itemize}
        \item \emph{Rapidité d'exécution et coûts compétitifs} : les entreprises chinoises réalisent les grands ouvrages (Kribi, Memve'ele) dans des délais souvent plus courts et à moindre coût que les partenaires traditionnels.
        \item \emph{Absence de conditionnalité politique} : la Chine applique le principe de non-ingérence (pas d'exigences sur les droits humains ou la gouvernance), ce qui préserve la marge de manœuvre de l'État.
        \item \emph{Accès à des financements et à une technologie} que les bailleurs classiques refusaient ou retardaient : infrastructures structurantes pour l'émergence (énergie, ports, routes).
    \end{itemize}
    \item \textbf{Deux risques majeurs :}
    \begin{itemize}
        \item \emph{Soutenabilité de la dette} : l'envolée de l'encours envers la Chine (environ 4 milliards USD) pèse sur le service de la dette et réduit l'espace budgétaire ; le FMI classe le Cameroun à « risque élevé de surendettement ».
        \item \emph{Opacité des clauses et faible contenu local} : contrats confidentiels, clauses de confidentialité, main-d'œuvre largement importée, faibles transferts de compétences, parfois impacts environnementaux et sociaux insuffisamment traités.
    \end{itemize}
    \item \textbf{Pertinence du « piège de la dette » :} le concept doit être nuancé. Une grande partie des prêts chinois au Cameroun est \emph{concessionnelle} (taux faibles, longues maturités), ce qui atténue le risque ; il n'y a pas eu de saisie d'actif stratégique comparable au port de Hambantota (Sri Lanka). Cependant, la part croissante de la \emph{dette commerciale} et le regroupement des échéances créent une vulnérabilité réelle : le risque est moins un « piège » délibéré qu'une dépendance financière limitant la souveraineté budgétaire. La vigilance s'impose donc sans céder à la caricature.
    \item \textbf{Deux mesures de gouvernance :}
    \begin{itemize}
        \item Instaurer des \emph{clauses de contenu local} (quotas de main-d'œuvre camerounaise, sous-traitance aux PME locales, transfert de technologie et formation) dans tout nouveau contrat.
        \item Publier les contrats et soumettre les grands projets à un \emph{audit indépendant} (Cour des comptes, CAA) et à une évaluation d'impact ex ante ; négocier si nécessaire des rééchelonnements (comme ceux obtenus en 2019 et dans le cadre de l'ISSD du G20).
    \end{itemize}
\end{enumerate}
\end{corrige}

\begin{corrige}[Exercice 7 — Schéma synthèse : Architecture institutionnelle de la coopération]
\textbf{Organigramme attendu (trois niveaux) :}
\begin{center}
\begin{tabularx}{\linewidth}{|l|X|X|}\hline
\rowcolor{popblueL}
\textbf{Niveau} & \textbf{Acteurs / Cadres (sigles)} & \textbf{Objectif stratégique principal} \\ \hline
Universel (multilatéral global) & ONU (AG, CS, ECOSOC, OMS, UNESCO, PNUD, OMP) ; FMI ; BM ; OMC & Paix et sécurité collective ; financement et stabilité macroéconomique ; légitimité internationale \\ \hline
Régional et sous-régional (Afrique) & UA (CPS) ; CEEAC (COPAX, FOMAC) ; CEMAC (BEAC, COBAC) ; CLC / FMM (Lac Tchad) & Intégration économique et monétaire ; sécurité collective sous-régionale ; libre circulation \\ \hline
Bilatéral et interrégional & Historiques : France, Royaume-Uni ; UE ; Émergents : Chine (FOCAC), Japon (TICAD), Brésil, Turquie, Inde, Corée ; Blocs : OIF, Commonwealth, OCI & Diversification des financements ; investissements et infrastructures ; désendettement et rayonnement diplomatique \\ \hline
\end{tabularx}
\end{center}
\textbf{Points de vigilance :} placer le COPAX sous la CEEAC (et non la CEMAC) ; la BEAC et la COBAC relèvent de la CEMAC ; la FMM est un cadre ad hoc du bassin du Lac Tchad adossé à l'UA ; le TICAD est le cadre de coopération avec le Japon, le FOCAC celui avec la Chine. Chaque niveau doit être associé à un objectif distinct (sécurité, intégration, diversification).
\end{corrige}

\begin{corrige}[Exercice 8 — Argumentation dirigée : Diversification vs Dépendance]
\textbf{Proposition de paragraphe argumenté (corrigé-type) :}

La diversification des partenariats internationaux constitue d'abord, pour le Cameroun, une stratégie d'élargissement de ses marges de manœuvre. En s'ouvrant à la Chine, au Brésil, à la Turquie, à l'Inde ou à la Russie, Yaoundé dispose d'alternatives de financement (prêts concessionnels d'Exim Bank, coopérations Sud-Sud) qui lui évitent de dépendre du seul arbitrage des bailleurs traditionnels et renforcent son pouvoir de négociation : la concurrence entre partenaires permet d'obtenir de meilleures conditions et de financer des infrastructures (Kribi, Memve'ele) que l'aide classique ne finançait plus. Cette politique s'inscrit dans une tradition de non-alignement et de diplomatie équilibrée héritée de l'indépendance. Toutefois, cette diversification engendre de nouvelles vulnérabilités : la multiplication des créanciers fragmente l'aide et complique sa coordination ; une partie de la dette contractée auprès des nouveaux partenaires est opaque et moins concessionnelle, ce qui dégrade la soutenabilité budgétaire (risque élevé de surendettement selon le FMI) ; les clauses de confidentialité, la main-d'œuvre importée et l'affaiblissement des normes sociales et environnementales réduisent les retombées locales ; enfin, dans un contexte de rivalité entre grandes puissances, le Cameroun s'expose à des pressions d'alignement géopolitique contraires à sa souveraineté. En synthèse, la véritable souveraineté ne réside ni dans la dépendance aux anciens partenaires ni dans la simple substitution de nouveaux créanciers, mais dans la capacité de l'État à \emph{négocier}, \emph{coordonner} et \emph{hiérarchiser} l'ensemble des flux de coopération : cela suppose un cadre national de partenariat transparent, la publication des contrats, des audits indépendants et l'ancrage de tous les accords aux priorités de la SND30 et de l'Agenda 2063. La diversification n'est donc vertueuse que si elle est gouvernée.

\textbf{Points de vigilance :} respecter la structure thèse / antithèse / synthèse ; mobiliser au moins deux exemples précis (Kribi, FMI, FOCAC) ; la synthèse ne doit pas être un simple compromis mais un dépassement (la gouvernance comme clé).
\end{corrige}

\begin{corrige}[Exercice 9 — Sujet type Bac : Analyse de documents — Diversification des partenariats et nouveaux risques]
\textbf{Barème indicatif (20 pts) :} analyse des documents 6 pts (2 pts par question) ; mise en relation et synthèse 8 pts ; note de synthèse 6 pts.

\textbf{1. Analyse des documents (6 pts)}

a) \emph{Part combinée des partenaires traditionnels :}
\[ \text{UE} + \text{France} + \text{FMI} + \text{BM} = 35 + 18 + 5 + 20 = 78\,\% \]
Contre $12\,\%$ pour la Chine. En montants : $350 + 180 + 50 + 200 = 780$ millions USD contre 120 millions USD. \textbf{Interprétation :} malgré le discours sur la diversification, l'APD reste massivement dominée par les partenaires traditionnels (près des 4/5 des flux) : la diversification est réelle mais encore marginale dans l'aide publique ; elle est en revanche plus forte dans les prêts et investissements (Document 2), ce qui distingue APD et endettement.

b) \emph{TCAM de la dette envers la Chine, 2010-2023 :} valeur initiale $V_0 = 0{,}3$ milliard USD (2010), valeur finale $V_1 = 4{,}0$ milliards USD (2023), sur $n = 13$ ans :
\[ \text{TCAM} = \left(\frac{4{,}0}{0{,}3}\right)^{1/13} - 1 = (13{,}33)^{1/13} - 1 \approx 1{,}221 - 1 \approx 22\,\% \text{ par an} \]
(Vérification : $0{,}3 \times 1{,}22^{13} \approx 0{,}3 \times 13{,}3 \approx 4{,}0$.) La dette envers la Chine a donc été multipliée par environ 13 en 13 ans. \emph{Infléchissement 2023} (4,2 $\rightarrow$ 4,0) : ralentissement des nouveaux emprunts, début de remboursement du principal, effet des rééchelonnements (ISSD du G20) et de la discipline imposée par le programme FMI ; le Cameroun atteint une phase de consolidation de son endettement.

c) \emph{Trois principes directeurs (Document 3) :} (1) le refus que l'Afrique soit l'enjeu des rivalités de puissances (non-alignement) ; (2) des partenariats « gagnant-gagnant » respectueux de la souveraineté et des priorités africaines (Agenda 2063) ; (3) l'attachement au multilatéralisme efficace comme garant de la sécurité collective.

\textbf{2. Mise en relation et synthèse (8 pts) — éléments attendus :}
\begin{itemize}
    \item La diversification réduit la dépendance : alternatives de financement (Chine : 12 \% de l'APD mais surtout prêts d'infrastructures hors APD), pouvoir de négociation accru, non-alignement revendiqué (Doc. 3).
    \item Mais les documents montrent les limites : les partenaires traditionnels pèsent encore 78 \% de l'APD (Doc. 1) ; la dette envers la Chine a explosé (TCAM $\approx 22\,\%$, Doc. 2), créant une nouvelle dépendance financière.
    \item Nouveaux risques à citer : soutenabilité de la dette (risque élevé de surendettement), fragmentation de l'aide et coûts de coordination, opacité des contrats, conditionnalités différenciées (conditionnalités économiques du FMI vs non-ingérence chinoise), pressions d'alignement géopolitique.
    \item Conclusion : la diversification n'émancipe que si elle est gouvernée (transparence, coordination, ancrage SND30/Agenda 2063).
\end{itemize}
\textbf{Points de vigilance :} croiser explicitement les trois documents ; distinguer APD (Doc. 1) et dette (Doc. 2) ; ne pas opposer mécaniquement « bons » et « mauvais » partenaires.

\textbf{3. Note de synthèse au Ministre (6 pts) — corrigé-type :}

\emph{À l'attention de M. le Ministre des Relations Extérieures — Objet : gouvernance de la coopération à l'horizon 2030.}

Notre coopération reste dominée par les partenaires traditionnels (78 \% de l'APD) tandis que la dette envers la Chine a crû d'environ 22 \% par an depuis 2010. Trois axes prioritaires sont proposés :
\begin{enumerate}
    \item \textbf{Transparence et soutenabilité de la dette} : publication systématique des contrats de prêts, audit indépendant ex ante des grands projets, plafonnement des emprunts non concessionnels et recours prioritaire aux dons et prêts concessionnels.
    \item \textbf{Coordination nationale de l'aide} : création d'un cadre national de partenariat unique (base de données des flux, revue annuelle avec les bailleurs) aligné sur la SND30 et l'Agenda 2063, afin de réduire la fragmentation et les doublons.
    \item \textbf{Maximisation des retombées locales} : clauses de contenu local (emploi, sous-traitance, transfert de technologie) dans tout accord bilatéral, et diplomatie économique active pour transformer les financements en industrialisation et en exportations.
\end{enumerate}
Ces mesures concilient diversification des partenaires et préservation de la souveraineté budgétaire et décisionnelle du Cameroun.

\textbf{Points de vigilance :} forme de la note (destinataire, objet, ton administratif) ; exactement trois axes, chacun assorti d'instruments concrets ; ancrage dans les documents.
\end{corrige}

\begin{corrige}[Exercice 10 — Sujet type Bac : Dissertation courte — OMP et Intérêt national]
\textbf{Barème indicatif (20 pts) :} introduction 4 pts ; développement (plan et argumentation) 10 pts ; conclusion 3 pts ; qualité rédactionnelle et exemples 3 pts.

\textbf{1. Analyse du sujet :} \cle{OMP} : opérations déployées par l'ONU (mandat du Conseil de sécurité) pour interposer des forces, protéger les civils et consolider la paix. \emph{Devoir moral} : solidarité internationale, responsabilité de protéger (R2P), vocation pacifiste du Cameroun. \emph{Intérêt national} : formation et équipement des troupes, revenus (remboursements onusiens), rayonnement diplomatique, sécurisation du voisinage (la paix en RCA protège l'Est camerounais). \emph{Contrainte capacitaire} : limites logistiques et financières, pertes humaines, fatigue opérationnelle d'une armée déjà engagée sur trois fronts internes (Boko Haram, crise anglophone, Est).

\textbf{2. Plan détaillé :}
\begin{itemize}
    \item \textbf{I. Un devoir moral et juridique assumé} : attachement au multilatéralisme ; solidarité africaine ; contributions concrètes (MINUSCA, FPU, observateurs).
    \item \textbf{II. Un intérêt national bien compris} : bénéfices militaires (entraînement, équipement, remboursements ONU), diplomatiques (crédibilité, candidatures, réforme du CS) et sécuritaires (stabiliser le voisinage, c'est sécuriser le Cameroun).
    \item \textbf{III. Des contraintes capacitaires qui imposent des choix} : coûts humains et financiers, tension avec les besoins internes, dépendance logistique ; nécessité d'une contribution ciblée et soutenable.
\end{itemize}

\textbf{3. Introduction (corrigé-type) :} Le 20 septembre 1960, le Cameroun nouvellement indépendant adhère à l'ONU et s'engage dans la sécurité collective. Devenu l'un des principaux contributeurs africains de troupes, notamment à la MINUSCA, il déploie ses soldats et policiers sous le casque bleu. Cette participation soulève une question : relève-t-elle d'un devoir moral de solidarité, d'un calcul d'intérêt national, ou se heurte-t-elle aux limites capacitaires d'un État lui-même confronté à des crises sécuritaires ? Nous montrerons que l'engagement camerounais dans les OMP conjugue devoir moral et intérêt national, avant d'en examiner les contraintes capacitaires.

\textbf{4. Développement — éléments attendus :}
\begin{itemize}
    \item \emph{Devoir moral} : Charte de l'ONU, R2P, solidarité panafricaine ; hôpital de campagne camerounais, unités de police constituées (FPU) protégeant les civils en RCA.
    \item \emph{Intérêt national} : la MINUSCA sécurise la frontière Est et le corridor Douala-Bangui (voie d'approvisionnement) ; remboursements onusiens par soldat déployé ; formation et expérience acquises ; prestige diplomatique soutenant les revendications camerounaises (réforme du Conseil de sécurité). Exemples : MINUSCA (bataillon), contribution passée à la MINUSMA (Mali, retrait en 2023), observateurs à la MONUSCO (RDC) et à la MINURSO (Sahara occidental) ; engagement dans la FMM du bassin du Lac Tchad (cadre ad hoc adossé à l'UA).
    \item \emph{Contraintes} : pertes humaines, coûts logistiques avancés avant remboursement, armée sollicitée sur les fronts internes (Extrême-Nord, NOSO), limites d'équipement ; d'où une contribution sélective (RCA prioritaire car voisine).
\end{itemize}

\textbf{5. Conclusion (corrigé-type) :} La participation du Cameroun aux OMP n'oppose pas devoir moral et intérêt national : elle les articule, la solidarité internationale servant la sécurité et le rayonnement du pays. Mais les contraintes capacitaires imposent une gestion rigoureuse des engagements. L'ouverture : la réforme en cours des opérations de paix (partenariats ONU-UA, financement des opérations africaines) offre au Cameroun l'occasion de plaider pour des missions mieux dotées, où les contributeurs de troupes auraient voix au chapitre.

\textbf{Points de vigilance :} définir les termes dès l'introduction ; équilibrer les parties ; citer des exemples précis (MINUSCA, FMM, FPU) ; ne pas confondre OMP de l'ONU et opérations sous mandat de l'UA ; la conclusion doit répondre à la problématique et ouvrir sans introduire d'argument nouveau.
\end{corrige}

\begin{corrige}[Exercice 11 — Sujet type Bac : Croquis commenté — Le Cameroun, carrefour de coopérations]
\textbf{Barème indicatif (20 pts) :} fond de carte exact et orienté 4 pts ; item 1 (ancrages) 4 pts ; item 2 (flux et projets) 5 pts ; item 3 (enjeux sécuritaires) 4 pts ; légende hiérarchisée et titre 3 pts.

\textbf{1. Fond de carte attendu :} contour simplifié du Cameroun (forme en triangle allongé vers le nord), flèche du nord, échelle indicative ; Yaoundé (capitale politique, étoile) et Douala (capitale économique, point carré) ; villes : Garoua, Maroua, Ngaoundéré, Bertoua, Bafoussam, Bamenda, Buea, Ebolowa, Kribi ; pays limitrophes nommés : Nigeria (ouest), Tchad (nord-est), RCA (est), Congo, Gabon, Guinée équatoriale (sud) ; océan Atlantique.

\textbf{2. Item 1 — Ancrages institutionnels (figurés ponctuels) :} étoile rouge sur Yaoundé : sièges des représentations (délégation UE, agences ONU, ambassades), institutions sous-régionales ; symboles/sigles ONU, UA, CEEAC, CEMAC, OIF, Commonwealth, OCI placés en encadré-rappel reliés à Yaoundé ; flèches entrantes fines figurant les représentations diplomatiques.

\textbf{3. Item 2 — Flux financiers et projets structurants (figurés linéaires et ponctuels) :}
\begin{itemize}
    \item Flèches d'épaisseur proportionnelle convergeant vers Douala et Kribi : bleu épais (UE + France, 53 \% de l'APD), vert (Banque mondiale/BAfD, 20 \%), rouge (Chine, 12 \% de l'APD mais premier créancier bilatéral, environ 4 milliards USD d'encours).
    \item Pictogrammes localisés : port en eau profonde de Kribi (ancre) ; pipeline Tchad-Cameroun (trait gras Tchad $\rightarrow$ Kribi) ; barrage de Memve'ele et centrale de Nachtigal (éclair, Sud) ; autoroute Yaoundé-Nsimalen (double trait) ; stade d'Olembé et hôpital de référence (Yaoundé).
    \item Bulles chiffrées : « APD totale $\approx$ 1 milliard USD/an », « Dette envers la Chine $\approx$ 4 Mds USD ».
\end{itemize}

\textbf{4. Item 3 — Enjeux sécuritaires et diplomatiques (figurés de surface et flèches) :}
\begin{itemize}
    \item Surfaces hachurées : Extrême-Nord (insécurité Boko Haram, FMM), Nord-Ouest/Sud-Ouest (crise anglophone), Est (réfugiés centrafricains).
    \item Flèches sortantes vers la RCA (MINUSCA), le Mali (MINUSMA, contribution passée), la RDC (MONUSCO) avec icône casque bleu.
    \item Symbole « main serrée » sur Yaoundé : médiation et vocation de paix ; rappel du COPAX/FOMAC (Afrique centrale).
\end{itemize}

\textbf{5. Légende hiérarchisée (3 items numérotés) :} 1. Ancrages institutionnels (ponctuels : étoiles, sigles) ; 2. Flux et projets (linéaires proportionnels + pictogrammes + chiffres) ; 3. Enjeux sécuritaires (surfaces hachurées + flèches sortantes). Titre obligatoire en haut : « Le Cameroun, carrefour de coopérations multilatérales et bilatérales ».

\textbf{Points de vigilance :} tout figuré utilisé doit figurer dans la légende (et réciproquement) ; hiérarchiser les flèches selon les montants (proportionnalité) ; ne pas surcharger : un figuré = une information ; placer correctement les villes (Kribi au sud-ouest sur la côte, Maroua à l'Extrême-Nord) ; orienter la carte (nord) ; les chiffres cités doivent être cohérents avec le cours (parts d'APD, encours de la dette chinoise).

\textbf{Proposition de réalisation cartographique (TikZ) :}
\begin{popfigure}
\begin{tikzpicture}[scale=0.85, every node/.style={font=\footnotesize}]
    % Fond de carte - Contour simplifié du Cameroun
    \draw[popdark, thick, fill=popcream!30] (0,0) -- (0,8) -- (2,10) -- (5,10) -- (8,7) -- (8,0) -- cycle;
    % Mer
    \draw[popblue, thick, fill=popblueL!30] (0,0) -- (8,0) -- (8,-0.5) -- (0,-0.5) -- cycle;
    \node[popblue, rotate=90] at (-0.5, -0.25) {Océan Atlantique};
    % Flèche Nord
    \draw[->, thick, popdark] (7.5, 9.5) -- (7.5, 10.2) node[above] {N};
    % Villes principales
    \node[star, star points=5, star point ratio=2.25, fill=popgold, draw=popdark, inner sep=2pt] (Ya) at (3.5, 3.5) {};
    \node[above right] at (Ya) {Yaoundé};
    \node[rectangle, fill=poporange, draw=popdark, inner sep=2pt] (Do) at (1.2, 1.8) {};
    \node[below right] at (Do) {Douala};
    \foreach \name/\pos in {Garoua/{(4,8)}, Maroua/{(3,9)}, Ngaoundere/{(4.5,5.5)}, Bertoua/{(6,3.5)}, Bafoussam/{(2.5,4.5)}, Bamenda/{(1.5,5.5)}, Buea/{(0.8,2.5)}, Ebolowa/{(4,1.5)}, Kribi/{(0.5,0.3)}} {
        \node[circle, fill=popdark, inner sep=1.5pt] at \pos {};
        \node[above right, font=\tiny] at \pos {\name};
    }
    % Pays limitrophes
    \foreach \name/\pos/\anchor in {Nigeria/{(-0.5,4)}/east, Tchad/{(3.5,10.5)}/south, RCA/{(6.5,8.5)}/west, Congo/{(8.5,4)}/west, Gabon/{(8.5,1)}/west, Guinée Eq./{(8.5,-0.5)}/west} {
        \node[font=\small\bfseries, anchor=\anchor] at \pos {\name};
    }
    % Item 1: Ancrages institutionnels (Yaoundé)
    \node[draw=poppurple, rounded corners, fill=poppurpleL!30, inner sep=3pt, anchor=south west] at (3.7, 3.7) {ONU, UA, CEEAC, CEMAC, OIF, Commonwealth, OCI};
    % Item 2: Flux et projets
    % Flèches entrantes proportionnelles
    \draw[popblue, line width=3pt, ->] (-1, 2) -- (Do); % UE/France
    \draw[popblue, line width=3pt, ->] (-1, 5) -- (Ya);
    \node[popblue, font=\tiny, rotate=90] at (-1.2, 3.5) {UE/France (53\%)};
    \draw[popgreen, line width=1.5pt, ->] (4, -1) -- (Do); % BM/BAD
    \node[popgreen, font=\tiny] at (4, -1.3) {BM/BAD (20\%)};
    \draw[popred, line width=2.5pt, ->] (9, 3) -- (Ya); % Chine
    \draw[popred, line width=2.5pt, ->] (9, 1) -- (0.5, 0.3); % Chine -> Kribi
    \node[popred, font=\tiny, rotate=-90] at (9.2, 2) {Chine (12\% APD, ~4 Mds \$ dette)};
    % Projets
    \node[anchor=west, font=\tiny] at (0.5, 0.1) {⚓ Port Kribi};
    \draw[poporange, thick] (3.5, 9.5) -- (0.5, 0.3); % Pipeline
    \node[poporange, font=\tiny, rotate=-20] at (2, 5) {Pipeline Tchad-Kribi};
    \node[font=\tiny] at (4, 1.2) {⚡ Memve'ele};
    \node[font=\tiny] at (3.5, 4) {⚡ Nachtigal};
    \node[font=\tiny] at (3, 3.8) {⛶ Autoroute Ydé-Nsimalen};
    \node[font=\tiny] at (3.7, 3.3) {🏟️ Stade Olembé};
    % Bulles chiffrées
    \node[draw=popblue, rounded corners, fill=popblueL!30, font=\tiny, anchor=north east] at (7.8, 9.8) {APD $\approx$ 1 Md \$/an};
    \node[draw=popred, rounded corners, fill=popred!10, font=\tiny, anchor=north west] at (0.2, 9.8) {Dette Chine $\approx$ 4 Mds \$};
    % Item 3: Enjeux sécuritaires
    % Hachures (patterns)
    \pattern[pattern=north east lines, pattern color=popred!50] (2, 8.5) rectangle (4.5, 10); % Extrême-Nord
    \node[popred, font=\tiny, align=center] at (3.25, 9.2) {Extrême-Nord\\(Boko Haram, FMM)};
    \pattern[pattern=north west lines, pattern color=poporange!50] (0.5, 5) rectangle (2, 7); % NOSO
    \node[poporange, font=\tiny, align=center] at (1.25, 6) {NOSO\\(Crise anglophone)};
    \pattern[pattern=dots, pattern color=poppurple!50] (5.5, 2.5) rectangle (7.5, 5); % Est
    \node[poppurple, font=\tiny, align=center] at (6.5, 3.7) {Est\\(Réfugiés RCA)};
    % Flèches sortantes OMP
    \draw[->, popblue, thick] (Ya) -- (7, 5) node[midway, above, font=\tiny] {MINUSCA (RCA)};
    \draw[->, popblue, dashed, thick] (Ya) -- (-0.5, 9) node[midway, left, font=\tiny] {MINUSMA (Mali, passé)};
    \draw[->, popblue, thick] (Ya) -- (8.5, 4) node[midway, above, font=\tiny] {MONUSCO (RDC)};
    % Médiation
    \node[font=\Large] at (Ya) {🤝};
    \node[font=\tiny, align=center] at (3.5, 2.8) {COPAX / FOMAC};
\end{tikzpicture}
\legende{Croquis schématique : Le Cameroun, carrefour de coopérations multilatérales et bilatérales. 1. Ancrages institutionnels (Yaoundé). 2. Flux financiers proportionnels (Bleu: UE/France, Vert: BM/BAD, Rouge: Chine) et projets structurants. 3. Zones d'insécurité (hachures) et engagements OMP sortants (flèches bleues).}
\end{popfigure}
\end{corrige}

\newpage

\coursec{Fiche bilan}

\begin{popfigure}
\begin{tikzpicture}[
  mindmap,
  concept color=popblue,
  level 1 concept/.append style={level distance=4.5cm,sibling angle=51.4},
  level 2 concept/.append style={level distance=3cm,sibling angle=45},
  every node/.style={font=\small, text=popdark},
  concept/.append style={minimum size=1.8cm, inner sep=4pt},
  grow cyclic
]
\node[concept, concept color=popblue, text=white, minimum size=2.2cm] {Cameroun \& Coopération Multilatérale}
  child[concept color=popgreen] { node[concept] {ONU \& Maintien Paix}
    child { node[concept, scale=0.7] {Casques bleus} }
    child { node[concept, scale=0.7] {Résolutions CS} }
    child { node[concept, scale=0.7] {Dossier 2} }
  }
  child[concept color=poporange] { node[concept] {Coopération Régionale}
    child { node[concept, scale=0.7] {UA \& Agenda 2063} }
    child { node[concept, scale=0.7] {CEMAC/CEEAC} }
    child { node[concept, scale=0.7] {COPAX (D4)} }
  }
  child[concept color=poppurple] { node[concept] {Partenariats Institutionnels}
    child { node[concept, scale=0.7] {UE (Accords Cotonou/Samoa)} }
    child { node[concept, scale=0.7] {Francophonie} }
    child { node[concept, scale=0.7] {Commonwealth} }
    child { node[concept, scale=0.7] {OCI} }
  }
  child[concept color=poppink] { node[concept] {Grands Partenaires Bilatéraux}
    child { node[concept, scale=0.7] {France \& R-U (Historiques)} }
    child { node[concept, scale=0.7] {Chine \& Japon (Asie)} }
    child { node[concept, scale=0.7] {USA \& Brésil (Amériques)} }
  }
  child[concept color=popgold] { node[concept] {Aide au Développement (D3)}
    child { node[concept, scale=0.7] {APD : Dons/Prêts} }
    child { node[concept, scale=0.7] {Conditionnalités} }
    child { node[concept, scale=0.7] {Efficacité (PDSC)} }
  }
  child[concept color=popred] { node[concept] {Gestion Conflits (TD2)}
    child { node[concept, scale=0.7] {Médiation} }
    child { node[concept, scale=0.7] {Force multinationale (MNJTF)} }
    child { node[concept, scale=0.7] {Refugiés/IDPs} }
  }
  child[concept color=popmuted] { node[concept] {Problèmes \& Défis (D5)}
    child { node[concept, scale=0.7] {Endettement} }
    child { node[concept, scale=0.7] {Crises sécuritaires (NOSO, Extrême-Nord)} }
    child { node[concept, scale=0.7] {Diplomatie économique} }
  };
\end{tikzpicture}
\legende{Carte mentale de synthèse : Le Cameroun dans la coopération multilatérale (7 axes majeurs).}
\end{popfigure}

\begin{bilan}
\courssub{Le Cameroun à l'ONU : Engagement universel}
\begin{itemize}
  \item \textbf{Adhésion :} 20 septembre 1960. Respect de la \cle{Charte des Nations Unies} (souveraineté, non-ingérence, règlement pacifique).
  \item \textbf{Participation aux missions (Dossier 2) :} Fournisseur majeur de \cle{casques bleus} (police, militaires, observateurs) : MINUSCA (RCA), MINUSMA (Mali), MONUSCO (RDC), UNMISS (Soudan du Sud).
  \item \textbf{Positions diplomatiques :} Candidatures aux organes (CSNU, ECOSOC, CDH), défense des pays en développement (G77+Chine), réforme du Conseil de Sécurité (représentativité africaine).
\end{itemize}

\courssub{Coopération Régionale : UA, CEMAC/CEEAC, COPAX}
\begin{itemize}
  \item \textbf{Union Africaine (UA) :} Siège de la \cle{Commission de l'UA} (Addis-Abeba) -- Cameroun membre actif. Ratification \cle{Agenda 2063}, ZLECAf (Zone de Libre-Échange Continentale Africaine), Protocole sur la libre circulation.
  \item \textbf{CEMAC/CEEAC :} Siège des institutions à \cle{Yaoundé} (Secrétariat CEMAC, Parlement, Cour de Justice). Objectifs : Union douanière, marché commun, monnaie unique (FCFA $\rightarrow$ Eco ?).
  \item \textbf{COPAX (Dossier 4) :} Conseil de Paix et Sécurité de l'Afrique Centrale. Mécanisme : \cle{MICOPAX} (Mission de consolidation de la paix), \cle{FOMAC} (Force multinationale). Rôle clé en RCA (2013, 2020).
\end{itemize}

\courssub{Partenariats Historiques \& Institutionnels}
\begin{center}
\begin{tabularx}{\linewidth}{|l|X|X|}\hline
\rowcolor{popblueL} \textbf{Cadre} & \textbf{Principaux Instruments} & \textbf{Enjeux Cameroun} \\ \hline
\textbf{Union Européenne} & Accords de Cotonou (2000) $\rightarrow$ Accord de Samoa (2023) & APD, Accords de Partenariat Économique (APE), Appui gouvernance, Migration. \\ \hline
\textbf{Francophonie (OIF)} & Sommets, Conférences ministérielles & Promotion français, diversité culturelle, appui éducation/formation (IFEF). \\ \hline
\textbf{Commonwealth} & Adhésion 1995 (Harare) & Anglophonie, échanges universitaires (ACU), commerce intra-Commonwealth. \\ \hline
\textbf{OCI} & Adhésion 1974 & Solidarité islamique, appui causes nationales, développement banque (BID). \\ \hline
\end{tabularx}
\end{center}

\courssub{Grands Partenaires Bilatéraux}
\begin{itemize}
  \item \textbf{France :} 1er partenaire historique (dette, APD, défense, culture). Accords de défense (1960, rénovés). Investissements (Bolloré, Total, Orange).
  \item \textbf{Royaume-Uni :} Partenaire anglophone historique (CDC, Pamol). Post-Brexit : Accord de partenariat commercial (2021). Commonwealth.
  \item \textbf{Chine :} 1er partenaire commercial \og tout flux\fg{} (import/export). FOCAC. Infrastructures (barrages, routes, stades, Assemblée nationale). Prêts concessionnels.
  \item \textbf{Japon :} TICAD. APD (dons, expertise technique), santé, éducation, riziculture (SEMRY).
  \item \textbf{USA :} AGOA (accès marché), appui sécurité (BIR, formation), santé (PEPFAR, PMI), gouvernance (USAID).
  \item \textbf{Brésil :} Coopération Sud-Sud. Agriculture (Embrapa), bioéthanol, défense, coton.
\end{itemize}

\courssub{L'Aide au Développement (Dossier 3)}
\begin{itemize}
  \item \textbf{Acteurs :} Bilatéraux (France, USA, Chine, Japon, Allemagne...), Multilatéraux (BM, FMI, BAD, UE, ONU), ONG, Fondations (Gates, Clinton).
  \item \textbf{Instruments :} \cle{APD} (Dons, Prêts concessionnels), Assistance technique, Allègement dette (PPTE), Appui budgétaire sectoriel.
  \item \textbf{Conditionnalités :} \cle{Gouvernance} (transparence, lutte corruption), \cle{Droits humains}, \cle{Réformes structurelles} (FMI: facilité élargie de crédit - FEC), \cle{Alignement} sur la SND30 (Stratégie Nationale de Développement).
  \item \textbf{Efficacité :} Principes de \cle{Paris (2005)}, \cle{Accra (2008)}, \cle{Busan (2011)} : Appropriation, Alignement, Harmonisation, Gestion axée résultats, Responsabilité mutuelle.
\end{itemize}

\courssub{Gestion des Conflits \& Sécurité Collective (TD2)}
\begin{itemize}
  \item \textbf{Crises internes :} \cle{Crise anglophone (NOSO)} depuis 2016 (revendications, sécessionnisme, DDR -- Désarmement Démobilisation Réintégration), \cle{Insurrection Boko Haram/ISWAP} (Extrême-Nord, Force Multinationale Mixte - FMM/MNJTF), \cle{Crise centrafricaine} (refugiés, insécurité frontalière Est).
  \item \textbf{Diplomatie préventive :} Médiation dans la sous-région (Tchad, Guinée Équatoriale, RCA).
  \item \textbf{Contribution forces :} Armée nationale, Police, Gendarmerie dans missions ONU/AU/CEEAC. \cle{Rôle de puissance d'équilibre} régional.
\end{itemize}

\courssub{Problèmes \& Défis des Relations Internationales (Dossier 5)}
\begin{itemize}
  \item \textbf{Endettement extérieur :} Risque de surendettement (ratio dette/PIB), négociations Club de Paris / Chine / Eurobonds. Gestion de la dette vs investissements.
  \item \textbf{Image internationale :} Crises sécuritaires (NOSO, Extrême-Nord) $\rightarrow$ impact sur notation financière, investissements, tourisme, droits humains (rapports ONG/ONU).
  \item \textbf{Diplomatie économique :} Diversification partenaires (Asie, Golfe, Turquie, Israël), attraction IDE (Investissements Directs Étrangers), promotion \cle{Made in Cameroon}, ZLECAf.
  \item \textbf{Intégration sous-régionale inachevée :} Barrières non tarifaires, libre circulation personnes/biens effective, monnaie unique CEMAC retardée.
  \item \textbf{Changement climatique :} Vulnérabilité (Sahel, côtière), financement adaptation (Fonds Vert), COP.
\end{itemize}
\end{bilan}

\begin{aretenir}[Checklist avant le Bac]
\begin{itemize}
  \item $\square$ Citer la date d'adhésion du Cameroun à l'ONU et nommer 3 missions de maintien de la paix actuelles (Dossier 2).
  \item $\square$ Expliquer le rôle du COPAX et citer ses deux bras opérationnels (MICOPAX, FOMAC) (Dossier 4).
  \item $\square$ Différencier les apports de la coopération avec l'UE (APE, Samoa) vs la Francophonie vs le Commonwealth vs l'OCI.
  \item $\square$ Classer les 6 grands partenaires bilatéraux (France, R-U, Chine, Japon, USA, Brésil) par domaine de prédilection (infra, commerce, sécurité, santé, éducation, agriculture).
  \item $\square$ Définir l'APD et distinguer don / prêt concessionnel / appui budgétaire (Dossier 3).
  \item $\square$ Citer 3 conditionnalités majeures de l'aide (Gouvernance, Droits humains, Réformes FMI).
  \item $\square$ Analyser la contribution du Cameroun à la sécurité régionale : FMM (Boko Haram), MICOPAX (RCA), Casques bleus ONU (TD2).
  \item $\square$ Identifier 3 défis majeurs des RI du Cameroun : Endettement, Crises sécuritaires (image), Intégration CEMAC (Dossier 5).
  \item $\square$ Relier la SND30 à l'alignement de l'aide extérieure (appropriation nationale).
  \item $\square$ Maîtriser le vocabulaire : \cle{Souveraineté}, \cle{Non-ingérence}, \cle{Multilatéralisme}, \cle{Bilatéralisme}, \cle{Conditionnalité}, \cle{PPTE}, \cle{ZLECAf}, \cle{DDR}.
  \item $\square$ Savoir rédiger une réponse structurée : Introduction (définitions/contexte), Développement (arguments illustrés par dossiers), Conclusion (bilan/perspectives).
\end{itemize}
\end{aretenir}

\findecours

\end{document}
